\documentclass[11pt,a4paper]{article}
\usepackage[T1]{fontenc}
\usepackage{lmodern,microtype}
\usepackage[a4paper,margin=27mm]{geometry}
\usepackage{amsmath,amssymb,amsthm,mathtools,booktabs,array}
\usepackage[dvipsnames]{xcolor}
\usepackage[colorlinks=true,linkcolor=MidnightBlue,citecolor=MidnightBlue,urlcolor=MidnightBlue]{hyperref}
\newtheorem{theorem}{Theorem}[section]
\newtheorem{proposition}[theorem]{Proposition}
\newtheorem{lemma}[theorem]{Lemma}
\newtheorem{corollary}[theorem]{Corollary}
\newtheorem{maintheorem}{Theorem}
\renewcommand{\themaintheorem}{A}
\newtheorem{classification}{Theorem}
\renewcommand{\theclassification}{B}
\newtheorem*{sharpbounds}{Sharp bound}
\numberwithin{equation}{section}
\setlength{\emergencystretch}{2em}
\setcounter{tocdepth}{1}
\hypersetup{pdftitle={Grothendieck rings of complex fusion categories: rank 5},pdfauthor={S\'ebastien Palcoux},pdfsubject={Spherical descent, formal codegrees and central twists}}
\title{Grothendieck rings of complex fusion categories:\\ rank 5}
\author{S\'ebastien Palcoux}\date{}
\begin{document}
\maketitle
\begin{abstract}
We classify the Grothendieck rings of fusion categories of rank five over
$\mathbb C$. There are exactly seventeen based rings, and each has a
unitary realization. Every such fusion category has a pseudo-unitary
Galois conjugate and admits a spherical structure. The sharp bounds on
the largest fusion multiplicity and the global Frobenius--Perron
dimension are four and $48+24\sqrt3$, respectively. The proof first
compares the ordinary and signed dimension characters through central
induction. Fourier estimates and integral traces establish Galois
pseudo-unitarity. We then bound the fusion rules by the action of the
invertible objects and by the degree of the dimension field. Exact
finite certificates complete the classification.
\end{abstract}
\begin{center}\begin{minipage}{0.88\linewidth}\small\itshape
GPT-6.2 Astra assisted with the research. Complete human verification is pending, and
the results should be treated with caution.
\end{minipage}\end{center}

{\small\tableofcontents}
\clearpage

\section{Introduction}
A fusion category has rank five when it has five isomorphism classes of
simple objects. Its Grothendieck ring records their tensor products,
with the simple classes as its distinguished basis. Our classification
is up to isomorphism preserving this basis. We use the standard terminology
for fusion categories, module categories and centers in
\cite[Chapters~4, 7 and~8]{EGNO}.

A pivotal structure is a monoidal natural isomorphism
$\mathrm{Id}_{\mathcal C}\to(-)^{**}$; it is spherical when its left and
right traces agree \cite[Section~2]{BW}. Etingof, Nikshych and Ostrik conjectured the
existence of a pivotal structure for every complex fusion category
\cite[Conjecture~2.8]{ENO}. We prove a stronger conclusion in rank five.
\begin{maintheorem}\label{thm:r5-galpu}
Every complex fusion category of rank five has a pseudo-unitary Galois
conjugate. Consequently it admits a spherical structure, and in
particular a pivotal structure.
\end{maintheorem}
Galois conjugation preserves the fusion rules and their
Frobenius--Perron dimensions. A pseudo-unitary conjugate has its
canonical positive spherical structure \cite[Proposition~8.23]{ENO}, which transports back to a
spherical structure on the original category.

Write $G$ for the group of invertible simple objects and set
\[
 D=\operatorname{FPdim}(\mathcal C),\qquad
 \mu=\max_{i,j,k}N_{ij}^k.
\]
\begin{sharpbounds}
Every complex fusion category of rank five satisfies
$\mu\le4$ and $D\le48+24\sqrt3$. Both bounds are attained.
\end{sharpbounds}
The multiplicity bound is attained by the two near-group rings
$C_4+4$ and $C_2^2+4$. The dimension bound is attained by the
$C_2$-equivariantization of the Izumi $C_3+6$ category.
\begin{classification}\label{thm:main-classification}
The Grothendieck rings of complex fusion categories of rank five are
exactly the seventeen based rings in Table~\ref{tab:models-five}.
For $|G|=1,2,3,4,5$ their numbers are $3,7,1,5,1$, respectively.
Each ring has a unitary realization.
\end{classification}
Thus complex, pivotal, spherical, pseudo-unitary and unitary
categorifiability give the same list of based rings at rank five.

% ===== models_table.tex =====
\begin{table}[ht]\centering\small
\begin{tabular}{@{}clp{.64\linewidth}@{}}\toprule
$|G|$ & $\mu$ & Realizations, one per based ring\\\midrule
5&1&$\mathrm{Vec}_{C_5}$\\
4&1&The two Tambara--Yamagami rings, for $C_4$ and $C_2^2$\\
 &3&$\operatorname{Rep}(\operatorname{AGL}(1,5))$\\
 &4&The two Izumi near-group rings, $C_4+4$ and $C_2^2+4$\\
3&2&$\operatorname{Rep}(C_7\rtimes C_3)$, for the faithful action\\
2&1&$\operatorname{Rep}(D_{14})$, $\operatorname{Rep}(S_4)$, $\mathrm{SU}(2)_4$, $\operatorname{Ad}(\mathrm{SU}(2)_8)$, and the Liu--Palcoux--Ren rank-five rings nos.~7 and~11 of dimensions $12$ and $24$ respectively \cite[\S3.5]{LPR}\\
 &3&$C_2$-equivariantization of the Izumi $C_3+6$ category\\
1&1&$\operatorname{Ad}(\mathrm{SU}(2)_9)$\\
 &2&$\operatorname{Rep}(A_5)$ and $\operatorname{Ad}(\mathrm{SU}(3)_4)$\\\bottomrule
\end{tabular}
\caption{The seventeen Grothendieck rings of complex fusion categories of rank five. Here $D_{14}$ has order fourteen, and $\operatorname{Ad}$ denotes the adjoint fusion subcategory.}
\label{tab:models-five}
\end{table}

\subsection*{Model conventions and sources}
The notation $\mathrm{SU}(n)_k$ refers to the unitary quantum-group
fusion category at $q=\exp(\pi i/(k+n))$: the construction and fusion
rules come from semisimplified tilting modules
\cite{AndersenParadowski}, and unitarity from \cite{Wenzl}.
For $\mathrm{SU}(2)_k$, the adjoint subcategory has the even labels
among $0,\ldots,k$. For $\mathrm{SU}(3)_4$, it has highest weights
$(0,0),(3,0),(0,3),(2,2),(1,1)$.

The Tambara--Yamagami models come from \cite{TambaraYamagami};
their $C^*$ realizations are also explicit in
\cite[Corollary~4]{EvansGannon}. The affine-group near-group model
is the case $q=5$ of \cite[\S4.1]{Siehler}.
The two additional multiplicity-one models are constructed in
\cite[\S\S4.2 and~4.4]{LPR}; the multiplicity-four near-group
models are given in \cite[Examples~9.5--9.6]{Izumi}.
The pointed and finite-group representation models carry their
standard unitary structures.

The $C_3+6$ example was first recognized by Liu and Snyder in the
zero-graded part of the $E_9$ quantum $\mathrm{SU}(3)$ subgroup
\cite[\S10, Remark preceding Definition~10.1]{LiuYB};
the concrete $C^*$ realization and symmetry used here are those of
\cite[Theorem~10.19, Lemma~13.1 and Example~13.10]{Izumi}.

\subsection*{The proof in three steps}
\emph{Galois descent.}
We compare the ordinary Grothendieck ring with the signed quotient of
its sphericalization. Their dimension codegrees are $D$ and the
intrinsic global dimension $E$. Central induction makes $E/D$ an
algebraic integer \cite[Propositions~8.21--8.22]{ENO} and supplies integral classes with its conjugate
values. A simultaneous-column identity and Fourier bounds show that
some conjugate of $E$ equals $D$. This proves Theorem~A directly,
including Galois pseudo-unitarity. The field degree is at most five;
the quadratic, cubic, quartic and quintic cases use the same central
identities, with separate arithmetic arguments.

\emph{Bounds on the fusion rules.}
We henceforth use positive spherical dimensions. The published
integral classification \cite{ABDP} and near-group results
\cite{SchNG,MengYu} handle two established families.
For two invertibles, comparing central induction in Morita-equivalent
categories bounds the remaining exchanged-pair family. When the unit
is the only invertible, every nonunit simple generates the category.
A generator of dimension at most two gives the adjoint
$\mathrm{SU}(2)_9$ ring. Otherwise all four nonunit dimensions exceed
two, so $D>17$; character, local-order and twist-moment restrictions
then control each possible dimension field.

\emph{Finite classification.}
The sharp bounds leave a finite enumeration of based rings.
Of 269 rings with $\mu\le4$, exactly 114 also satisfy the dimension
bound. Exact necessary conditions exclude 97, and explicit basis
permutations identify the remaining seventeen with unitary models.
The computations, their finite ranges and their mathematical inputs
are specified in the proofs. The source package contains the
verification programs and certificates.

\subsection*{Relation to earlier work}
Ostrik classified fusion categories of rank two and pivotal fusion
categories of rank three \cite{Ostrik2,Ostrik3}. Hagge and Hong showed
that every pivotal fusion category of rank at most four admits a
spherical structure \cite[Lemma~4]{HH}. In rank four, Larson reduced
the non-self-dual pseudo-unitary case to seven candidate based rings
\cite{Larson}; Edie-Michell, Izumi and Penneys classified
$\mathbb Z/2\mathbb Z$-quadratic unitary fusion categories and completed
the unitary case with a dual pair \cite{EIP}. Dong, Palcoux and Plessis
classified the fifteen Grothendieck rings under the pivotal hypothesis
\cite[Theorem~1.1]{DPP}. The preceding papers in the present series
\cite{Rank2,Rank3,Rank4} treat ranks two through four over $\mathbb C$;
in particular, the rank-four sphericality theorem removes the pivotal
hypothesis from that ring classification.

For categories with braiding, earlier milestones include Ostrik's
premodular rank-three classification \cite{OstrikPremod3}, the unitary
modular classification through rank four of Rowell, Stong and Wang
\cite{RSW}, and Bruillard's rank-four premodular classification up to
Grothendieck equivalence \cite{Bruillard4}. At rank five, Hong and
Rowell treated the non-self-dual modular case \cite{HR}, and Bruillard,
Ng, Rowell and Wang completed the modular classification \cite{BNRW}.
Bruillard and Ortiz-Marrero classified pseudo-unitary premodular
rank-five categories up to Grothendieck equivalence \cite{BO5}.
Without a braiding assumption, prior results include multiplicity-one
rings \cite{LPR}, integral rings \cite{ABDP}, and partial results for
categories with nontrivial faithful grading \cite{WD}.

The dimension-block estimate is due to Schopieray
\cite{SchopierayBlock}. We derive the ordinary estimate and the
ordinary/signed versions used here directly from a common modular
projection bound. The bounded classification \cite{Rank5Bounded}
supplies model data and finite certificates; the accompanying verifier
regenerates the bounded ring list and checks the relevant certificates.

\subsection*{Sources and reproducibility}
The current \href{https://github.com/sebastienpalcoux/grothendieck-ring-classification/blob/main/manuscripts/rank5.tex}{LaTeX source},
\href{https://github.com/sebastienpalcoux/grothendieck-ring-classification/blob/main/manuscripts/rank5/README.md}{reading guide},
and \href{https://sebastienpalcoux.github.io/grothendieck-ring-classification/}{project website}
are publicly available. The structural checks, exact census generator,
certificates, model tensors and recorded outputs are preserved in the
\href{https://github.com/sebastienpalcoux/grothendieck-ring-classification/tree/b479c65c702c6a09a674df7c2f8468b1a48915d3/supplements/rank5/sharp}{verification snapshot}
at commit \texttt{b479c65c702c}. From that directory, run
\texttt{python3 verify.py}; its README gives dependencies and optional
independent checks. The finite verifier establishes the stated arithmetic
and tensor computations; the categorical reductions are proved in the text.


% ===== structure_r5.tex =====
\section{Rank five and proper fusion subcategories}

Let $R=K_0(\mathcal C)$, and let $\mathcal S\to\mathcal C$ be the ENO
sphericalization, with central sign object $\varepsilon$ of dimension $-1$
\cite[Remark~3.1 and Proposition~5.14]{ENO}. Put $T=K_0(\mathcal S)/(\varepsilon+1)$.
Choose one lift of each simple, compatibly with duality, including the unit.
This is possible because the two lifts have nonzero opposite spherical
dimensions, and duality preserves spherical dimension. The coefficient of the
unit gives a trace $\tau$ on either ring with
$\tau(b_i b_j^*)=\delta_{ij}$. In $T$, the two lift multiplicities have
sum $N_{ij}^k$ and difference $T_{ij}^k$. Consequently
\[
 |T_{ij}^k|\leq N_{ij}^k,\qquad T_{ij}^k\equiv N_{ij}^k\pmod 2.
\]
If a change of lift signs makes $T=N$, the selected lifts span a tensor
subcategory on which the forgetful functor is an equivalence. This transports
a spherical structure to $\mathcal C$.

\begin{proposition}\label{prop:rank-five-commutative}
The ordinary and signed rank-five algebras are commutative.
The degree of either dimension character field is at most the number of
self-dual simples of $\mathcal C$. In particular, a rank-five category with
two nontrivial dual pairs is integral and pseudo-unitary.
\end{proposition}
\begin{proof}
Ordinary commutativity is also \cite[Proposition~9.2]{ABDP}. We give one
argument for both algebras. Their complexifications are semisimple, since
the involution makes left multiplication by $b_i^*$ adjoint to left
multiplication by $b_i$ in the orthonormal basis of simples. A noncommutative
semisimple algebra of complex dimension five with a character is
$\mathbb C\oplus M_2(\mathbb C)$. Its unique linear character $\chi$ would
be Galois invariant and integral on the distinguished basis. In $R$ it is
the Frobenius--Perron character and in $T$ it is the spherical dimension
character. In particular all five entries are nonzero integers. Its formal
codegree $h=\sum_i|\chi(b_i)|^2$ would be at least five.

For clarity, write the coefficient-of-unit trace on the simple matrix
blocks as $\tau=\sum_V f_V^{-1}\operatorname{Tr}_V$. Each $f_V$ is an
algebraic integer even for $T$. In fact,
$\sum_i\operatorname{Tr}_V(b_i)b_i^*$ is $f_V$ times the central idempotent
for $V$, by pairing against the basis. All its coefficients are algebraic
integers, since the eigenvalues of $b_i$ on $V$ occur in its integral regular
matrix. Hence $f_V$ is an algebraic integer. On the two proposed blocks,
evaluation at $1$ gives $1=1/h+2/f$, and thus
$f=2+2/(h-1)$ is a nonintegral rational number. This proves commutativity.

The linear extension of $b_i\mapsto b_i^*$ permutes the primitive
idempotents by complex conjugation of their characters. Its trace counts
both self-dual basis elements and real-valued characters. Every Galois
conjugate of the relevant dimension row is fixed by $b_i\mapsto b_i^*$:
this follows from the same equality on the original row, since the
involution is defined over $\mathbb Q$. All such conjugates are real-valued,
which proves the degree bound. With two nontrivial dual pairs the sole
self-dual basis element is the unit, so all Frobenius--Perron dimensions
are rational algebraic integers. Apply \cite[Propositions~8.23--8.24]{ENO}.
\end{proof}

\begin{proposition}\label{prop:proper-integral}
Every proper fusion subcategory of a rank-five complex fusion category is
integral. Consequently a rank-five category with a nontrivial faithful
grading is weakly integral and pseudo-unitary.
\end{proposition}
\begin{proof}
Let $\mathcal D$ be a proper fusion subcategory. Its simple classes act on
the free abelian group on the complementary simple classes, since Frobenius
reciprocity prevents a product of an inside and an outside simple from
containing an inside constituent. At rank four there is a single outside
simple $Y$. Each product $X Y=m_XY$ gives
$\operatorname{FPdim}(X)=m_X\in\mathbb Z$ for $X$ in $\mathcal D$;
cf.\ \cite[Proposition~3.2]{DCW}.

A unital field subalgebra of $R\otimes\mathbb Q$ must have degree dividing
five by restriction of scalars. This excludes a Fibonacci subcategory in
rank two, leaving only the pointed possibility by \cite{Ostrik2}.
In rank three, a non-self-dual based ring is pointed
\cite[Proposition~3.1]{Ostrik3}, so assume all simples are self-dual.
The rational ring of $\mathcal D$ cannot be a field, by the same
divisibility. Hence its codegree polynomial is reducible. If it
splits, the global Frobenius--Perron dimension of $\mathcal D$ is integral,
and \cite{ENO} gives sphericity. If it has a quadratic factor, the
unconditional descent theorem \cite[Theorem~5.1]{Ostrik3} gives sphericity.
We can therefore use Ostrik's pivotal classification
\cite[Theorem~1.1]{Ostrik3}. The rational algebra of the adjoint
$\mathrm{SU}(2)_5$ ring is a cubic field: the spin-one generator has
irreducible characteristic polynomial $t^3-2t^2-t+1$. It is therefore
excluded. The remaining integral ring is $K_0(\operatorname{Rep}(S_3))$. Each of the other two has basis $1,g,x$ and
\[
 g^2=1,\quad gx=x,\quad x^2=1+g+mx,\qquad m\in\{0,2\}.
\]
On the two outside simples let $P$ be the permutation matrix of $g$ and
$M$ the symmetric nonnegative integral matrix of $x$. If $P$ exchanges
them, $PM=MP=M$ makes $M$ a constant matrix with entry $a$, and
$2a^2=1+ma$ is impossible modulo two. If $P=I$, then
$M^2=mM+2I$. Both candidate quadratics are irreducible over $\mathbb Q$,
so $\operatorname{Tr}(M)=m$ and $\det(M)=-2$. Writing the diagonal as
$a,m-a$ gives the square of the off-diagonal entry as $2+a(m-a)$, which
is either two or three. This rules out both cases. Rank one is immediate.

For a nontrivial faithful grading, the neutral subcategory is proper.
The homogeneous components have equal Frobenius--Perron dimensions
\cite[Proposition~8.20]{ENO}, so the total dimension is the product of the
group order and the integer dimension of the neutral subcategory.
Pseudo-unitarity and sphericity follow from \cite{ENO}.
\end{proof}



% ===== central.tex =====
\section{Codegrees and the dimension field}\label{sec:central}
Write $d_i=\operatorname{FPdim}(X_i)$ and choose lifts $X_i^+$ of positive
spherical dimension $x_i$, their negatives being $X_i^-=\varepsilon X_i^+$.
Set
\[
 D=\sum_i d_i^2,\qquad E=\sum_i x_i^2,\qquad L=\mathbb Q(x_i:i).
\]
Thus $E$ is the intrinsic global dimension of $\mathcal C$ and $2E$ is the
global dimension of $\mathcal S$. The $x_i$ and their conjugates are real
and nonzero \cite[Corollary~2.10]{ENO}. Write $f_\chi=\sum_i|\chi(b_i)|^2$ for a formal codegree of
$R$ or $T$, with the usual normalization for $R$
\cite[Section~2.2]{OstrikCodegrees}. In each ring $\sum_\chi 1/f_\chi=1$ by expansion of the
coefficient-of-unit trace. The complex algebra of $\mathcal S$ is the
product of the two complex algebras, and its codegrees are twice those of
the factors. All character values are cyclotomic \cite[Corollary~8.53]{ENO} and $L$ is
totally real and abelian. We recall the central formulas that will be used
with signs and at different embeddings.

Let $F:Z(\mathcal S)\longrightarrow\mathcal S$ be forgetful, with right
adjoint $I$. Each character $\chi$ of either factor indexes a different
simple constituent $A_\chi$ of $I(1)$, with multiplicity one, dimension
$E/f_\chi$ and twist one \cite[Theorems 2.5 and 2.13]{Ostrik3}.
Consequently $f_\chi\in L$, and $E/f_\chi$ is an algebraic integer.
The full set of codegrees is Galois invariant. These assertions hold
when $E$ is replaced by any of its conjugates. Each $f_\chi$ is a
$d$-number, meaning that its principal ideal in the ring of all algebraic
integers is Galois invariant \cite[Theorem~1.2]{OstrikCodegrees}. Indeed the theorem
applies to $2f_\chi$, and division by a rational integer preserves
invariance of fractional ideals.

If $\sigma$ is an embedding of $L$, choose signs on the lifts so that
$\beta_i=|\sigma(x_i)|$ are their dimensions in the transported spherical
data. The signed multiplication matrices are dominated entrywise in
absolute value by the ordinary matrices. Spectral comparison consequently
gives $\beta_i\le d_i$, hence $\sigma(E)\le D$. If equality holds, then
$\beta_i=d_i$ for every $i$. The two dimension equations then yield
$\sum_k(N_{ij}^k-T_{ij}^k)d_k=0$. Each term is nonnegative, so $T=N$
in this choice of lifts and the section descends. The conjugate category is then pseudo-unitary. Transporting its spherical
structure back gives a spherical structure on $\mathcal C$.
Thus the common goal of the following sections is to prove that a Galois
conjugate of $E$ equals $D$.

By M\"uger's center theorem \cite[Theorem~1.2]{MugerCenter},
$Z(\mathcal S)$ is modular, has dimension $(2E)^2$, and has
both Gauss sums equal to $2E$. Let $S$ be its modular matrix
normalized by $2E$. For a
central simple $W$ the central-character formula gives
\begin{equation}\label{eq:modchar}
 S_{A_\chi,W}=\frac{\chi(F(W))}{2f_\chi},\qquad
 S_{0W}=\frac{\dim W}{2E}.
\end{equation}
Here $0$ is the tensor unit of the center.
The central lift $\widehat\varepsilon$ of the sign has dimension and twist
$-1$, and $I(\varepsilon)=\widehat\varepsilon I(1)$. The constituents of
these two induction objects are disjoint, since each constituent of $I(1)$ has positive
dimension and each of its translates has negative dimension.
Balancing in \eqref{eq:modchar} gives
$\theta_{\widehat\varepsilon A_\chi}=\chi(\varepsilon)=\pm1$.
These are the sign and central-induction identities used in
\cite[Section 2]{Rank4}, and their proofs apply in every rank in which both
character algebras are commutative.

\begin{lemma}[Simultaneous columns]\label{lem:simultaneous-columns}
For $\sigma\in\operatorname{Gal}(L/\mathbb Q)$, let $B_\sigma$ be the
constituent of $I(1)$ indexed by the $\sigma$-conjugate of the signed
dimension character. Then every ordinary or signed character satisfies
\begin{equation}\label{eq:simultaneous-columns}
 \chi(F(B_\sigma))=\frac{f_\chi}{\sigma(f_\chi)}.
\end{equation}
In particular,
\[
 \dim F(B_\sigma)=\frac E{\sigma(E)},\qquad
 \operatorname{FPdim}F(B_\sigma)=\frac D{\sigma(D)},\qquad
 \frac ED\le\sigma\!\left(\frac ED\right).
\]
If all signed codegrees outside the dimension-character orbit are
rational, the signed class $U_\sigma$ of $F(B_\sigma)$ is a unit in $T$.
\end{lemma}
\begin{proof}
The modular entry between $B_\sigma$ and $A_\chi$ is
$\sigma(E/f_\chi)/(2\sigma(E))=1/(2\sigma(f_\chi))$.
Symmetry and \eqref{eq:modchar} give
\eqref{eq:simultaneous-columns}. Applying the two dimension characters
gives the displayed ratios. The inequality follows from
$|\dim V|\le\operatorname{FPdim}V$ for objects of the sphericalization.
This proves the simultaneous-column formula of
\cite[Proposition~6.1]{Rank4} in the present setting.

Under the final hypothesis, $U_\sigma$ has value $E/\sigma(E)$ on the
dimension factor and value one on every complementary character. Its
regular determinant is therefore $\operatorname{N}_{L/\mathbb Q}
(E/\sigma(E))=1$. Its multiplication matrix has an integral inverse,
so $U_\sigma^{-1}\in T$.
\end{proof}

We use the universal grading of Gelaki and Nikshych
\cite[Theorem~3.5 and Corollary~3.7]{GN}.
\begin{proposition}\label{prop:field}
If a rank-five category is not weakly integral, its universal grading is
trivial, and both dimension-character fields equal $K=\mathbb Q(D)$.
Every pivotal rank-five fusion category admits a spherical structure.
When its universal grading is trivial, every pivotal structure is spherical.
\end{proposition}
\begin{proof}
The grading assertion follows from Proposition~\ref{prop:proper-integral}.
Put $K_1=\mathbb Q(d_i:i)$. The dimensional grading of
\cite[Proposition~1.8]{GannonSchopieray} is faithful with group
$\operatorname{Gal}(K_1/\mathbb Q(D))$, a quotient of the universal
grading group. Consequently $K_1=\mathbb Q(D)$.
For the signed field, central induction gives $D\in L$. Suppose
$\sigma$ fixes $D$ but acts nontrivially on $L$. Then
$B_\sigma\ne1$, whereas Lemma~\ref{lem:simultaneous-columns} gives
$\operatorname{FPdim}F(B_\sigma)=1$. The forgetful image contains the
unit by adjunction, so $F(B_\sigma)=1$. Write $B=B_\sigma$. A central structure on the unit is
a monoidal natural automorphism $\gamma$ of the identity on $\mathcal S$.
It is nontrivial since $B\ne1$. Its value on $\varepsilon$ is one:
the central character of $B$ at $\widehat\varepsilon$ is $-1$ by
\eqref{eq:modchar}, and balancing evaluates it as
$-\gamma_\varepsilon$.
Thus the scalars of $\gamma$ agree on the two lifts of every simple.
They give a nontrivial grading of $\mathcal C$ by a finite subgroup of
$\mathbb C^\times$, a contradiction.

Finally, for a given pivotal structure the ratios of left to right
simple dimensions define the slope, a monoidal automorphism of the
identity \cite[Section~4]{BV}. A nontrivial slope gives a nontrivial
grading. Thus every pivotal structure is spherical when the universal
grading is trivial. In the remaining case the category is weakly
integral by Proposition~\ref{prop:proper-integral}, so it admits the
canonical spherical structure supplied by pseudo-unitarity.
The distinction is necessary: on $\mathrm{Vec}_{C_5}$, a pivotal
structure given by a character of order five is not spherical.
\end{proof}



In the descent arguments we may first transport the category by a
Galois automorphism so that $E$ is its largest conjugate, and then
choose the positive lifts again. This transport preserves the integer
fusion rules and therefore the distinguished positive numbers $d_i$
and $D$; it changes $E$ to the chosen conjugate. Independently,
$|\tau(d_i)|\le d_i$ for every embedding $\tau$, by the spectral-radius
bound for the integral fusion matrix of $X_i$. Hence $\tau(D)\le D$.
At this embedding, $1/D$ and $1/E$ are both minimal, and
Lemma~\ref{lem:simultaneous-columns}, applied to the transported
category, makes $E/D$ minimal as well.

% ===== hecke_fourier.tex =====
\section{Fourier bounds from the modular representation}\label{sec:hecke}
Schopieray obtains a Dirichlet-character bound for the reciprocal intrinsic global dimension of a spherical category \cite{SchopierayBlock}. We apply the congruence representation of the center of the sphericalization to ordinary and signed induction constituents together.
\begin{lemma}\label{lem:hecke-projection}
Let $\pi$ be a unitary representation of $\mathrm{SL}_2(\mathbb Z/N)$,
let $U(x)=\left(\begin{smallmatrix}1&x\\0&1\end{smallmatrix}\right)$,
$L(x)=U(x)^{\mathsf t}$ and
$S=\left(\begin{smallmatrix}0&-1\\1&0\end{smallmatrix}\right)$.
Suppose that $v$ is a unit vector fixed by $\pi(U(1))$, and that
$\pi(\mathrm{diag}(u,u^{-1}))v=\eta(u)v$ for a Dirichlet character
$\eta$ of conductor $c$. If $P$ projects orthogonally onto the vectors
fixed by $\pi(U(1))$, then $\|P\pi(S)v\|^2\le 1/c$.
\end{lemma}
\begin{proof}
Write $a(b)=\langle v,\pi(L(b))v\rangle$. Averaging the upper unipotent
subgroup and conjugating by $S$ gives
$\|P\pi(S)v\|^2=N^{-1}\sum_{b\bmod N}a(b)$. For $u=1+bx$ a unit,
\[
 U(x)L(b)=L(b/u)\,\mathrm{diag}(u,u^{-1})\,U(x/u).
\]
Since $v$ is fixed by both upper unipotents, $a(b)$ differs from
$a(b/u)$ by the scalar $\eta(u)$ or its complex conjugate, according to
the convention for the inner product. On the stratum $\gcd(b,N)=d$,
every unit $u\equiv1\pmod d$ can be used for every $b$ in that stratum:
the congruence $bx=u-1\pmod N$ is solvable. The map $b\mapsto b/u$
permutes the stratum. If $c\nmid d$, the character is nontrivial on
some such $u$, for otherwise it would factor through $(\mathbb Z/d)^\times$.
It follows that the sum on that stratum is zero. The remaining strata
have $c\mid b$, of which there are $N/c$ residues; $|a(b)|\le1$ proves
the inequality.
\end{proof}

\begin{proposition}\label{prop:induction-fourier}
Let $K=\mathbb Q(D)$ equal the signed dimension-character field and the
ordinary Frobenius--Perron character field, $G=\operatorname{Gal}(K/\mathbb Q)$,
and $d=|G|$. For a Dirichlet character $\eta$ of $G$ of conductor $c$ and
a codegree $f_\rho$, write
$\widehat y_{\rho}(\eta)=\sum_{\sigma\in G}\eta(\sigma)/\sigma(f_\rho)$.
Assume the character algebras are commutative. Then
\[
  \sum_{\rho\ {\rm signed}}|\widehat y_{\rho}(\eta)|^2
   +2\sum_{\rho\ {\rm ordinary}}|\widehat y_{\rho}(\eta)|^2
  \ \le\ 4d/c .
\]
In particular $p=1/D$, $q=1/E$ satisfy
$2|\sum_\sigma\eta(\sigma)\sigma p|^2+
 |\sum_\sigma\eta(\sigma)\sigma q|^2\le 4/c$.
\end{proposition}
\begin{proof}
Use the anomaly-free modular representation of the Drinfeld center of the
sphericalization. It factors through $\mathrm{SL}_2(\mathbb Z/N)$ for
the order $N$ of its twist matrix, and the diagonal matrices act by the
Galois signed permutations, up to replacing the parameter $u$ by $u^{-1}$
\cite[Theorem II and Remark 4.5]{DLN}. This representation is unitary in the simple-object basis:
$T$ is diagonal with root-of-unity entries, and $S$ is symmetric with
$S^2=C$, the duality permutation. Moreover $\overline{S}=CS$, since
modular characters take conjugate values at dual objects and the
normalized vacuum row is real. Hence $S^*S=1$.

For $\sigma\in G$, let $B_\sigma$ be the induction constituent associated
to the $\sigma$-conjugate of the signed spherical character. Its modular
row at a central simple $W$ is precisely
$\sigma(\dim F(W)/(2E))$. Consequently Galois action permutes these rows
through its action on $G$, with positive sign even in a nonunitary
category. The unit vector
$v=d^{-1/2}\sum_{\sigma\in G}\eta(\sigma)e_{B_\sigma}$ has the required
diagonal character, up to inversion; all its coordinates have twist one.
By the central-character formula,
$S_{B_\sigma,A_\rho}=1/(2\sigma(f_\rho))$. The coordinate of $Sv$ at
$A_\rho$ therefore has squared absolute value
$|\widehat y_{\rho}(\eta)|^2/(4d)$. All these constituents have twist
one. If $\rho$ is ordinary, $\widehat\varepsilon A_\rho$ also has twist
one, and lies in a disjoint induction object. Its entry differs only by
a sign: the modular character of $B_\sigma$ takes the value $-1$ on the
central sign. Summing these squared coordinates and applying
Lemma~\ref{lem:hecke-projection} gives the first bound. Within either
full dimension-character orbit the Fourier coefficients have constant
absolute value, which gives the second bound.
\end{proof}

\begin{proposition}\label{prop:parity-fourier}
Under the hypotheses of Proposition~\ref{prop:induction-fourier}, the two
separate bounds
\[
 \frac1d\sum_{\rho\ {\rm signed}}|\widehat y_\rho(\eta)|^2\le
 \frac{\gcd(c,4)}{c},\qquad
 \frac1d\sum_{\rho\ {\rm ordinary}}|\widehat y_\rho(\eta)|^2\le
 \frac{\gcd(c,4)}{c}
\]
hold. In particular, for $c$ odd, each full dimension-orbit Fourier
coefficient has absolute value at most $1/\sqrt c$.
\end{proposition}
\begin{proof}
First use the proof of Lemma~\ref{lem:hecke-projection} with the
unipotent subgroup generated by $U(2)$, both for the vector and the
projection. The average now runs over $b\in2(\mathbb Z/N)$, and $x$
may range over $2(\mathbb Z/N)$. On the stratum $\gcd(b,N)=d_0$ this
allows exactly the units $u\equiv1\pmod{\gcd(2d_0,N)}$ in the same
factorization. Hence only the strata with $c\mid 2b$ can contribute.
Among $b$ in this subgroup their fraction is $\gcd(c,4)/c$: equivalently
sample $b=2j$ with $j$ uniform modulo $N$ and count $c\mid4j$.

Let $J_0=\widehat\varepsilon$. Its modular character ratios
$z_W=S_{J_0,W}/S_{0,W}$ are $\pm1$, since $J_0^2=1$. They satisfy
$z_{J_0 W}=z_W$: the scalar double braiding of $J_0$ with itself is
$\theta_{J_0^2}/\theta_{J_0}^2=1$. At $A_\rho$ one has
$z_{A_\rho}=\rho(\varepsilon)$ by the central-character formula.
The modular relation for an invertible simple reads
$S_{J_0B,W}=z_W S_{B,W}$. Because $z_W$ is rational, the Galois action
on the rows indexed by $J_0B_\sigma$ is again the unsigned permutation
of $G$. Both unit vectors
\[
 v_\pm=(2d)^{-1/2}\sum_\sigma\eta(\sigma)
    (e_{B_\sigma}\pm e_{J_0B_\sigma})
\]
are fixed by $T^2$ and have diagonal character $\eta$ up to inversion.
At $A_\rho$, the coordinate of $Sv_\pm$ is proportional to
$(1\pm z_{A_\rho})\widehat y_\rho(\eta)$; it vanishes in the
unselected sector and equals $\widehat y_\rho(\eta)/\sqrt{2d}$ in the
selected sector. At $J_0 A_\rho$ it has the same absolute value, since
$z_{J_0 A_\rho}=z_{A_\rho}$ and $z_{B_\sigma}=-1$. All these coordinates
belong to distinct simples with squared twist one. Applying the subgroup
bound to $v_-$ and $v_+$ proves the assertions.
\end{proof}




\begin{proposition}[Ordinary Fourier bound]
\label{prop:ordinary-fourier}
Let $\mathcal C$ be a fusion category with positive spherical dimensions
and commutative Grothendieck ring. Suppose
$K=\mathbb Q(d_i:i)=\mathbb Q(D)$, and write
$G=\operatorname{Gal}(K/\mathbb Q)$ and $d=|G|$.
For a character $\eta$ of $G$ of Dirichlet conductor $c$, put
\[
 \widehat y_\rho(\eta)=
 \sum_{\sigma\in G}\frac{\eta(\sigma)}{\sigma(f_\rho)}.
\]
Then
\[
 \frac1d\sum_\rho|\widehat y_\rho(\eta)|^2\le\frac1c,
 \qquad
 \left|\sum_{\sigma\in G}\frac{\eta(\sigma)}{\sigma(D)}\right|
 \le\frac1{\sqrt c}.
\]
\end{proposition}
\begin{proof}
Use the center of $\mathcal C$ itself, with modular matrix normalized
by $D$. Each character $\rho$ indexes a distinct twist-one induction
constituent $A_\rho$ of dimension $D/f_\rho$. In particular
$f_\rho\in K$. The field $K$ is abelian by cyclotomicity.
Let $B_\sigma$ correspond to the $\sigma$-conjugate of the dimension
character. The ordinary version of the simultaneous-column identity is
\[
 S_{B_\sigma,A_\rho}=\frac1{\sigma(f_\rho)}.
\]
These $d$ labels are distinct and have twist one. Their rows are the
Galois conjugates of the vacuum row, with positive vacuum entries
$1/\sigma(D)$; consequently the diagonal Galois matrices permute
them without signs. If $N$ is the twist order of the center, inflate
$\eta$ along $\operatorname{Gal}(\mathbb Q(\zeta_N)/\mathbb Q)\to G$.
The vector
\[
 v=d^{-1/2}\sum_{\sigma\in G}\eta(\sigma)e_{B_\sigma}
\]
is a unit vector fixed by $T$ and has diagonal character $\eta$ or its
inverse. At $A_\rho$, its $S$-transform has squared absolute value
$|\widehat y_\rho(\eta)|^2/d$. Summing these twist-one coordinates and
applying Lemma~\ref{lem:hecke-projection} proves the first inequality.
The $d$ characters in the dimension orbit have Fourier coefficients
of the same absolute value, which proves the second.
\end{proof}


% ===== small_degree_descent.tex =====
\section{Galois descent in degrees two and three}
\label{sec:cubic-permutation}

Write $K$ for the common ordinary and signed dimension-character field,
assumed totally real and Galois, $d=[K:\mathbb Q]$, and
$q=1/E$, $t=\operatorname{Tr}_{K/\mathbb Q}(q)$,
$s=\operatorname{Tr}_{K/\mathbb Q}(1/D)$.
Characters outside the Galois orbit of the chosen dimension character
will be called complementary in this section. We use a commutative
ordinary and signed algebra, as already proved for rank five. For an
ordinary or signed character $\chi$, write $f_\chi$ for its codegree in
that algebra, $A_\chi$ for its induction constituent, and $F$ for the
forgetful functor from the center of the sphericalization. In particular
$A_D$ denotes the constituent for the ordinary Frobenius--Perron row.

\begin{lemma}\label{lem:rational-E-descent}
Suppose every complementary signed codegree is rational. Then $s\le t$.
Under the same hypothesis on the complementary codegrees, rationality of $E$ gives $E=D$, in every degree.
\end{lemma}
\begin{proof}
For every signed complementary character $\rho$, the modular character
formula and symmetry of the $S$-matrix give
\[
 \rho(F(A_D))=f_\rho\operatorname{FPdim}(F(A_\rho))/D>0 .
\]
On each Galois orbit of complementary characters, these values run through
all conjugates of a totally positive algebraic integer with a common
multiplicity. Their arithmetic mean is at least one, by their positive
integral norm. This argument applies to nonreal character rows as well:
their displayed values on this particular forgetful object are positive,
including all Galois conjugates. The weight $1/f_\rho$ is constant on
each orbit. The unit coefficient of the signed class of $F(A_D)$ is one,
and its field contribution is
$\operatorname{Tr}((E/D)/E)=s$. Thus
\[
 1-s=\sum_{\rho}\rho(F(A_D))/f_\rho
       \ge\sum_{\rho}1/f_\rho=1-t .
\]
If $E$ is rational this gives $\operatorname{Tr}(E/D)=Es\le Et=d$.
Since $E/D$ is a totally positive algebraic integer by central induction,
its trace is at least $d$ by the arithmetic--geometric mean inequality.
Equality forces every conjugate to be one. Hence $E=D$, giving pseudo-unitarity.
\end{proof}

\begin{theorem}\label{thm:cubic-full}
Suppose $d\in\{2,3\}$ and every complementary signed codegree is rational.
Then some Galois conjugate of the category is pseudo-unitary.
In particular this holds for every rank-five complex fusion category
whose global Frobenius--Perron dimension has degree three.
\end{theorem}
\begin{proof}
By Lemma~\ref{lem:rational-E-descent} we may assume $E$ irrational.
Write $j_-$ for passage to the signed quotient on Grothendieck rings.
For $\sigma\in\mathrm{Gal}(K/\mathbb Q)$, let $B_\sigma$ and
$U_\sigma=j_-F(B_\sigma)$ be as in
Lemma~\ref{lem:simultaneous-columns}. That lemma gives
\begin{equation}\label{eq:ratio-minimal}
 \alpha:=E/D\le\sigma(\alpha)
 \qquad\text{and}\qquad U_\sigma^{-1}\in T.
\end{equation}
The value of $U_\sigma$ in $K$ is $E/\sigma(E)$; its value at every
complementary character is one.

Let $C_\sigma$ be the ordinary induction constituent for the
$\sigma$-conjugate Frobenius--Perron character. The unit coefficient
of $V_\sigma=j_-F(C_\sigma)$ is one, since $F(C_\sigma)$ contains
the positive unit once and the negative unit zero times. Its value in
$K$ is $E/\sigma(D)$. The field contribution to its unit coefficient
is therefore $s$, and its complementary contributions sum to $1-s$.
Multiplication by $U_\sigma^{-1}$ leaves these contributions unchanged
and changes the value in $K$ to $\sigma(\alpha)$. The unit coefficient
of the resulting integral class is an integer. Since $s\le1$ and the
trace below is positive,
\begin{equation}\label{eq:shifted-trace}
  \operatorname{Tr}_{K/\mathbb Q}(q\sigma(\alpha))
        \in s+\mathbb Z_{\ge0}
       \qquad\text{for every }\sigma .
\end{equation}
Equality holds at the identity.

Choose the spherical embedding to make $E$ its largest conjugate; the
ordinary Frobenius--Perron character is unchanged by Galois transport.
Order the embeddings by $q_0<q_1$ when $d=2$ and $q_0<q_1<q_2$ when
$d=3$, with $q_0=1/E$. These are strict because either degree is
prime and $E$ is irrational. If $\alpha$ is rational, it is a positive
algebraic integer at most one and thus equals one. Otherwise \eqref{eq:ratio-minimal} gives
$\alpha_0<\alpha_i$ for every $i\ne0$. In degree two, the nonidentity
shift in \eqref{eq:shifted-trace} gives a trace $T$ with
\[
 s-T=(\alpha_1-\alpha_0)(q_1-q_0)>0 .
\]
In degree three both three-cycles are available. For the shift
$T=q_0\alpha_1+q_1\alpha_2+q_2\alpha_0$ one obtains
\[
 s-T=(\alpha_1-\alpha_0)(q_1-q_0)
       +(\alpha_2-\alpha_0)(q_2-q_1)>0 .
\]
Both inequalities contradict \eqref{eq:shifted-trace}. Hence $E=D$ at
the chosen embedding, proving the assertion.

Finally, when $D$ has degree three in rank five, both dimension fields
contain $D$ and have degree at most five, so they equal $K=\mathbb Q(D)$.
This field is totally real, cyclotomic and hence cyclic. Each
complementary signed character has degree at most two. Its codegree
belongs both to its character field and, by central induction, to
$K$, and is therefore rational. The general statement applies.
\end{proof}



% ===== quadratic.tex =====
\section{Galois descent in quadratic degree}\label{sec:quadratic}
\begin{theorem}\label{thm:quadratic-full}
If $D$ has degree two over $\mathbb Q$, then some Galois conjugate of
$\mathcal C$ is pseudo-unitary.
\end{theorem}
Suppose no Galois conjugate is pseudo-unitary.
Proposition~\ref{prop:field} identifies both dimension-character fields
with $K=\mathbb Q(D)$. Choose the embedding with $E$ largest.
Then $E<D$.

If there are two nontrivial dual pairs, Proposition~\ref{prop:rank-five-commutative} gives integrality. With one such pair, the signed ring has three real characters, of which two form the quadratic dimension orbit. The third real character is rational, since reality is the Galois-stable condition $\chi(b_i)=\chi(b_i^*)$ for every $i$. The two other characters generate an imaginary quadratic field. Their common codegree lies both in that field and in $K$, so is rational. Theorem~\ref{thm:cubic-full} gives a pseudo-unitary conjugate, a
contradiction.

We may therefore assume that every simple is self-dual. The tensor of coefficients on each ring is then fully symmetric by reciprocity. Theorem~\ref{thm:cubic-full} also shows that some complementary signed codegree is irrational. It belongs to $K$ and to its character field, of degree at most three, so that character field is another copy of $K$. The signed rational algebra is consequently $K\times K\times\mathbb Q$. Write $E,E'$ for the codegrees of the signed dimension pair, $G,G'$ for those of the other quadratic pair, and $g$ for the rational codegree. We already have $E\ge E'$, and label the other two characters so that
$G>G'$. We first exclude $g=2$, then bound $E$ and test the remaining codegrees.



% ===== no_signed_two.tex =====
\subsection{A rational signed character of codegree two}
\begin{theorem}\label{thm:no-signed-two}
Suppose that the ordinary and signed dimension-character fields are quadratic, and that the three complementary signed characters form a rational singleton and a quadratic pair. Then the rational singleton cannot have codegree two.
\end{theorem}
\begin{proof}
Suppose $\psi$ is the rational singleton of codegree two. Let $A=A_\psi$ be its central induction constituent in the center of the sphericalization, and let $s=\widehat\varepsilon$ be the central sign. The normalized modular matrix has a sparsely supported rational row
\[
 \delta_B=4S_{A,B}=\psi(FB)\in\mathbb Z,\qquad
 \sum_B\delta_B^2=16,\qquad \delta_{sB}=-\delta_B .
\]
Here $A$ is self-dual since its row is real, and $S^2$ is charge conjugation. Also $\theta_A=1$. For each character $\chi$ of either rank-five algebra, write $f_\chi$ for its codegree in that algebra and $A_\chi$ for its central label. Modular symmetry gives
\[
              \delta_{A_\chi}=2\chi(F(A))/f_\chi .
\]
The integer entries on these labels are consequently constant on Galois orbits. Their sum is two on each side, since $\tau(F(A))=1$ in each quotient. The two signed dimension labels both have entry one, as the unit does. The two ordinary dimension labels have the same strictly positive integer entry by positivity of the Frobenius--Perron dimension of $F(A)$. The $A_\chi$ and their $s$-translates are disjoint, so the norm bound forces this positive integer to be one as well. It leaves a budget of four for the sum of squares of the six complementary label entries before their translates. The signed three have sum zero and two are Galois conjugate, so they must be $(0,0,0)$: a nonzero triple $(b,b,-2b)$ would cost six. In particular $\delta_A=0$.

Let $M_B=\pm1$ be monodromy with $s$. Then $\theta_{sB}=-M_B\theta_B$. The modular rows at $A$ and $sA$ are orthogonal, so the squared weights $\delta_B^2$ have mass eight in each monodromy sector. The center has central charge one, and the modular relation $STS=T^{-1}ST^{-1}$ at $(A,A)$ gives $\sum_B\delta_B^2\theta_B=0$. The sector $M_B=+1$ cancels in sign pairs. In $M_B=-1$ the two ordinary dimension labels and their translates already contribute four at twist one. Since the total weight of the sector is eight, all its remaining nonzero weight is at twist minus one. In particular the four known labels exhaust its nonzero twist-one entries.

Now use the same modular relation at $(A,0)$, retaining the signs of the spherical dimensions. Put $v_B=\delta_B\dim B$. Orthogonality of the row at $A$ with the rows at $0$ and $s$ gives $\sum_{M_B=-1}v_B=0$ and $\sum_{M_B=+1}v_B=0$ separately. The two members of every sign pair have equal $v_B$, so $M_B=+1$ contributes zero to $\sum_B v_B\theta_B$. Write $p=1/D+1/D'$. In $M_B=-1$ the known twist-one labels contribute $2Ep$: their dimensions before translation are $E/D$ and $E/D'$, and their entries are one. The remaining $v_B$ in this sector sum to $-2Ep$ and all have twist minus one. It follows that
\[
 4Ep=\sum_B v_B\theta_B
    =8E(STS)_{A,0}=8E S_{A,0}=2E,
 \qquad p=\tfrac12 .
\]

Finally let $C$ be the ordinary positive dimension induction constituent (codegree $D$). The coefficient of the unit in $j_-F(C)$ is one. Its two signed dimension characters contribute $p$ and its rational character contributes $\psi(F(C))/2=1/2$. Every complementary character $\rho$ contributes a strictly positive number, since
\[
 \frac{\rho(F(C))}{f_\rho}
     =\frac{\operatorname{FPdim}(F(A_\rho))}{D}>0
\]
by modular symmetry. This contradicts the coefficient-of-unit identity and completes the proof.
\end{proof}



% ===== quadratic_units.tex =====
\subsection{A unit test on the complementary quadratic factor}\label{sec:quad-units}
Put  $c=\operatorname{disc}(K)$ and $s=\operatorname{Tr}(1/D)$, $t=\operatorname{Tr}(1/E)$ and $W=\operatorname{Tr}(1/G)=1-1/g-t$. Write $\alpha=E/D$ and $r=E'/D'$, where primes denote conjugation in $K$. The totally positive algebraic integer $\alpha$ satisfies $0<\alpha<1$ and $r\ge\kappa=(3+\sqrt5)/2$: its trace and norm are positive integers, the smaller conjugate is below one, and $r$ is smallest at trace three and norm one. More explicitly, putting $k_0=\alpha r$, $k_0=1$ gives integral trace at least three; $k_0=2$ gives trace at least four since trace three would have roots one and two; and $k_0\ge3$ gives $r>3$. Writing $d=1/D$, $e=1/E$, we have $s=\alpha e+r(t-e)\ge r(t-e)$ and $2s\ge(\alpha+r)t\ge3t$ because $1/E\le1/E'$. In $K=\mathbb Q(\sqrt2)$, $r=a+b\sqrt2$ with $b$ a positive integer and $0<a-b\sqrt2<1$, so $r\ge2+\sqrt2$.

Comparison of the discriminant of the signed order with that of its maximal order, $c^2$, shows
$g\operatorname{N}(E)\operatorname{N}(G)$ is a square. Hence $g$ is
a rational norm from $K$, since $\operatorname{N}(E/G)/g$ is a rational
square. For $c=5$, neither $2$ nor $3$ is a norm since both are inert,
and $6$ is excluded as well. For $c=12$, neither $2$ nor $3$ is
a norm: the equations $a^2-3b^2=2m^2$ and $a^2-3b^2=3m^2$ contradict
primitivity by reduction modulo three and division. For $c=8$, three
is inert. Thus $g\ge4$ at $c=5,12$ and also $g\ge4$ at
$c=8$, by Theorem~\ref{thm:no-signed-two}.

Evaluate the three complementary signed characters on the ordinary
Frobenius--Perron induction constituent. Denote the resulting positive
algebraic integers by $\beta,\beta',\gamma$, with $\gamma$ rational.
The unit coefficient gives
\[
 s+\beta/G+\beta'/G'+\gamma/g=1.
\]
The strict inequality $s>t$ (quadratic comparison) requires
$\beta/G+\beta'/G'<W$. Since $\beta\beta'$ is a positive integer, $\beta+\beta'\ge2$. The smaller conjugate must occur at the larger weight $1/G'$, for reversing this order would make the weighted sum at least $W$. It follows that $0<\beta'<1<\beta$.
The ratio $G/G'$ is a positive norm-one algebraic unit by the
$d$-number property. Write it as $\epsilon^j$ with $j\ge1$, where
the smallest norm-one totally positive units of the three fields are
respectively $(3+\sqrt5)/2,3+2\sqrt2,2+\sqrt3$. These values follow by
testing the least admissible integral trace in the Pell equation.

We record the small unit tests explicitly. If
$G/G'=(U+V\sqrt c)/2$ and $\beta=(m+n\sqrt c)/2$ with
$m\equiv nc\pmod 2$, the preceding strict inequalities imply
\[
  n\sqrt c<m<n\sqrt c+2,\qquad mV-(U-2)n<2V.
\]
Thus $0<n<2V/(\sqrt{U^2-4}-U+2)$. At $j=1$, substituting
$(U,V)=(3,1),(6,2),(4,1)$ gives no possibility. At $j=2$, one has
$G=M\epsilon$ for a positive integer $M$ since $G/\epsilon$ is a
rational algebraic integer. At $c=5$ the sole permissible $\beta$
is $\epsilon$. At $c=8$, substitution of $(U,V)=(34,12)$ gives the
four pairs $(m,n)=(4,1),(6,2),(12,4),(20,7)$, for which
$J=\operatorname{Tr}(\beta/\epsilon)\in\{2,4\}$.




% ===== quadratic_120.tex =====
\begin{proposition}\label{prop:quadratic-small-conductor-120}
Suppose the common dimension field is quadratic of conductor $c$.
If no Galois conjugate is pseudo-unitary, then
$\max(E,E')\le120$.
\end{proposition}
\begin{proof}
Use the notation and quadratic-integer tests of
Section~\ref{sec:quad-units}. By
Theorem~\ref{thm:no-signed-two}, $g\ge3$; the norm obstructions
in Section~\ref{sec:quad-units} give $g\ge4$ when
$c\in\{5,8,12\}$. Take $E\ge E'$, and write
$s=\operatorname{Tr}(1/D)$, $t=\operatorname{Tr}(1/E)$,
$W=1-1/g-t$, $e=1/E$, $d=1/D$, $\alpha=E/D<1$,
$r=E'/D'\ge k$, where $k=(3+\sqrt5)/2$ in all conductors
except eight, for which we use $k=2+\sqrt2$. The differences in the
two dimension packets satisfy the exact comparison
\[
       s-2d=r(t-2e)+(r-\alpha)e\ge k(t-2e).
\]
Let $j$ be the positive exponent of the totally positive fundamental
unit in $G/G'$, with $G>G'$, and
$L=((G-G')/(G+G'))^2$; the small-unit tests at conductors $5,8,12$ give $j\ge2$ there.
Suppose $e<1/120$.

First suppose $c$ is outside $\{5,8,12\}$. The trace
$U=\operatorname{Tr}(G/G')$ satisfies $U^2-4=cV^2$ for a nonzero
integer $V$, and $L=(U-2)/(U+2)$. If $c$ is odd then $c\ge13$,
$U\ge5$ and $L\ge3/7$. Proposition~\ref{prop:parity-fourier}
would give
\[
 1/13\ge(t-2e)^2+LW^2
 \ge \tfrac3{10}(2/3-1/60)^2>1/13 .
\]
If $c$ is even then $c\ge24$, $U\ge6$, and $L\ge1/2$.
Proposition~\ref{prop:induction-fourier} would give
\[
 1/6\ge (2k^2+1)(t-2e)^2+LW^2
 \ge \frac{2k^2+1}{4k^2+3}(2/3-1/60)^2>1/6.
\]
Both lower bounds use $t+W\ge2/3$ and completion of a square.

For $c=5$ Proposition~\ref{prop:parity-fourier} gives the signed bound
$(t-2e)^2+LW^2\le1/5$. If $j\ge3$, then $L\ge4/5$, whence its
left side is at least
$\frac49(3/4-2/120)^2>1/5$ by completing a square. If $j=2$,
put $G=M\epsilon$ with $M$ a positive integer. The same unit tests give
$t=1-1/g-3/M$ and $s=1-\gamma/g-2/M$, $\gamma\ge1$.
Hence $s-t\le1/M$ and $s\ge k(t-e)$ imply
\[
 M\big[3(k-1)/4-k/120\big]\le1+3(k-1),
\]
which forces $M\le4$. But $LW^2=5/M^2>1/5$.

For $c=8,12$, Proposition~\ref{prop:induction-fourier} gives
$2(s-2d)^2+(t-2e)^2+LW^2\le4/c$. Using the exact comparison and
completing a square gives the necessary inequality
\[
  \frac{(2k^2+1)L}{2k^2+1+L}
               (1-1/g-2/120)^2 < 4/c .
\]
At $c=12$, $j\ge2$ gives $L\ge3/4$, and the left side is
greater than $1/3$ already at $g=4$. At $c=8$, $j\ge3$ gives
$L\ge49/50$, and the left side is greater than $1/2$ at $g=4$.
For $j=2$, $L=8/9$, and the left side is greater than $1/2$
if $g\ge5$. These strict numerical comparisons follow already from
$k>13/5$ and $k>17/5$ at conductors twelve and eight, respectively.

The remaining case has $c=8$, $j=2$, $g=4$. Put $G=M\epsilon$.
Here $t=3/4-6/M$ and $s=1-\gamma/4-J/M$, $J\in\{2,4\}$.
Thus $s-t\le4/M$. As above $s\ge k(t-e)$ gives
$M[3(k-1)/4-k/120]\le4+6(k-1)$, hence $M\le10$, while
$t>0$ gives $M\ge9$. For $M=10$,
$2k^2(3/20-2/120)^2+32/100>1/2$ contradicts the mixed
inequality. For $M=9$, we have $t=1/12$ and
$s\ge k(1/12-1/120)>1/4$. Among the possible values
$1-\gamma/4-J/9$, the least value above $1/4$ is $5/18$,
attained at $(\gamma,J)=(2,2)$. Since $d<1/120$, one obtains
\[
 2(5/18-2/120)^2+32/81=\frac{8609}{16200}>\frac12,
\]
again impossible.
\end{proof}



% ===== quadratic_small.tex =====
\subsection{The finite quadratic calculation}
\begin{proof}[Completion of the proof of Theorem~\ref{thm:quadratic-full}]

Suppose no Galois conjugate is pseudo-unitary. By the preceding reductions it
suffices to exclude a signed rational algebra $K\times K\times\mathbb Q$
with $E\ge E'$, irrational $G>G'$, $g\ge3$ and $E\le120$.
We give the entire finite calculation.
Represent any quadratic integer by a pair $(h,b)$, with value
$(h+b\sqrt c)/2$, $h\equiv bc\pmod2$, and norm
$n=(h^2-cb^2)/4$. For a full pair of conjugate characters with common
field $K$, both conjugates of its formal codegree $H$ exceed one.
Moreover $c\mid h$: writing the real character values as
$(a_i+b_i\sqrt c)/2$, orthogonality of the conjugate rows gives
$\operatorname{Tr}(H)=c\sum_i b_i^2$. This also applies when the
codegree of the conjugate pair is rational. The $d$-number condition
implies $n\mid h^2$, since $h^2/n=H/H'+H'/H+2$ is a rational
algebraic integer.

Enumerate positive fundamental discriminants $5\le c\le240$; a larger
$c$ cannot divide $\operatorname{Tr}(E)\le240$. For positive multiples
$h_E$ of $c$ up to $240$, use the integers
$0\le b_E\le\lfloor h_E/\sqrt c\rfloor$, the preceding tests, and
$E\le120$. This produces $331$ possibilities, including rational $E$.
For each one, enumerate $3\le g\le\sqrt{n_E}$ with $g\mid h_E$ and
$g^2\mid n_E$, necessary because $E/g$ is integral.
For each $n_G\mid n_E$, put
\[
 h_G=n_G(1-1/g-h_E/n_E)
\]
by the unit-coefficient identity. Require a positive integral $h_G$
and a positive integral $b_G$ with $h_G^2-4n_G=c b_G^2$. Apply the
codegree tests and require both $E/G,E/G'$ integral, equivalently
here $2n_G\mid h_E h_G\pm c b_E b_G$. The order discriminant requires
$g n_E n_G$ square. Exactly $50$ profiles remain.
For $X=E,G,D$ write $v_X=c b_X^2/n_X^2$; this is
$(1/X-1/X')^2$. Proposition~\ref{prop:parity-fourier} requires
$v_E+v_G\le\gcd(c,4)/c$, leaving $13$ profiles.

Finally take $n_D\mid n_E$ and $0<h_D\le n_D$ with $c\mid h_D$;
the upper bound follows from $(D-1)(D'-1)>0$. Require
$h_D^2-4n_D=c b_D^2$ with $b_D>0$, the codegree tests, $D>E$ and
$2n_D\mid h_Eh_D-cb_Eb_D$, since $E/D$ is integral.
This gives $22$ tuples. Propositions~\ref{prop:parity-fourier}
and~\ref{prop:induction-fourier} require
$v_D\le\gcd(c,4)/c$ and $2v_D+v_E+v_G\le4/c$. The sole survivor is
\[
 (c,(h_E,b_E),(h_G,b_G),g,(h_D,b_D))
     =(5,(20,4),(10,2),4,(40,16)).
\]
Here $G/G'=(3+\sqrt5)/2$ is the first positive power of the
fundamental totally positive unit, excluded by the small-unit test in
Section~\ref{sec:quad-units}.
The integer and rational arithmetic just described is reproduced by
\texttt{verify\_quadratic\_120.py}.
\end{proof}



% ===== quartic.tex =====
\section{Galois descent in quartic degree}\label{sec:quartic}
\begin{theorem}\label{thm:quartic-full}
If $D$ has degree four over $\mathbb Q$, then some Galois conjugate of
$\mathcal C$ is pseudo-unitary.
\end{theorem}
The field $K$ is totally real and abelian. Its Galois group is the Klein four group or the cyclic group of order four. The two cases are settled in Theorem~\ref{thm:biquartic-descent} and Proposition~\ref{prop:cyclic-quartic-complete}.
\subsection{The two rational characters}
\begin{proposition}\label{prop:quartic-conductor}
Suppose $D$ has degree four. Let $F$ be the smallest conductor among quadratic subfields $k$ of $K=\mathbb Q(D)$ and $J$ the smaller conductor of the two characters outside the character group of $k$. The ordinary and signed rational codegrees satisfy $4\le f\mid g$. For $(r,H)=(f,D),(g,E)$ we have $\operatorname{N}(H)\ge r\operatorname{disc}(K)$ and $\operatorname{disc}(K)\ge\max(F^3,FJ^2)$.
\end{proposition}
\begin{proof}
The dimension fields both equal $K$ by degree, so both rational algebras
are $\mathbb Q\times K$. All simples are self-dual. Let $g$ be the
rational signed codegree. For the signed rational character $\rho$ and an
ordinary induction constituent $A$ from the quartic factor, coefficient
of the unit gives
\[
 1=\operatorname{Tr}_{K/\mathbb Q}(1/D)+\rho(F(A))/g
   =1-1/f+\rho(F(A))/g .
\]
Indeed $F(A)$ has the positive unit once and the sign zero times. Thus
$f\mid g$. Both are at least two. We first show that $g=2$ is impossible. If $A$ is the signed rational central label, the integers $\delta_B=4S_{A,B}$ satisfy $\sum_B\delta_B^2=16$ and have opposite signs on the translates by $\widehat\varepsilon$, as in Theorem~\ref{thm:no-signed-two}. On the four signed conjugates of the dimension label they are one, since the entry at the unit is one. On the five signed labels their sum is two by the coefficient-of-unit formula applied to $F(A)$. The fifth entry is thus $-2$. Those five entries and their translates already exhaust the squared norm. Yet the ordinary Perron label has strictly positive entry by modular symmetry. This excludes $g=2$.

In particular a nonidentity invertible $h$ cannot exist in quartic degree. The kernel of the ordinary dimension map is one-dimensional and contains $1-h$. A second such invertible, or an exterior simple moved by $h$, would give a linearly independent kernel vector. Therefore $h$ has order two and fixes each of the three exterior simples. Its positive lift fixes their positive lifts as well and interchanges the unit and $h$. The projector $(1-h)/2$ in the signed ring has rank one and yields a rational character of codegree two, which was excluded. If $f=2$, the ordinary rational row would have exactly one nonunit entry $\pm1$. Orthogonality to the positive row forces that entry to be $-1$ on an invertible. Thus $f\ge3$.

For an order $A\subset\mathbb Z\times\mathcal O_K$ of the present
product type, let $r=f$ or $g$, and let $H=D$ or $E$. Its trace
discriminant is $r\operatorname{N}(H)$. The rational primitive
idempotent has basis coefficients $\rho(b_i)/r$, including $1/r$ at the
unit, so $A\cap(\mathbb Q\times0)=r\mathbb Z\times0$. The index of
$A$ in the maximal product order is at least $r$. Thus
\[
 \operatorname{N}(H)\ge r\operatorname{disc}(K)\ge rF^3 .
\]
The last inequality follows from the conductor--discriminant formula:
every nontrivial character either is quadratic or has quadratic square
\cite{Was}.

The four characters have conductors $1,F,J,J'$ with $J'\ge J$ and $J\ge F$, which also gives $\operatorname{disc}(K)\ge FJ^2$.

It remains to exclude $f=3$. At the prime three, ramification in $K$ is tame. Its ramification index $e$ is one or two: total ramification of index four would require a quotient of order four of $\mathbb F_3^\times$, since $K/\mathbb Q$ is abelian. The exponent of three in $\operatorname{disc}(K)$ is therefore zero or two. Since $D$ is a $d$-number, its valuations at primes above three agree. The exponent of three in $\operatorname{N}(D)$ is a multiple of $4/e$, hence even. The equality of the order discriminant $f\operatorname{N}(D)$ with $\operatorname{disc}(K)$ times a square forces the exponent of three in $f$ to be even. Thus $f\ge4$.
\end{proof}



% ===== quartic_fourier.tex =====
\begin{corollary}\label{cor:quartic-effective-conductor}
In the notation of Proposition~\ref{prop:quartic-conductor}, $F<40$ and
$J<729$. In particular, the conductor of $K$ is at most $39\cdot 728$.
\end{corollary}
\begin{proof}
Let $p_* ,q_*$ be the minima of the reciprocal quartets and $s,t$ their
sums, which are at least $3/4$ since $4\le f\mid g$. The discriminant
bound $\operatorname{N}(D)\ge f\operatorname{disc}K$ and its signed
analogue give $p_*,q_*\le (4F^3)^{-1/4}$. Apply Fourier inversion at a
minimum in each quartet and then the triangle inequality in $\mathbb C^2$:
\[
 \big\|(\sqrt2(s-4p_*),t-4q_*)\big\|_2
 \le \sum_{\eta\ne1}
      \big(2|\widehat p(\eta)|^2+|\widehat q(\eta)|^2\big)^{1/2}
 \le 6/\sqrt F .
\]
The two minimizing embeddings may differ; the Fourier inversion factors
have modulus one in both coordinates. For $F\ge40$ this contradicts
$\sqrt3\,(3/4-4(4\cdot40^3)^{-1/4})>6/\sqrt{40}$ and monotonicity.

Now $H=\operatorname{Gal}(K/k)$ has two cosets. For either real list
$y$, character orthogonality gives
\[
 \sum_{\eta|_H\ne1}|\widehat y(\eta)|^2
   =2\sum_{\{i,i'\}\text{ an }H\text{-coset}}(y_i-y_{i'})^2.
\]
The two characters in this sum have conductors at least $J$.
Proposition~\ref{prop:induction-fourier} therefore bounds every
within-pair difference, in both lists, by $2/\sqrt J$.
The discriminant bound $\operatorname{disc}K\ge FJ^2$ also gives
$p_*,q_*\le(4FJ^2)^{-1/4}<1/(2\sqrt J)$, since $F\ge5$.
Whichever pair contains the minimum of its list has sum at most
$3/\sqrt J$. The character of $k$ distinguishes the pair cosets, so
both its Fourier coefficients have absolute value at least
$3/4-6/\sqrt J$. If $J\ge729$, this is at least $19/36$, contrary to
$3(19/36)^2>4/5\ge4/F$ and Proposition~\ref{prop:induction-fourier}.
The characters of conductors $F$ and $J$ generate the character group,
so the conductor of $K$ divides $\operatorname{lcm}(F,J)$.
\end{proof}



% ===== residual.tex =====
\subsection{A stabilizer of the dimension ratio}
\begin{proposition}\label{prop:stab}
Suppose that $D$ has degree four. If a Galois conjugate of $E$ divided
by $D$ belongs to a proper subfield of $K=\mathbb Q(D)$, then some Galois
conjugate of the category is pseudo-unitary.
\end{proposition}
\begin{proof}
Transport the category by Galois action so the given quotient is
$E/D$. Its proper subfield has a nonidentity stabilizer element
$\sigma\in\mathrm{Gal}(K/\mathbb Q)$. For the signed dimension
constituent $B_\sigma$, Lemma~\ref{lem:simultaneous-columns} gives
\[
 \dim F(B_\sigma)=E/\sigma E=D/\sigma D=\operatorname{FPdim}F(B_\sigma)>1.
\]
The last inequality uses that $D$ is the maximum among its conjugates,
and that these four conjugates are distinct. The object $F(B_\sigma)$
contains the positive unit once and does not contain the negative
unit. At least one other simple therefore occurs. Each summand must
saturate $|\dim X|\le\operatorname{FPdim}X$, by the triangle inequality.

In the original category, the objects for which either lift saturates
the bound generate a nontrivial fusion subcategory consisting entirely
of such objects. To see this, take two saturated simples, choose their
lifts of positive dimension, and decompose the product. Since the
categorical and Frobenius--Perron dimensions of that product agree,
the same triangle inequality forces every simple summand to be a
positive saturated lift. Positive lifts are also closed under duality.
If it is the whole category, all positive lift dimensions equal the
Frobenius--Perron dimensions. Thus $E=D$, as required. A proper subcategory of rank three or four would be integral by Proposition~\ref{prop:proper-integral}. Three rational simple dimensions would give at least two independent rational relations, whereas the dimension map has kernel of dimension one. A rank-two integral subcategory is pointed by the rank-two classification, and Proposition~\ref{prop:quartic-conductor} excluded a nontrivial invertible. The saturated subcategory must therefore have rank five, and the chosen
Galois conjugate is pseudo-unitary.
\end{proof}




% ===== quartic_finite.tex =====
\begin{lemma}\label{lem:quartic-moments}
Suppose no Galois conjugate is pseudo-unitary. Set $p=1/D$, $q=1/E$ and let $\alpha=E/D$ at the embedding maximizing $E$. Then $\alpha$ is a totally positive degree-four integral $d$-number, and its smallest conjugate is $\alpha_0<1$. If $r=g/f$, then
\[
 m_\alpha(X)\equiv(X-r)^4\pmod g,\qquad
 n_\sigma=\operatorname{Tr}(q\sigma\alpha)-(1-1/f)\in\mathbb Z_{\ge0}.
\]
Also $n_{\mathrm{id}}=0$ and $\sum_\sigma n_\sigma=(1-1/g)\operatorname{Tr}(\alpha)-4(1-1/f)$.
\end{lemma}
\begin{proof}
Central induction gives integrality, and the quotient of the two Galois-invariant principal ideals $(E)$ and $(D)$ is Galois invariant. Proposition~\ref{prop:stab} gives degree four. The minimality assertion is Lemma~\ref{lem:simultaneous-columns}.
Put $r=g/f$. In the signed order identified with a subring of $K\times\mathbb Q$, each pair $(\sigma\alpha,r)$ is integral. Indeed the signed class of the ordinary constituent of codegree $\sigma(D)$ has field value $E/\sigma(D)$ and rational value $r$ by its unit coefficient. Divide by the signed integral unit $(E/\sigma(E),1)$ from
Lemma~\ref{lem:simultaneous-columns} to obtain the claim. Elementary symmetric functions of these four pairs and the equality $T\cap(0\times\mathbb Q)=0\times g\mathbb Z$ yield the coefficientwise congruence
\[
 m_\alpha(X)\equiv(X-r)^4\pmod g.
\]
The unit coefficient of $(\sigma\alpha,r)$ is $\operatorname{Tr}(q\sigma\alpha)+r/g$. It is a positive integer and is one at the identity. Subtracting one gives $n_\sigma$. The sum formula follows by using $\operatorname{Tr}(q)=1-1/g$.
\end{proof}



% ===== biquartic.tex =====
\subsection{Descent in the biquadratic case}
\begin{theorem}\label{thm:biquartic-descent}
A rank-five complex fusion category whose Frobenius--Perron dimension
generates a biquadratic field has a pseudo-unitary Galois conjugate.
\end{theorem}
\begin{proof}
Suppose no Galois conjugate is pseudo-unitary.

\medskip\noindent\emph{Field and trace bounds.}
Write $G=\operatorname{Gal}(K/\mathbb Q)$
and let $c_1<c_2<c_3$ be the fundamental discriminants of its three
quadratic subfields, so $\Delta=\operatorname{disc}(K)=c_1c_2c_3$.
Retain the notation $p_i,q_i,\alpha_i=p_i/q_i$, with $i=0$ indexing
the minima in all three lists. Set $s=1-1/f$, $t=1-1/g$, and
$r=g/f$. By the preceding reductions $f\ge4$, $r$ is a positive
integer, and $\alpha$ is a totally positive primitive integral
$d$-number with $\alpha_0<1$. Write
\[
 b(c)=\sqrt{\gcd(c,4)/c},\qquad m_g=(g\Delta)^{-1/4}.
\]
For each nontrivial character with conductor $c$,
$|\widehat q(\chi)|\le b(c)$ and $|\widehat p(\chi)|\le b(c)$.
The norm estimates give $p_0,q_0\le m_4$, and $q_0\le m_g$.
Consequently a necessary pair of field tests is
\[
 \tfrac14\left(\tfrac34-\sum_{j=1}^3b(c_j)\right)\le m_4,
 \qquad
 \sqrt3\left(\tfrac34-4m_4\right)\le
       \sum_{j=1}^3\frac{2}{\sqrt{c_j}} .
\]
The first follows by Fourier inversion, and the second is the joint
inversion estimate from Corollary~\ref{cor:quartic-effective-conductor}.
For each $i\ne0$, the pair $\{0,i\}$ is the kernel of exactly one
quadratic character. Thus $q_0+q_i\ge(t-b(c))/2$. In particular
\[
 q_i\ge l_g:=\min_j\max\left\{
       \tfrac14(t-b(c_j)),\quad\tfrac12(t-b(c_j))-m_g\right\},
 \qquad A:=\operatorname{Tr}(\alpha)<1+s/l_g .
\]
Here $l_g>0$ since $t\ge3/4$ and $b(c_j)\le1/\sqrt2$.
This sharper trace bound distinguishes the three quadratic characters.
The bounds $c_1<40$, $c_2<729$ give $2097$ biquadratic fields before
the two tests. Exactly $52$ pass both.

\medskip\noindent\emph{Possible ratios.}
Write $d_j$ for the squarefree radicand of conductor $c_j$. Up to
conjugation,
\[
 4\alpha=A+x\sqrt{d_1}+y\sqrt{d_2}+z\sqrt{d_3},
 \qquad x,y\in\mathbb Z_{\ge0},\quad z\in\mathbb Z .
\]
Indeed each irrational summand is a character-weighted sum of
conjugates of $\alpha$, hence a trace-zero algebraic integer in the
corresponding quadratic subfield. Total positivity bounds its absolute
value strictly by $A$. Primitivity requires at least two nonzero
coefficients. First enumerate $4<A<1+1/l_4$ and all these coefficients.
If $P(X)=X^4-AX^3+BX^2-CX+H$ is their conjugate polynomial, all
four coefficients must be integral, $H>0$, and
$H^j\mid e_j^4$ for the elementary symmetric coefficients $e_j$.
Also require $\min_i\alpha_i<1$. For the remaining polynomials,
$P(X)\equiv(X-r)^4\pmod g$ gives $4r\le A$ because $4r\le g$.
Try the integers $1\le r\le A/4$, and those $g\ge4r$ divisible by
$r$ and dividing
\[
 \gcd(A-4r,B-6r^2,C-4r^3,H-r^4).
\]
This gcd is nonzero by primitivity. Reapply $A<1+s/l_g$ with the
actual $f,g$. The remaining numbers of tuples $(A,x,y,z,f,g)$ are
\[
\begin{array}{c|rrrrrr}
(c_1,c_2,c_3)&(5,8,40)&(8,12,24)&(8,17,136)&(8,21,168)&(8,28,56)&(8,60,120)\\\hline
\text{tuples}&6&12&1&2&4&5
\end{array}
\]
All thirty have $f=g=4$.

\medskip\noindent\emph{Integral traces.}
The integral unit-coordinate tests on
$(\sigma\alpha,r)$ imply
\[
 n_\sigma=\operatorname{Tr}(q\sigma\alpha)-s\in\mathbb Z_{\ge0},
 \qquad n_1=0,\qquad\sum_\sigma n_\sigma=At-4s.
\]
Enumerate these compositions for each tuple, relabeling the conjugates
so $\alpha_0$ is least. For each nontrivial character the convolution
identity
\[
 \widehat n(\chi)=\widehat q(\chi)\widehat\alpha(\chi)
\]
uses $G$ of exponent two. First require
$|\widehat n(\chi)|^2\le b(c)^2|\widehat\alpha(\chi)|^2$ for all
three characters. Every composition passing these rational inequalities
has all three $\widehat\alpha(\chi)\ne0$. Division and Fourier inversion
then determine every $q_i$ because $\widehat q(1)=t$. Put $p_i=q_i\alpha_i$.
Require $q_i>0$, that $p_0,q_0$ are least in their lists, and, for each
character, both
\[
 |\widehat p(\chi)|^2\le b(c)^2,\qquad
 2|\widehat p(\chi)|^2+|\widehat q(\chi)|^2\le4/c.
\]
No composition survives. The supplementary program
\texttt{verify\_biquartic.py} performs all enumerations in this proof.
Its only real inequalities use outward rational enclosures for the
square roots, computed by integer square roots.
\end{proof}



% ===== cyclic_quartic.tex =====
\subsection{The cyclic quartic calculation}\label{sec:cyclic-quartic}
\begin{proposition}\label{prop:cyclic-quartic-complete}
If $K=\mathbb Q(D)$ is cyclic quartic, the category has a pseudo-unitary
Galois conjugate.
\end{proposition}
\begin{proof}
Suppose no Galois conjugate is pseudo-unitary. Write $\sigma$ for a generator
of the Galois group,
$F$ for the conductor of its quadratic subfield and $J$ for the field
conductor. The conductor estimates give $F<40$, $J<729$. Let
$q=1/E$, $p=1/D$, $\alpha=p/q$, $s=1-1/f$, $t=1-1/g$ and $r=g/f$.
Use the signed embedding with $E$ largest; index its position by zero.
The preceding results give $0<\alpha_0<1$, $[\mathbb Q(\alpha):\mathbb Q]=4$,
$4\le f\mid g$, and the coefficientwise congruence
$m_\alpha(X)\equiv(X-r)^4\pmod g$.
The ratio $\alpha$ is a totally positive integral $d$-number by
Lemma~\ref{lem:quartic-moments}.

\medskip\noindent\emph{Conductor and trace bounds.}
An even primitive Dirichlet character of order four and its inverse
specify one real cyclic quartic field. At an odd prime $\ell$ its
nontrivial local character has conductor $\ell$ and sends a generator
to $i^a$ where $a\in\{1,2,3\}$ and $4\mid a(\ell-1)$.
It contributes $\ell$ to $F$ exactly for odd $a$. At two the primitive
possibilities have conductor $4$, $8$ or $16$. At $4$ the negative
unit has value $-1$. At $8$ the possibilities have value $-1$ on $5$
and either sign on $-1$. At $16$ the value at $5$ is $\pm i$, and $-1$
can have either sign. Only the last possibility contributes a factor
eight to $F$. Multiplying local characters and keeping even products
with nontrivial square gives 105 fields:
\[
\begin{array}{c|rrrrrr}
 F&5&8&13&17&29&37\\\hline
 \text{number of fields}&38&19&18&13&10&7 .
\end{array}
\]
Order-four even characters have conductor at least 15, as follows also
by checking the real cyclotomic degrees below 15.

Set $B(c)=\sqrt{\gcd(c,4)/c}$ and
$u=(gFJ^2)^{-1/4}$, so $q_0\le u$ by the norm bound.
The squared-twist Fourier estimates bound the two nontrivial Fourier
amplitudes of $q$ by $B(F)$ and $B(J)$.
The sum of entries at $0$ and $2$ is at least
$v_2=(t-B(F))/2$. The sum at $0$ and either neighbor is at least
$v_1=(t-\sqrt{2}B(J))/2$: the deficit relative to the other pair is
bounded by the sum of the absolute real and imaginary parts of the
order-four Fourier coefficient. Thus
$q_i\ge b_i:=\max(v_i/2,v_i-u)$ for $i=1,2$, and also $q_3\ge b_1$.
The conductor pairs in the displayed finite list give positive $b_i$.
For $A=\operatorname{Tr}(\alpha)$ this gives
\[
 A<1+s/\min(b_1,b_2).
\]
If $b_2\le b_1$, put $w=\min(s,(s+\sqrt{2/F})/2)$.
The coupled Fourier bound gives $p_2\le p_0+p_2\le w$.
Hence also $A<1+(s-w)/b_1+w/b_2$.
The verifier uses $s=1$ and $g=4$ as initial bounds before $f,g$ are
known, and then uses their actual values. All radicals in the bounds
are rounded outward using integer square roots.

\medskip\noindent\emph{Possible ratio polynomials.}
Write $z$ for the squarefree radicand of the quadratic subfield; thus
$z=F$ or $F/4$. Two opposite roots of $m_\alpha$ have sum and product
\[
 (A+B\sqrt z)/2,\qquad (C+D_1\sqrt z)/2.
\]
Here $B,C,D_1$ are integers and conjugation permits $B\ge0$.
Integrality implies the necessary parity conditions
$A\equiv zB\pmod2$ and $C\equiv zD_1\pmod2$.
For even $z$ these conditions are deliberately weaker than membership
in the quadratic integer ring; this only enlarges the search.
The other opposite pair has the conjugate sum and product.
Put $H=A^2+zB^2-8C$, $L=2AB-8D_1$. The two root differences have
squares $(H\pm L\sqrt z)/4$. The product of the two differences changes
sign under the generating four-cycle, so it is a rational multiple of
$\sqrt z$. Thus $H^2-zL^2=zW^2$ for a rational $W>0$, and $W$ is an
integer: a denominator in lowest terms would have its square dividing
the squarefree integer $z$. Consequently $H=zh$, $h\in\mathbb Z_{>0}$,
and $L^2+W^2=zh^2$. Positive pair sums and products give
$B^2z<A^2$, $0<H<A^2+zB^2$. These inequalities exhaust the finite
search on $B,h,L,W$ for each possible trace $A$. At $B=0$ the sign of
$L$ can also be fixed by conjugation. Recover $C,D_1$ from $H,L$.
The elementary symmetric functions of the four roots are
\[
 A,\qquad u=(A^2-zB^2+4C)/4,\qquad
 v=(AC-zBD_1)/2,\qquad w_0=(C^2-zD_1^2)/4 .
\]
They are integral. We require all roots positive and some root below one.
These sign tests use the pair sums and products and the coefficients of
$m_\alpha(X+1)$, all of whose roots are real.
Moreover $27w_0<A^3$ by arithmetic--geometric mean on the three roots
other than the minimum, and
$w_0^j\mid e_j(\alpha)^4$ for $1\le j\le3$ by the $d$-number criterion.

The index of $\mathbb Z[\alpha]$ gives
$J\mid\sqrt{\mathrm{disc}(m_\alpha)/F}$. This square root is
$|z(2D_1^2-AB D_1+B^2C)|W/8$ if $F=z$, and half this quantity otherwise:
multiply the two within-pair differences by the resultant of the two
quadratic factors. These conditions leave 221 parameter tuples before
the rational-axis test. That test requires $A\ge4r$, since
$A\equiv4r\pmod g$ and $g\ge4r$. Also $r w_0$ is a square by the two
order discriminants: their ratio is $g\operatorname{N}(E)/
(f\operatorname{N}(D))$. For each $r\le A/4$ the candidate $g$ divides
the gcd of the four coefficients of $m_\alpha(X)-(X-r)^4$ and $f=g/r$.
Reapplying the trace cutoff leaves 32 profiles $(\alpha,f,g)$,
six for $F=5$ and 26 for $F=8$.

\medskip\noindent\emph{Recovering the reciprocal dimension.}
Represent elements in the basis $1,\sqrt z,T,\sqrt zT$, with
$T^2=H+L\sqrt z$. Choose the generator sending
$\sqrt z\mapsto-\sqrt z$ and $T\mapsto W\sqrt z/T$.
In this representation $\alpha=(A+B\sqrt z+T)/4$.
The signed order contains $(\sigma^j\alpha,r)$ for all $j$ by the
preceding rational-axis argument. Its unit coefficient and total
positivity imply
\[
 n_j:=\operatorname{Tr}(q\sigma^j\alpha)-s
       \in\mathbb Z_{\ge0},\qquad
 n_0=0,\qquad \sum_{j=0}^3 n_j=At-4s .
\]
If $B\ne0$, then
$1,\alpha,\sigma\alpha,\sigma^2\alpha$ form a rational basis:
$\alpha+\sigma^2\alpha$ contributes the quadratic direction and the
two root differences are independent over $\mathbb Q$. Thus the traces
against these four elements, namely $t,s,s+n_1,s+n_2$, determine $q$.
If $B=0$, replace the last element by $\alpha^2$, which has a nonzero
$\sqrt z$ coefficient since $L\ne0$. Its trace against $q$ differs from
an integer by $-r^2/g$, using $(\alpha,r)^2$ in the signed order.
The integer is strictly positive and below $A$, because
$\operatorname{Tr}(q\alpha^2)<A\operatorname{Tr}(q\alpha)=As$
and $r^2/g=r/f\le A/(4f)$.
In this case $n_2=At/2-2s$. All linear systems have rational entries.

For each resulting $q$, require $\operatorname{N}(q)^{-1}$ and
$\operatorname{N}(q)^{-1}/w_0$ to be positive integers meeting the
signed and ordinary field-discriminant bounds, and require
$Fg/\operatorname{N}(q)$ to be a square. Ten candidates survive.
For $y=a_0+a_1\sqrt z+(a_2+a_3\sqrt z)T$, the squared Fourier amplitudes
of orders two and four are respectively
\[
 16z a_1^2,\qquad
 8\big[H(a_2^2+za_3^2)+2zLa_2a_3\big].
\]
The second identity adds the squares of the two differences of opposite
embeddings. Both expressions are rational. For each of the ten pairs
$p=\alpha q,q$, at least one of the coupled bounds
$2|\widehat p|^2+|\widehat q|^2\le4/F$ or
$2|\widehat p|^2+|\widehat q|^2\le4/J$ fails. In all tests the verifier
allows the smallest $J$ compatible with the index and trace bounds, which
only weakens the necessary conditions. This final contradiction and all
the counts above are reproduced by \texttt{verify\_cyclic\_quartic.py}.
\end{proof}



% ===== quintic.tex =====
\section{The quintic field and Galois descent}\label{sec:quintic}
\begin{theorem}\label{thm:quintic}
If $[\mathbb Q(D):\mathbb Q]=5$, then
$\mathbb Q(D)=\mathbb Q(\zeta_{11})^+$ and some Galois conjugate of
$\mathcal C$ is pseudo-unitary.
\end{theorem}
The degree condition is equivalent to the ordinary rational fusion algebra
being a field. Indeed a degree-five dimension map identifies that algebra
with its image. Conversely, for a quintic field algebra the ordinary
dimension character is injective. If $D$ were rational, all five codegrees
would coincide, giving $D=5$ by the reciprocal identity. Every simple
would then be invertible, and the dimension map could not be injective.
A nonrational element of a quintic field generates it.

\subsection{Determination of the field}
When $[\mathbb Q(D):\mathbb Q]=5$ both character algebras identify with the
same totally real cyclic quintic field $K$. For the signed algebra this
follows since $D\in L$ and $[L:\mathbb Q]\le5$. All basis elements are
self-dual, since their values at every character are real. In each order
let $H$ denote its distinguished dimension codegree, either $D$ or $E$.
The reciprocal identity and the two trace forms give
\[
 \operatorname{Tr}_{K/\mathbb Q}(1/H)=1,\qquad
 \operatorname{N}_{K/\mathbb Q}(H)=
 [\mathcal O_K:A]^2\operatorname{disc}(K).
\]
Indeed the ordinary trace of multiplication by $t$ is
$\tau(t\sum_i b_i^2)$, by computing the diagonal entries in the
orthonormal basis. Its Gram determinant is therefore the determinant of
multiplication by the Casimir, whose image in $K$ is $H$.

Let $c$ be the conductor of $K$. Its four nontrivial Dirichlet characters
are powers of one primitive order-five character and have the same
conductor, so $\operatorname{disc}(K)=c^4$ by the conductor--discriminant
formula \cite{Was}. The conductor is odd: a character of order five is
trivial on the unit group of any power of two. The only possible conductor
below $25$ is $11$, as a modulus below $25$ supports such a primitive
character only through a prime congruent to $1$ modulo five. All four nontrivial characters have odd conductor $c$. By Proposition~\ref{prop:parity-fourier}, their Fourier coefficients on either of the lists of five reciprocal conjugates have squared absolute value at most $1/c$. The trivial coefficient is one. Parseval gives
\begin{equation}\label{eq:qmoment}
 u,v\le\frac15+\frac{4}{5c},\qquad
 u=\operatorname{Tr}(D^{-2}),\quad v=\operatorname{Tr}(E^{-2}).
\end{equation}
Suppose $c\ge25$. The smaller of $u,v$ is at most $29/125$. Write
$y_1\le\cdots\le y_5$ for the corresponding reciprocal conjugates.
They sum to one. Since $y_1^2+(1-y_1)^2/4\le29/125$, we have
$y_1\ge1/25$. If $t=y_1+y_2\le1/5$, then
$\sum y_i^2\ge t^2/2+(1-t)^2/3\ge7/30>29/125$.
Hence $y_2>1/10$. The norm of the corresponding codegree is
less than $25\cdot10^4$, whereas it must be at least $25^4$.
Thus $c=11$. The unique quintic field of this conductor is the real
subfield of the eleventh cyclotomic field. In particular
\eqref{eq:qmoment} gives $u,v\le3/11$ separately.



% ===== quintic_trace.tex =====
\subsection{Integral traces force descent}
\label{sec:quintic-trace}
Put $p=1/D$, $q=1/E$, and denote their simultaneous conjugates by
$p_i,q_i$. The preceding subsection gives
\[
 \operatorname{Tr}(p)=\operatorname{Tr}(q)=1,\qquad
 \operatorname{Tr}(p^2),\operatorname{Tr}(q^2)\le3/11.
\]
Choose the signed embedding maximizing $E$ and take the identity as index zero, so $p_0,q_0$ are both minima. If two of the $q_i$ were at most $1/20$, then $\mathrm{Tr}(q^2)\ge 2(1/20)^2+(9/10)^2/3>3/11$ by Cauchy--Schwarz on the remaining entries. The bound decreases in each of the two small coordinates up to $1/20$. Thus $q_i>1/20$ for $i\ne0$. The ratio $\alpha=E/D$ is a totally positive algebraic integer by central induction, with identity value at most one and
\[
   \mathrm{Tr}(\alpha)=\sum_i p_i/q_i<1+20=21 .
\]
By Lemma~\ref{lem:simultaneous-columns}, $\alpha_0$ is its least
conjugate. A rational ratio is a positive integer at most one, hence
one, and gives descent. Assume therefore that $\alpha$ is irrational.

Let $G=\operatorname{Gal}(K/\mathbb Q)$. For each $\sigma\in G$,
Lemma~\ref{lem:simultaneous-columns} supplies the signed integral unit
of value $E/\sigma(E)$. The ordinary induction constituent of codegree
$\sigma(D)$ has signed value $E/\sigma(D)$. Dividing by that unit places
$\sigma(\alpha)$ in the signed order. Its unit coefficient is therefore the integer
\[
  h_\sigma=\mathrm{Tr}(q\sigma(\alpha))>0 .
\]
Furthermore $h_{\mathrm{id}}=\mathrm{Tr}(p)=1$, and $\sum_\sigma h_\sigma=\mathrm{Tr}(\alpha)$ since $\mathrm{Tr}(q)=1$.

We describe an exact enumeration of these traces. Put $b_j=\zeta_{11}^j+\zeta_{11}^{-j}$ for $1\le j\le5$. They give an integral basis: $b_1,\ldots,b_4$ are monic polynomials of successive degrees in $b_1$, $\sum b_j=-1$, and $\mathcal O_K=\mathbb Z[b_1]$. Write $[k]$ for the representative of $\pm k\bmod 11$ between one and five. The embedding matrix $M_{ij}=b_{[ij]}$ satisfies $M^2=11I-2J$, $MJ=-J$, since $\mathrm{Tr}(b_i b_j)=11\delta_{ij}-2$; $J$ denotes the all-ones matrix. Thus $M^{-1}=(M-2J)/11$. Writing $\alpha=-\sum c_j b_j$ and applying this inverse to its vector of positive conjugates gives
\[
   1\le c_j<4A/11,\qquad \sum_j c_j=A=\mathrm{Tr}(\alpha)\le20.
\]
In particular $c_j\in\{1,\ldots,7\}$.

For each nonconstant tuple, retain it unless disjoint rational isolating intervals contradict positivity, $\alpha_0\le1$, or $\alpha_0\le\sigma(\alpha)$. There are 34 tuples. Let $C$ be the matrix whose rows are the coordinate vectors of all five conjugates of $\alpha$, starting with $\alpha$. It is nonsingular: the cyclic group of order five acts regularly on these five coordinates, and over $\mathbb Q$ its augmentation representation is irreducible. The vector $c$ is nonconstant and has positive sum, so its translates span both rational summands. If $q=-\sum v_jb_j$, its trace is $\sum v_j=1$, and the trace vector $h$ determines its coordinates by
\[
   11Cv=h+2A\mathbf1 .
\]
Enumerate the positive integer vectors $h$ with first coordinate one and sum $A$. For each solution $v$, the signed square-trace test is $121\sum_jv_j^2\le25$, by the displayed trace pairing. Compute the coordinates of $p=\alpha q$ using $b_i b_j=b_{i+j}+b_{i-j}$, $b_0=2=-2\sum b_j$, and apply the same test.

The certificate \texttt{verify\_quintic\_trace.py} uses integer determinants and rational isolating intervals for the five roots of $x^5+x^4-4x^3-3x^2+3x+1$, in descending order. Of 15,396 trace vectors, 39 pass the signed test and four pass both. In every one of the four cases the coordinates of $p$ are a Galois permutation of the coordinates of $q$. Hence a conjugate of $E$ equals $D$, proving Theorem~\ref{thm:quintic}.



% ===== conclusion.tex =====
\subsection{Completion of the proof of Theorem A}
By Proposition~\ref{prop:rank-five-commutative}, the degree of
$D$ is at most five. If $D$ is rational, the category is weakly integral
and therefore pseudo-unitary \cite[Proposition 8.24]{ENO}.
For degrees two, three, four and five, Theorems~\ref{thm:quadratic-full},
\ref{thm:cubic-full}, \ref{thm:quartic-full} and \ref{thm:quintic}
respectively give a Galois conjugate whose intrinsic global dimension
is $D$. That conjugate has a positive spherical structure, which
transports to a spherical structure on the original category.
This proves Theorem~A, including Galois pseudo-unitarity.

The four finite calculations for descent use integer and rational
arithmetic, with outward rational bounds for radicals. Each search
range is derived in the preceding proofs. The source archive contains
the programs; \texttt{python3 verify.py --bounds-only} runs them together
with the calculations for the dimension and multiplicity bounds.




% ===== ring_classes.tex =====
\section{Dimension bounds for weakly integral categories}
\label{sec:classify}
From this section onward, Theorem~A allows us to replace $\mathcal C$
by a pseudo-unitary Galois conjugate and use its canonical positive
spherical structure. This preserves the based ring and all
Frobenius--Perron dimensions. Put $D=\operatorname{FPdim}(\mathcal C)$
and $K=\mathbb Q(D)$. A category is \emph{perfect} when the unit is its
only invertible simple.

\begin{lemma}\label{lem:mu-d}
For $d_i=\operatorname{FPdim}(X_i)$ one has $N_{ij}^k\le\min(d_i,d_j,d_k)$.
\end{lemma}
\begin{proof}
The dimension identity and Frobenius reciprocity give
\[
 N_{ij}^k\le\min\{d_i d_j/d_k,\ d_i d_k/d_j,\ d_j d_k/d_i\}.
\]
The quotient with the largest denominator is at most the least of the three dimensions.
\end{proof}


% ===== integral_bound.tex =====
\begin{proposition}\label{prop:integral-five}
An integral fusion category of rank five satisfies $D\le60$ and
$\mu\le3$. Both inequalities are sharp.
\end{proposition}
\begin{proof}
The integral classification \cite[Theorem~1.1]{ABDP} gives nine
rank-five based rings. Their dimensions and largest multiplicities are
\[
\begin{array}{c|rrrrrrr}
 D&5&8&14&20&21&24&60\\\hline
 \text{number of rings}&1&2&1&1&1&2&1\\
 \mu&1&1&1&3&2&1&2
\end{array}
\]
The two rings of dimension eight are the Tambara--Yamagami rings;
the two of dimension twenty-four are the representation ring of
$S_4$ and the ring in \cite[\S3.5, no.~11]{LPR}.
The remaining entries are the representation rings of $C_5$,
$D_{14}$, $\operatorname{AGL}(1,5)$, $C_7\rtimes C_3$ and $A_5$.
Their standard fusion rules give the displayed multiplicities.
The endpoints are attained by $\operatorname{Rep}(A_5)$ and
$\operatorname{Rep}(\operatorname{AGL}(1,5))$, respectively.
\end{proof}


\begin{proposition}\label{prop:weak-five}
If the category is weakly integral, then $\mu\le3$ and $D\le60$. Both bounds are attained by integral categories.
\end{proposition}
\begin{proof}
The integral case is Proposition~\ref{prop:integral-five}. Otherwise use the square-class grading of Gelaki and Nikshych
\cite[Theorem~3.10]{GN}. The neutral component is integral. The grading group is elementary abelian of order two or four, with all components of equal dimension. At order four the neutral rank is at most two. Neutral rank one forces integrality. At neutral rank two each of the three other components would consist of one simple of dimension $\sqrt2$, placing all three in the same square class. Thus the grading has order two. Its neutral category has rank at most four and integral dimension at most twelve \cite{ABDP}. Hence $D\le24$. Each simple belongs to a component of dimension at most twelve, so $d_i\le\sqrt{12}<4$. Apply Lemma~\ref{lem:mu-d}.
\end{proof}


\section{The action of the invertible simples}

% ===== ng.tex =====
\subsection{Four invertible simples}
\label{sec:four-invertibles}
For a finite group $G$, let $R(G,m)$ have basis $G\sqcup\{x\}$ and rules
\[
 gx=xg=x,\qquad x^2=\sum_{g\in G}g+mx.
\]
If exactly four of the five simples are invertible, they fix the
remaining simple under tensor product. That simple is self-dual, so
reciprocity gives $R(G,m)$ with $|G|=4$.

\begin{corollary}\label{cor:four-final}
The rank-five Grothendieck rings with four invertibles are
\[
 R(C_2^2,0),\quad R(C_2^2,4),\quad
 R(C_4,0),\quad R(C_4,3),\quad R(C_4,4).
\]
They satisfy $\mu\le4$ and $D\le16+8\sqrt2$.
\end{corollary}
\begin{proof}
Schopieray's classification gives $m\in\{0,4\}$ for
$G=C_2^2$ \cite[Theorem~1.1]{SchNG}. Meng and Yu give
$m\in\{0,3,4\}$ for $G=C_4$ \cite[Theorem~1.1]{MengYu}.
We apply these results to the positive spherical Galois conjugate
already supplied by Theorem~A; no new pivotality assumption is needed.
All five rings have the unitary realizations in
Table~\ref{tab:models-five}. Finally, the noninvertible dimension
satisfies $x^2=4+mx$, and $D=4+x^2$ is largest at
$m=4$, where $x=2+2\sqrt2$.
\end{proof}



% ===== invertibles.tex =====
\subsection{Three invertible simples}
\begin{theorem}\label{thm:three-inv}
A rank-five fusion category with exactly three invertible simples is integral.
\end{theorem}
\begin{proof}
Let $G=C_3$, let $S=\sum_{g\in G}g$, and denote the remaining simples by $X,Y$.
The action of $G$ fixes both of them. If $X^*=Y$, the product $X^2$ has no
invertible constituents. Hence $\operatorname{FPdim}(X)^2$ is a nonnegative
integer multiple of $\operatorname{FPdim}(X)=\operatorname{FPdim}(Y)$, proving
integrality. We treat the case of two self-dual objects.

Reciprocity and the commutativity of rank-five fusion rings give
\[
 X^2=S+aX+bY,\qquad XY=bX+cY,\qquad Y^2=S+cX+dY.
\]
Associativity gives
\begin{equation}\label{eq:three-conic}
 b^2+c^2-ac-bd=3.
\end{equation}
The three-dimensional rational summand $A=(S/3)(K_0(\mathcal C)\otimes\mathbb Q)$
has basis $S/3,X,Y$. Write $1$ for its unit. With
$x=\operatorname{FPdim}(X)$, $y=\operatorname{FPdim}(Y)$ and
$D=3+x^2+y^2$, the element
\[
 H=1+(X^2+Y^2)/3=3+\big((a+c)X+(b+d)Y\big)/3
\]
has Frobenius--Perron value $D/3$, an algebraic integer by divisibility of the
dimension by the dimension of the pointed subcategory
\cite[Proposition~8.15]{ENO}.
If $A$ were a field, its dimension character would be injective. The regular
trace of $H$ would then be an integer. The regular traces of $X,Y$ on $A$ are
$a+c,b+d$, respectively, so
\[
 \operatorname{Tr}_{A/\mathbb Q}(H)=9+\big((a+c)^2+(b+d)^2\big)/3.
\]
Thus $a+c\equiv b+d\equiv0\pmod3$. Equation~\eqref{eq:three-conic} now gives
$2(b^2+c^2)\equiv0\pmod3$, so $3$ divides all four coefficients. This
contradicts~\eqref{eq:three-conic} modulo $9$.

Suppose that at least one of $x,y$ is irrational. The algebra $A$ is
commutative and semisimple. The preceding argument forces it to be
$\mathbb Q\times L$ for a real quadratic field $L$, and the dimension character
has image $L$ after rationalization. Write $(u,v)\in\mathbb Z^2$ for the
values of the rational character on $X,Y$; rational character values are
algebraic integers. Orthogonality of the two distinct characters of the fusion ring gives
\begin{equation}\label{eq:three-linear}
 3+ux+vy=0.
\end{equation}
This spans the space of rational linear relations among $1,x,y$.

The two nontrivial characters of $G$, extended by zero on $X,Y$, have
formal codegree $3$. Let $B$ be the corresponding central simple for either
one. Then $B\subset I(1)$ has twist one and positive dimension $D/3$.
For $g\ne1$, put mass $[I(g):V]\operatorname{FPdim}(V)$ at the
twist of every simple $V$ of the center.
Its total mass is $D$. Its first two moments vanish by the indicator trace
formula \cite[Theorem~4.1]{NSExp}, because $\operatorname{Hom}(1,g^j)=0$ for $j=1,2$. Integrating
$|1+z+z^2|^2$ bounds its mass at one by $D/3$.
By central induction and adjunction,
\[
 [I(1),I(g)]=[FI(1):g]=\#\{T\text{ simple}: gT=T\}=2.
\]
All summands of $I(1)$ occur once and have twist one. If $B$ occurred
in $I(g)$, one occurrence would already account for the whole allowed mass,
while the displayed inner product forces another occurrence of positive
weight. This proves $[F(B):g]=0$ for both $g\ne1$. Since $[F(B):1]=1$,
there are nonnegative integers $p,q$ such that
\[
 D/3=1+px+qy.
\]
Using the displayed formula for $H$, this gives
$6+(a+c-3p)x+(b+d-3q)y=0$. Comparison with
\eqref{eq:three-linear} implies
$a+c\equiv 2u$ and $b+d\equiv 2v\pmod3$. Evaluating the two square rules at
$(u,v)$ gives $u^2+v^2=6+(a+c)u+(b+d)v$. Reduction modulo $3$ gives
$u^2+v^2\equiv0$, so $3\mid u,v$. The contradiction to
\eqref{eq:three-conic} modulo $9$ established above applies again.
This proves integrality and reduces the bounds to Proposition~\ref{prop:weak-five}.
\end{proof}



% ===== inv2fixed.tex =====
\subsection{An involution fixing every noninvertible}
\begin{theorem}\label{thm:fixed-inv-two-classification}
A rank-five fusion category with exactly two invertibles, whose
nonidentity invertible fixes all three noninvertible simples, is integral.
\end{theorem}
\begin{proof}
Write $g$ for the involution. Its negative eigenspace gives the
character $(1,-1,0,0,0)$ of codegree two. The corresponding constituent
$B\subset I(1)$ has dimension $D/2$. Write $A_1,A_2,A_3$ for the
other nonunit constituents. Thus
\[
 I(1)=1\oplus B\oplus A_1\oplus A_2\oplus A_3,
 \qquad \sum_i\dim A_i=D/2-1.
\]
All five constituents have twist one.

Because $g$ fixes a simple, its pointed subcategory has a rank-one
module category and hence trivial associator class. Therefore
$\nu_1(g)=0$ and $\nu_2(g)=1$. The total dimension and second twist
trace of $I(g)$ both equal $D$. Positivity forces its twists to be
$\pm1$, and its vanishing first trace puts dimension $D/2$ at each sign.
Adjunction also gives
\[
 [I(1),I(g)]=3,\qquad [I(g),I(g)]=5.
\]
The first counts the simples fixed by $g$, and the second follows
from commutativity of the Grothendieck ring.

The object $B$ cannot occur in $I(g)$: it would exhaust the twist-one
mass, although the first pairing requires further common constituents.
The multiplicities of $A_1,A_2,A_3$ in $I(g)$ therefore sum to three
and their squares sum to at most five. The only possibilities are
$(1,1,1)$ and a permutation of $(2,1,0)$. The latter exhausts the
norm with twist-one objects, contradicting the zero first trace.
Hence
\[
 I(g)=A_1\oplus A_2\oplus A_3\oplus U\oplus V,
\]
with $U,V$ distinct and new. Comparing the masses at the two signs
shows that $\theta_U=1$, $\dim U=1$, $\theta_V=-1$ and
$\dim V=D/2$.

The invertible $U$ restricts to $g$ and is its unique invertible
central lift, so $U^*=U$. It generates a central
$\operatorname{Rep}(C_2)$ whose forgetful functor is faithful.
Thus $\mathcal C\simeq\mathcal E^{C_2}$ for the corresponding
de-equivariantization. A fixed simple of $\mathcal E$ produces two
equivariant simples interchanged by $g$, whereas a free orbit produces
one simple fixed by $g$. The unit of $\mathcal E$ accounts for $1,g$;
the three remaining simples of $\mathcal C$ therefore come from three
free orbits. The induced involution on the rank-seven fusion ring
of $\mathcal E$ fixes only its unit. By
\cite[Theorem~2.3(1)]{SchFixed}, $\mathcal E$ is integral.
Consequently the three noninvertible dimensions of $\mathcal C$
are even integers.
\end{proof}


\section{An involution exchanging two simples}
\subsection{A representation subring of rank three}

% ===== s3swap.tex =====
\begin{proposition}\label{prop:s3-swap-pell}
Suppose $1,g,z$ span the representation ring of $S_3$, with
$\operatorname{FPdim}(z)=2$, and $g$ interchanges the two exterior
simples $X,Y$. Write $d=\operatorname{FPdim}(X)=\operatorname{FPdim}(Y)$.
If $d$ is irrational, there exists $h\in\mathbb Z_{\ge0}$ such that
\[
 XX^*=1+z+3h(X+Y),\qquad d^2=3+6hd.
\]
If $h\ge7$, there is an integer $t>0$ satisfying
\[
 3h^2+1=t^2\qquad\text{or}\qquad 3h^2+1=7t^2 .
\]
\end{proposition}
\begin{proof}
Frobenius reciprocity implies that the action of $z$ preserves the
exterior two-dimensional span. Since $zg=z$, it acts there by a
matrix whose rows and columns are constant; dimension gives $zX=zY=X+Y$.
Hence $XX^*=1+z+aX+bY$ and $d^2=3+kd$ for $k=a+b$.
Set $D=6+2d^2$ and $K=\mathbb Q(d)$.
The central simple in
$I(1)$ corresponding to the $d'$ row has dimension
$D/D'=1+(k/3)d$. Since its forgetful image has dimension in
$\mathbb Z+\mathbb Zd$, we have $3\mid k$.
For a dual pair, $XX^*$ is self-dual and $a=b$.
If $X$ is self-dual, put $m=a-b$. On the two characters with
$g=-1$, the values of $z,X,Y$ are $0,T,-T$, with
$T^2=1+mT$. Their codegrees are $f=2(2+mT)$.
For either codegree the positive dimension $D/f$ is the
dimension of the forgetful image of a central simple, so $f\in K$.
If $m\ne0$, $m^2+4$ is nonsquare (for $|m|=1$ it equals five,
and for $|m|\ge2$ it lies strictly between consecutive squares).
Then $K=\mathbb Q(\sqrt{m^2+4})=\mathbb Q(\sqrt{k^2+12})$.
This is impossible because $v_3(k^2+12)=1$ and $3\nmid m^2+4$.
Thus $m=0$ here as well. As $k$ is even and divisible by three,
$k=6h$. When $h=0$ the formula gives $d=\sqrt3$; this is the
weakly integral case and is to be retained if weak integrality is
allowed.

For $h\ge7$ let $f_K$ be the conductor of $K$. The dimension-block
bound for the nontrivial quadratic character gives
\[
 \frac1{\sqrt{f_K}}\ge |D^{-1}-(D')^{-1}|
   =\frac{k}{6\sqrt{k^2+12}} .
\]
Thus $f_K\le36(1+12/k^2)<37$. Write $3h^2+1=Qt^2$
with $Q$ squarefree. It is prime to three and $Q\equiv1\pmod3$.
It is odd: for odd $h$, $v_2(3h^2+1)=2$, and for even $h$
the integer $3h^2+1$ is odd. Since $K=\mathbb Q(\sqrt{3Q})$,
the conductor equals $3Q$ or $12Q$ according as $3Q\equiv1$
or $3\pmod4$. Hence $Q\le12$ and the only odd squarefree
positive integers in this interval congruent to one modulo three
are $1$ and $7$.
\end{proof}

\begin{lemma}\label{lem:actual-s3}The fusion subcategory on $1,g,z$ in Proposition~\ref{prop:s3-swap-pell} is tensor equivalent to $\operatorname{Rep}(S_3)$.\end{lemma}
\begin{proof}
The module spanned by $X,Y$ also determines the tensor category on
$1,g,z$. By the quasi-Hopf description in \cite[Theorem~7.1]{EGO}, this subcategory is a $C_2$-equivariantization of
$\mathrm{Vec}_{C_3}^{\omega}$ for inversion.
Let $A=\operatorname{Fun}(C_2)=1+g$ be the algebra in its canonical
central $\operatorname{Rep}(C_2)$. The $A$-modules in the exterior
module category form a module category over the de-equivariantization.
There is just one simple $A$-module: the two free modules on $X$
and $Y$ coincide and their endomorphism spaces are one-dimensional
by the free--forgetful adjunction. They generate all the modules by
separability. Thus $\mathrm{Vec}_{C_3}^{\omega}$ has a rank-one
module category, forcing the associator class to vanish. The
untwisted quasi-Hopf description identifies the original subcategory
with $\operatorname{Rep}(S_3)$.

\end{proof}



% ===== s3-moments.tex =====
\begin{proposition}\label{prop:s3-complete}
In Proposition~\ref{prop:s3-swap-pell}, one has $h\le1$. Consequently $\mu\le3$ and $D\le48+24\sqrt3$, with the latter equality precisely when $h=1$.
\end{proposition}
\begin{proof}
Let $V=X+Y$, let $R=1+g+2z$ denote the regular object of
$\mathcal B\simeq\operatorname{Rep}(S_3)$, and put
$\Delta=36h^2+12$, $D=6+2d^2=12(1+hd)$.
For a central simple $W$ with forgetful image $F(W)$, put
$A_W=[R,F(W)]$, $B_W=[V,F(W)]$. Then
$\dim W=A_W+B_Wd$ with $A_W,B_W$ nonnegative integers.
If $A_W>0$, the image contains a simple of $\mathcal B$, so that
the order of $\theta_W$ divides the Frobenius--Schur exponent of
$\mathcal B$, namely six
\cite[Proposition~2.2(2)]{SWY}. For $j\ge1$ the indicator trace formula for $R$ reads
\[
 12(1+hd)\nu_j(R)=\sum_W A_W^2\theta_W^j+
                         d\sum_W A_W B_W\theta_W^j.
\]
Both sums belong to $\mathbb Q(\zeta_6)$, and $\nu_j(R)$ counts the
solutions of $s^j=1$ in $S_3$. Since the real irrational $d$ does not
belong to this imaginary quadratic field, coefficient comparison yields
$\sum_W A_W^2\theta_W^j=12\nu_j(R)$ and
$\sum_W A_W B_W\theta_W^j=12h\nu_j(R)$. Applying the indicator formula to $V$ and
using $D/d=2\sqrt\Delta$, $12(h/d+h^2)=2h\sqrt\Delta$ yields
\[
 M_j:=\sum_W (B_W-hA_W)^2\theta_W^j
     =2\sqrt\Delta\bigl(\nu_j(V)-h\nu_j(R)\bigr).
\]
To compute $M_0$, write $L_T$ for left multiplication by $T$.
For a virtual object $T$,
$[I(T),I(T)]=\operatorname{Tr}(L_T L_{T^*})$
by adjunction and $FI(T)=\sum_i b_i T b_i^*$.
Apply this to the self-dual $T=V-hR$. The two dimension characters
take values $\pm\sqrt\Delta$ there. The remaining three characters
vanish on both $R$ and $V$: the two with $g=-1$ have $Y=-X,z=0$;
the one with $g=1,z=-1$ has $X=Y=0$. Hence
$M_0=2\Delta=72h^2+24$.

Write $3h^2+1=Qt^2$ with $Q$ squarefree and put $P=3Q$, so
$\sqrt\Delta=2t\sqrt P$.
The first two indicators of $V$ have values zero and
$\epsilon\in[-2,2]$, while those of the regular $S_3$ representation
are one and four, counting solutions of $s^j=1$. Thus
\begin{equation}\label{eq:s3-moments}
 M_1=-4ht\sqrt P,\qquad M_2=4(\epsilon-4h)t\sqrt P.
\end{equation}
The coefficients defining these two moments are the same nonnegative
weights of total mass $M_0$.

Suppose $Q=7$. In the normalized cyclotomic trace down to
$\mathbb Q(\sqrt{21})$, the absolute value of the $\sqrt{21}$-coefficient
of any root of unity is at most $1/12$. Indeed, project first to
$\mathbb Q(\sqrt{-3},\sqrt{-7})$. A primitive third or seventh root
has quadratic-period average $(-1\pm\sqrt{-p})/(p-1)$ for
$p=3,7$, respectively. Higher primary orders average to zero and the
other primes give a rational factor of absolute value at most one.
The quadratic component surviving projection to $\mathbb Q(\sqrt{21})$
has the asserted bound. Applying this to the second moment gives
$48(4h-2)t\le72h^2+24$. For $h\ge7$ this is impossible since
$t>h\sqrt{3/7}$: even for $h\ge2$ the left side exceeds
$144\sqrt{3/7}\,h^2>78h^2$.

Now suppose $Q=1$, and write the normalized trace of $\theta^j$
to $\mathbb Q(\sqrt3)$ as $a_j+b_j\sqrt3$.
For every root of unity one has
\[
 -3b_1-2b_2-a_2\le1.
\]
To see this, decompose the order into its $2$-part, $3$-part and
prime-to-six part. The last component contributes the same factor
$\rho=\mu(l)/\varphi(l)$ to both averages, where $l$ is its order,
so $-1/4\le\rho\le1$. A nonzero $b_1$ requires exact primary
orders four and three. Here the product of the primary components
is a primitive twelfth root, with normalized trace $\pm\sqrt3/2$;
its square has normalized trace $1/2$. Thus the functional equals
$\rho$ or $-2\rho$. A nonzero $b_2$ requires exact primary orders
eight and three. The first trace vanishes by averaging over the
kernel of restriction to twelfth roots, and the square has trace
$\pm\rho\sqrt3/2$; here the functional equals $\pm\rho$.
In all other cases both irrational coefficients vanish and
$|a_2|\le1$. This proves the inequality. Summing it with the weights
of \eqref{eq:s3-moments}, we may save $4h^2$ against the total mass. Indeed the unit and the induction simple for the conjugate dimension character both have twist one, where the functional is $-1$. Their dimensions $1$ and $1+2hd$ show that each has weight $h^2$. This gives
\[
 (44h-16)t\le68h^2+24 .
\]
For $h\ge5$ this contradicts $t>h\sqrt3$, since
$(44-16/h)\sqrt3>68+24/h^2$ already at five and the two sides
are respectively increasing and decreasing. In particular both fields
are excluded throughout the range $h\ge7$. For $h=2,5,6$ the squarefree
parts of $3(3h^2+1)$ are $39,57,327$, with conductors $156,57,1308$.
Each exceeds $36(1+12/k^2)$. For $h=3$ the preceding $Q=7$ bound
applies directly. Finally $h=4$ has $Q=1,t=7$, and the last inequality
reads $1120\le1112$. This proves the assertion.
\end{proof}



% ===== exchangebound.tex =====
\subsection{A norm bound for the exchanged pair}
\begin{proposition}\label{prop:exchange-large}
Suppose $g$ exchanges $X,Y$ and fixes $Z$, and $D$ is irrational. The positive character on the $g=1$ summand is quadratic. Write $(1,1,u,u,v)$ for the remaining rational row and $f=2+2u^2+v^2$ for its codegree. If $F$ denotes the conductor of the dimension field $K$ and $L=\operatorname N(D/f)$, then $L$ is a positive integer and
\begin{equation}\label{eq:exchange-norm-gap}
 (f-2)^2-4f^2/F\le 16/L .
\end{equation}
In particular $f\le22$. When $uv\ne0$ the field is $K=\mathbb Q(\sqrt{fL})$.
\end{proposition}
\begin{proof}
Write $x=d_X=d_Y$ and $z=d_Z$.
On the $g=-1$ block, the characters have exterior values $\lambda,-\lambda,0$. If $X$ is self-dual, $\lambda^2=1+t\lambda$ for $t=N_{XX}^X-N_{XX}^Y$, by evaluation of $X^2$. Their codegrees are $2(1+\lambda^2)$. When $X^*=Y$, the nonunit exterior coefficients of $XX^*$ on $X,Y$ agree by duality, so evaluation gives $|\lambda|^2=1$ and both codegrees are four. Put $t=0$ in this case. In both cases the reciprocals of the two codegrees sum to $1/2$.

The $g=1$ factor has three rows with exterior coordinates $x,x,z$ and their character analogues. If the positive row has degree three, $1,x,z$ are a rational basis of its field. The negative-block codegrees must lie in that field by central induction, hence are rational. In the self-dual case $t\ne0$ would imply rational $\lambda=(\lambda^2-1)/t$ and thus integral roots of product $-1$, forcing $t=0$. Thus in both cases the negative-block codegrees are four. The dimension of a corresponding central simple is $D/4$, an integral combination of $1,x,z$. But the coefficient of $1$ in $D$ in this basis is six: the two invertibles contribute two, each of $XX^*,YY^*$ has one pointed constituent, and $ZZ^*$ has two. This contradicts integrality. Degree one would give rational $D$. The positive row and its conjugate therefore leave one rational row, with integral character values $u,u,v$.

Write $q=1/D+1/D'=1/2-1/f$. The central simple induced from the rational row has dimension $D/f$, a totally positive algebraic integer. Thus $L=DD'/f^2$ is a positive integer. The dimension-block inequality gives
\[
 1/F\ge(1/D-1/D')^2=q^2-4/(f^2L),
\]
which yields \eqref{eq:exchange-norm-gap}. Using $F\ge5$ and $L\ge1$ gives $f^2-20f-60\le0$, so $f\le 10+4\sqrt{10}<23$.

In the $g=1$ quotient write the products as
$X^2=1+sX+cZ$, $XZ=2cX+dZ$, $Z^2=2+2dX+eZ$. Then
\begin{equation}\label{eq:exchange-rational}
 u^2=1+su+cv,\qquad uv=2cu+dv,\qquad v^2=2+2du+ev .
\end{equation}
If $v=0$, orthogonality gives $ux=-1$, impossible for $x>1$. If $u=0$, then $vz=-2$, and the rational algebraic integer $z>1$ must be two. Thus $v=-1$, $c=e=1$, $d=0$, giving the $S_3$ subring. Otherwise $u,v$ are coprime by the first equation and have opposite signs by the second equation and orthogonality.

Suppose $uv\ne0$ and set $H=2u^2+v^2$. From $vz=-2-2ux$ and $2+2xx'+zz'=0$ one obtains
\[
 Hxx'+2u(x+x')+v^2+2=0,\qquad
 \operatorname{N}(2Hx+4u)=-4v^2 f .
\]
Subtracting the orthogonality relation from $D=2+2x^2+z^2$ gives
$D=(x-x')(2x-2uz/v)=(x-x')(2Hx+4u)/v^2$. Taking norms yields
\begin{equation}\label{eq:exchange-discriminant}
 (x-x')^2=fv^2L/4 .
\end{equation}
As $z$ is affine in $x$ and $D$ is irrational, $x$ generates $K$. The last assertion follows.
\end{proof}



% ===== eleven.tex =====

% ===== morita_exchange.tex =====
\subsection{The Morita dual of $1+g$}
The next construction will be used twice. It explains why the parity
of the coefficient of $Z$ in $Z^2$ matters.

\begin{lemma}\label{lem:exchange-morita}
Assume the irrational exchanged-pair situation of
Proposition~\ref{prop:exchange-large}, and suppose $e=N_{ZZ}^Z$ is odd.
Put $x=d_X=d_Y$, $z=d_Z$ and $q=1/D+1/D'$.
The algebra $A=1+g$ has four simple right modules, lying over
$1+g,X+Y,Z,Z$. Its Morita dual
$\mathcal B={}_A\mathcal C_A$ has eight simples, of dimensions
\[
 1,1,x,x,z/2,z/2,z/2,z/2.
\]
Let $h$ be the nontrivial invertible bimodule lying over $A$.
Then $h^2=1$, left multiplication by $h$ fixes no simple, and
exactly four simple classes commute with $h$.

The complexified Grothendieck algebra of $\mathcal B$ is
$\mathbb C^4\oplus M_2(\mathbb C)$. The two linear characters with
$h=1$ have codegrees $D,D'$. Write $H_1,H_2$ for the codegrees of
the two with $h=-1$. They belong to $K$ and satisfy
\[
 \frac1{H_1}+\frac1{H_2}=q,\qquad
 f_{2}=\frac2{1-2q},
\]
where $f_2$ is the codegree of the two-dimensional representation.
After identifying the centers by Morita equivalence
\cite{SchCenter}, induction satisfies
\begin{equation}\label{eq:exchange-morita-pairings}
\begin{aligned}
 [I_C(1),I_B(1)]&=4,&[I_C(1),I_B(h)]&=2,\\
 [I_B(1),I_B(h)]&=0,&[I_B(h),I_B(h)]&=4.
\end{aligned}
\end{equation}
\end{lemma}
\begin{proof}
The fixed simple $Z$ gives a rank-one module category over the pointed
subcategory generated by $g$. Its associator is therefore trivial in
cohomology, so $A=1+g$ is a separable algebra.
Its simple right modules are $A$, the free module $P$ on $X$,
and two modules $Z^+,Z^-$ whose actions differ by sign.
Their underlying objects are $1+g,X+Y,Z,Z$, respectively.

Left multiplication by $g$ fixes $A$ and $P$ and permutes $Z^+,Z^-$.
If it exchanged the latter pair, $Zg=Z$ would imply
$ZZ^+=ZZ^-$ in the module Grothendieck group. Their sum is the free
module on $Z^2$; the multiplicity of $Z^+$ in it is $e$, which is
odd. This is impossible. Thus all four right modules are fixed.
For each of them the two choices of a compatible left $g$-action
give two simple bimodules: the coherence obstruction is a nonzero
scalar, which has two square roots over $\mathbb C$. Dividing the
underlying dimensions by $\operatorname{FPdim}(A)=2$ gives the stated
list.

Multiplying by $h$ on the left or right twists the corresponding
$A$-action by sign. Twisting the left action fixes no bimodule,
because each underlying right module is simple. On the four classes
lying over $Z$, the left and right sign changes generate a regular
four-element orbit; hence none of these classes commutes with $h$.
On $A$ and $P$ the right sign change is undone by an automorphism
acting with opposite signs on their two simple constituents. It
also reverses the left $g$-action, so here the left and right twists
coincide. Exactly these four classes commute with $h$.

In particular the eight-dimensional Grothendieck algebra is
noncommutative. Semisimplicity and the existence of its dimension
character force the decomposition $\mathbb C^4\oplus M_2(\mathbb C)$.
On the matrix block, $h$ is a nonscalar involution and has trace zero.
Its regular trace is also zero, because it fixes no basis element.
Thus two linear characters take $h$ to each sign. The positive
character and its quadratic conjugate take $h=1$ and exhaust that
pair. Expanding the coefficient-of-unit trace at $h$ gives
$1/H_1+1/H_2=q$; expanding it at $1$ then gives the formula for $f_2$.
The dual of a pseudo-unitary category is pseudo-unitary, and its
simple dimensions lie in $K$. Central induction consequently places
all its formal codegrees in $K$.

For completeness the mixed induction formula is
\[
 [I_C(T),I_B(V)]
 =\sum_{M\in\operatorname{Irr}(\mathcal C_A)}[TM,MV].
\]
It follows from the induction--restriction formula in the linking
multifusion category with diagonal parts $\mathcal C,\mathcal B$
and off-diagonal parts $\mathcal C_A$ and its opposite.
The four right modules give the first pairing in
\eqref{eq:exchange-morita-pairings}. Right multiplication by $h$
fixes $A,P$ and exchanges $Z^+,Z^-$, giving the second.
The last two count, respectively, simples fixed by left multiplication
by $h$ and simples commuting with $h$. The identified centers have
their canonical positive spherical structures, so the identification
also preserves twists.
\end{proof}


\begin{proposition}\label{prop:exchange-eleven}
In the irrational exchanged-pair case, $|u|=2$, $|v|=1$ and
$K=\mathbb Q(\sqrt5)$ cannot hold together.
\end{proposition}
\begin{proof}
The relation $v^2=2+2du+ev$ makes $e$ odd. Apply
Lemma~\ref{lem:exchange-morita}; here $f=11$ and $q=9/22$.
First $H_1,H_2$ form a nontrivial quadratic orbit. Equal codegrees
would give the noninteger $44/9$. Distinct rational codegrees would
distinguish their two characters, making both characters integer-valued
on the basis. The left sign action pairs the two simples over $P$
with opposite values. The left and right sign actions give equal
absolute values on the four simples over $Z$. Thus each codegree
would be $2+2a^2+4b^2$. Since $2<H_i$ and
$\min H_i\le2/q<5$, one would be four and its partner $44/7$,
again impossible.

Put $n=H_1H_2$. The norm and trace of this pair are integers,
and its trace is $qn$. Hence $22\mid n$. Positivity gives
$n\ge4/q^2>22$, so $n\ge44$.

Apply Lemma~\ref{lem:hecke-projection} to the center of $\mathcal B$.
The difference of the vacuum basis vector and its quadratic Galois
translate, normalized by $\sqrt2$, has conductor five and is supported
at twist one. Its $S$-image at the four linear-character induction
constituents gives
\[
 (q-2/D)^2+q^2-4/n\le1/5.
\]
Here the two modular row entries at a constituent of codegree $k$
are $1/k$ and its conjugate. The two quadratic pairs therefore
contribute $(1/D-1/D')^2$ and $(1/H_1-1/H_2)^2$.

Finally $DD'=121L$ with $L\ge1$, so the displayed left side is
\[
 2q^2-\frac4{121L}-\frac4n
 \ge2(9/22)^2-\frac4{121}-\frac4{44}>\frac15,
\]
a contradiction.
\end{proof}



% ===== exchange_pointwise.tex =====
\begin{corollary}\label{cor:exchange-pointwise}
In the irrational exchanged-pair case, outside the $S_3$ subring, one of the following holds:
\[
\begin{array}{c|c|rrrrrr}
f&F&u&v&c&s&d&e\\\hline
5&5&u&-u&c&c&2c+u&4c+3u\\
8&8&-1&2&2&4&1&2\\
20&5&-1&4&4&16&1&4
\end{array}
\]
In the first row $u=\pm1$ and $c\ge1$.
\end{corollary}
\begin{proof}
By Proposition~\ref{prop:exchange-large}, $2+2u^2+v^2\le22$ with $u,v$ nonzero and coprime. The possible $f$ are $5,8,11,13,19,20,21$. Inequality~\eqref{eq:exchange-norm-gap} forces $F=5$ for $f\ge11$, and $F\in\{5,8,12\}$ at $f=8$. The norm $-4v^2f$ excludes $(f,F)=(8,5),(13,5),(5,8)$ because $2,13,5$, respectively, are inert in the field in question and occur to odd valuation. Proposition~\ref{prop:exchange-eleven} excludes $f=11$. For $f=19,21$, the identity $K=\mathbb Q(\sqrt{fL})=\mathbb Q(\sqrt5)$ requires $L$ to be $5f$ times a positive integer square. In contrast, \eqref{eq:exchange-norm-gap} gives $L\le80,1$, respectively.

For $f=8$ and $F=12$, the same inequality gives $L=1$, inconsistent with $K=\mathbb Q(\sqrt{8L})$. For $(f,F)=(8,8),(20,5)$, $L$ is a square at most four. If $L=1$, the minimal polynomial of $D/f$ has trace $(f-2)/2$ and norm one, giving discriminant five or $77$, respectively, in the wrong field. Thus $L=4$. In both cases put $w=|v|$, so $u=\pm1$, $v=-uw$, $s=cw$. Using $vz=-2-2ux$ in $x^2=1+sx+cz$ gives
\[
 (x-x')^2=(w+2/w)^2c^2+8uc/w+4.
\]
Its value from \eqref{eq:exchange-discriminant} is $32$ when $w=2$ and $320$ when $w=4$. These quadratics respectively force $(c,u)=(2,-1),(4,-1)$, giving the last two rows.

Finally let $f=5$. Then $w=1$ and $(x-x')^2=9c^2+8uc+4=5L/4$. In particular it is an integer divisible by five. If $F\ge12$, inequality~\eqref{eq:exchange-norm-gap} gives $L\le24$, whence $9c^2+8uc+4\le30$ and $c\le2$. Among $c=0,1,2$ only $c=1,u=-1$ gives a multiple of five within this range, and its discriminant is five. Thus $F=5$. Relations~\eqref{eq:exchange-rational} give the remaining row; $c=0$ would give $x^2=1$.
\end{proof}



% ===== swap_five_prime.tex =====
\begin{proposition}
In the exchanged-pair case with $v=-u\in\{1,-1\}$ and $f=5$, suppose $\mathbb Q(x,z)=\mathbb Q(\sqrt5)$ and $c\ge2$. Then $uc\equiv4\pmod9$.
\end{proposition}
\begin{proof}
The rational character relations give
$s=c$, $d=2c+u$ and $e=4c+3u$, while $z=2(x+u)$ and
$x^2=1+3cx+2uc$. Consequently
\[
 D=2+2x^2+z^2=12+12uc+(18c+8u)x,\qquad
 D/D'=3D/10-1,
\]
where the second identity uses $1/D+1/D'=1/2-1/5$.
Let $N$ be the Frobenius--Schur exponent, equal to the order of the
center's twist \cite[Theorem~5.5]{NSExp}. If $3\nmid N$, the automorphism
$\zeta_N\mapsto\zeta_N^{3^{-1}\bmod N}$ acts nontrivially on
$\mathbb Q(\sqrt5)$, the conductor-five dimension field of the center:
the center contains an induced summand of dimension $D/5$. By the
higher Gauss sum formula of \cite[Theorem~4.4, equation~(4.17), in the arXiv version]{NSW} and the central indicator
identity, one obtains
\[
 \left|1+x(\nu_3(X)+\nu_3(Y))+z\nu_3(Z)\right|=D/D'.
\]
Here $\nu_3(g)=0$ since $g^3\ne1$. Every third indicator is rational:
the rotation formula places it in $\mathbb Q(\zeta_3)$ and the central
formula places it in $\mathbb Q(\zeta_N)$, using that the central
dimensions lie in the latter field. These fields are disjoint.
Set $A=\nu_3(X)+\nu_3(Y)$ and put $B=\nu_3(Z)$.
Comparison of the rational and irrational coordinates in
\[
 1+2uB+(A+2B)x
 =\pm\bigl((13+18uc)/5+(27c+12u)x/5\bigr)
\]
gives $A=u\pm(9c-u)/5$, so $|A|\ge(9c-6)/5>c$.
On the other hand, the third indicators are traces of finite-order
operators on the respective invariant spaces, whose dimensions sum to
$[X^3:1]+[Y^3:1]=s=c$. Indeed $Y^2=X^2$, and $X^*,Y^*$ are $X,Y$
in some order. This contradiction proves $3\mid N$.

Categorical Cauchy now gives $3\mid N_{K/\mathbb Q}(D)$, with
$K=\mathbb Q(\sqrt5)$. Taking the trace in the expression for $D$ above
and using $1/D+1/D'=3/10$ gives
\[
 N_{K/\mathbb Q}(D)=20(9c^2+8uc+4).
\]
The expression in parentheses is the discriminant of the quadratic
polynomial of $x$, so equals $5b^2$ for an integer $b$. Therefore $3\mid b$
and $8uc+4\equiv0\pmod9$, as claimed.
\end{proof}



% ===== five_morita.tex =====
\subsection{The remaining quadratic family}
\begin{theorem}\label{thm:exchange-five-effective}
In the irrational exchanged-pair case with $f=5$ the only possible parameters $(c,u)$ are $(1,-1)$ and $(4,1)$.
\end{theorem}
\begin{proof}
Corollary~\ref{cor:exchange-pointwise} gives $K=\mathbb Q(\sqrt5)$, $v=-u$, $s=c$, $d=2c+u$, $e=4c+3u$ and $z=2(x+u)$. Here $gX=Y$, $gZ=Z$, $x=d_X=d_Y$, $z=d_Z$, and $q=1/D+1/D'=3/10$.

Since $e=4c+3u$ is odd, Lemma~\ref{lem:exchange-morita}
applies. Retain its notation $\mathcal B,h,H_1,H_2$.
The four linear-character induction constituents have codegrees
$D,D',H_1,H_2$, and the two-dimensional representation has codegree
$2/(1-2q)=5$.
The numbers $H_1,H_2$ belong to $K$ by central
induction, and their Galois-invariant product $n$ is an integer.
The trace $qn$ is integral, so $10\mid n$. Since $n\ge4/q^2$,
the only possibilities below $70$ are $50,60$. For $n=60$ the
two roots are $9\pm\sqrt{21}$, outside $K$. For $n=50$ they
are $5,10$. More generally neither $H_i$ can be five: its partner
would be ten, so the corresponding $h=-1$ linear characters are
rational because their codegrees distinguish them. Their integral
values pair in sign-opposite pairs under the left action of $h$, so
both codegrees must be even. This proves $n\ge70$.

The vacuum and its quadratic Galois translate occur in both
inductions of the unit: the normalized row at the latter is the
Galois conjugate of the vacuum row, in either description of the
same center. Besides these two, $I_C(1)$ has one label of codegree five
and two labels from the negative $g$ block, whose reciprocal codegrees
sum to $1/2$. Neither negative-block codegree equals five: they are
four when $X^*=Y$ and otherwise are $2(1+\lambda^2)$ for
$\lambda^2-t\lambda-1=0$, $t\in\mathbb Z$.
The codegree-five labels on the two sides must coincide.
Otherwise their dimensions match none of the available labels on
the other side; of the remaining two labels on each side at most
one can match, since the sums of their reciprocal codegrees are
$1/2$ and $q$ respectively. This would give
$[I_C(1),I_B(1)]\le3$. The common codegree-five label occurs twice
in $I_B(1)$, accounting for the entire remaining overlap. The two
$H_i$ labels are therefore distinct from all labels of $I_C(1)$.

Let $P,N_-$ denote the masses at twists $1,-1$ in the center, normalized
by $D^2$ and weighted by squared simple dimensions. The union of both
unit inductions lies at twist one. Cauchy's inequality on the negative
$g$ block and the bound on $n$ give
\[
 P\ge D^{-2}+(q-D^{-1})^2+1/25+1/8+q^2-2/70.
\]
By \eqref{eq:exchange-morita-pairings}, the squared multiplicities
of $I_B(h)$ sum to four. It cannot be a
single doubled label, since its first twist trace is zero. Thus it has
four distinct labels. The two mixed pairings force it to contain
both negative $g$-block labels of $I_C(1)$, disjoint from $I_B(1)$.
Their dimensions total $D/2$ and their twists are one; the other two
dimensions total $D/2$. Vanishing of the first twist trace and
positivity force the other two twists to be $-1$. Hence $N_-\ge1/8$.

For the same normalized measure, $m_1=1/D$ and
$|m_2|\le(2+\sqrt{3(D-2)})/D$ by the central indicator formula.
Integrating $|3+2w+3w^2|^2$ gives
\[
 64P+16N_-\le22+24/D+18(2+\sqrt{3(D-2)})/D .
\]
Using $D^{-2}+(q-D^{-1})^2=q^2-2/\operatorname N(D)$ in the lower bound for $P$ gives
\[
 \frac{44}{175}\le
 \frac{128}{\operatorname N(D)}+\frac{60}{D}
       +\frac{18\sqrt{3(D-2)}}{D}.
\]
The discriminant $9c^2+8uc+4$ of $x$ is divisible by five, so $uc\equiv4\pmod5$. For $c\ge2$ the preceding third-indicator congruence yields $uc\equiv4\pmod9$. Hence $uc\equiv4\pmod{45}$, and either $(c,u)=(4,1)$ or $c\ge41$. In the latter case $x=(3c+\sqrt{9c^2+8uc+4})/2\ge3c-1\ge122$. Thus $D\ge2(122)^2$ and $\operatorname N(D)=20(9c^2+8uc+4)\ge20(121)^2>2(122)^2$ by \eqref{eq:exchange-discriminant}. The right-hand side is less than $94/122^2+18\sqrt{3/2}/122<1/4<44/175$, a contradiction. At $c=1$, divisibility of the discriminant by five forces $u=-1$.
\end{proof}



% ===== exchange_sharp.tex =====
\subsection{Sharp bounds for the exchanged pair}
\begin{proposition}\label{prop:sharp-exchange}
Suppose $g$ exchanges $X,Y$ and fixes $Z$, and $D$ is irrational. Then $\mu\le3$ and $D\le48+24\sqrt3$. Outside the $S_3$-subring case, $\mu=1$ and $D=15+5\sqrt5$. At the dimension maximum the $S_3$ parameter is $h=1$.
\end{proposition}
\begin{proof}
Set $x=d_X=d_Y$, $z=d_Z$. When $u=0$, Proposition~\ref{prop:s3-complete} gives the stated bounds. Otherwise Corollary~\ref{cor:exchange-pointwise} and Theorem~\ref{thm:exchange-five-effective} leave four sets of parameters $(f,F,c,u)$. The first, $(5,5,1,-1)$, has $s=c=d=e=1$ and $|t|=1$, so is multiplicity free. Its dimensions $x=(3+\sqrt5)/2$, $z=2(x-1)$ give $D=10x=15+5\sqrt5$.

\par\smallskip\noindent\emph{Saturation at conductor eight.}\enspace
At $(f,F,c,u)=(8,8,2,-1)$ one has $x=3+2\sqrt2$, $z=x-1$, $p=D^{-1}=(3-2\sqrt2)/16$. In particular $(p-p')^2=1/8$. Let $A$ be the induction constituent for the conjugate dimension character and $v_0=(e_1-e_A)/\sqrt2$ in the center's modular representation. The entries of $Sv_0$ at $1,A$ already contribute $1/8$ to the squared norm of its twist-one projection. This attains the conductor-eight bound. More precisely, in the proof by averaging $a(b)=\langle v_0,L(b)v_0\rangle$ over lower unipotents modulo the twist exponent, all strata with $8\nmid b$ cancel, and $|a(b)|\le1$. Equality forces $a(b)=1$ for all $8\mid b$. Thus $L(8)v_0=v_0$, and $Sv_0$ is supported on twists of order dividing eight.

For a central simple $W$ of twist order sixteen it follows that $\dim(W)/D$ is fixed by quadratic conjugation: the two row entries at $W$ agree. Neither $1$ nor $g$ occurs in $F(W)$. Indeed $I(1)$ has twist one, and $I(g)$ has only twists $\pm1$ by the second indicator trace; the associator on $\langle g\rangle$ is trivial since $gZ=Z$. If $F(W)$ contains in total $b_W$ copies of $X,Y$ and $c_W$ copies of $Z$, then
\[
 \dim(W)/D=(b_W+c_W)/16-c_W/(16x).
\]
It is rational only if $c_W=0$. For the virtual object $T_0=X+Y-2Z$ the character values are $2,2,-6,0,0$. Thus $[I(T_0),I(T_0)]=44$ by induction--restriction, and $[I(T_0),W]=b_W$ on this set. Its normalized squared-dimension mass is consequently at most $44/256$.

In contrast, let $M_j=\sum_W(\dim(W)/D)^2\theta_W^j$. By the Cauchy theorem all twist orders are powers of two, since $\operatorname{N}(D)=256$. Under the two quadratic coset averages, the real parts of $\theta^2$ differ only for $\theta$ of exact order sixteen, and there by $\pm\sqrt2$. Here and below the Galois relation for the squared vacuum row gives $\sigma_a(M_j)=M_{j/a}$, with exponents modulo the twist order. The second indicator formula gives
\[
 M_2=p(2+s_x x+s_z z),\qquad s_z=\nu_2(Z)=\pm1,\quad
 s_x=\nu_2(X)+\nu_2(Y)\in\{-2,0,2\}.
\]
Therefore $|M_2-M_2'|=\sqrt2(2-s_z)/4\ge\sqrt2/4$. The required order-sixteen mass is at least $1/4$, contradicting $44/256$.

\par\smallskip\noindent\emph{The two conductor-five parameters.}\enspace
The remaining parameters have $(f,c)=(20,4),(5,4)$. For the first $(u,v,c,s,d,e)=(-1,4,4,16,1,4)$ and for the second they are $(1,-1,4,4,9,19)$. Thus all that remains is
\[
\begin{array}{c|c|c|c|c}
 f&x&z&p&\operatorname{N}(D)\\\hline
 20&9+4\sqrt5&4+2\sqrt5&9/40-\sqrt5/10&1600\\
 5&6+3\sqrt5&14+6\sqrt5&3/20-\sqrt5/15&3600 .
\end{array}
\]
The following finite moment check disposes of both parameters, independently of $t$ and duality. It uses only the displayed dimension characters. For $f=5$ it also uses the twist masses already proved in the two-center argument of Theorem~\ref{thm:exchange-five-effective}. Denote the normalized masses at $1,-1$ by $P,N$. At $f=20$, induction of the unit gives $P\ge p^2+(p')^2+1/400+1/8$, $N\ge0$; the two negative-block reciprocal codegrees have sum $1/2$. At $f=5$ that same argument gives
\[
 P\ge p^2+(p')^2+1/25+1/8+9/100-2/70,\qquad N\ge1/8.
\]
Write $m_j=\operatorname{Re}\sum_W(\dim(W)/D)^2\theta_W^j$.
For each $j$ consider both quadratic coset averages of
$\operatorname{Re}(\theta^j)$. The signed Galois permutation identity
above gives $\sigma_k(m_j)=m_{j/k}$ for $k$ coprime to the twist
exponent. Since the moments used below belong to $K$ and inversion
preserves each quadratic coset, the two integrals are $m_j,m_j'$.
In both cases $m_1=p$ and $m_2=p(2+s_xx+s_z z)$ with the same
indicator ranges as before. At $f=5$ also use
\[
 \begin{gathered}
 m_3=p(1+a x+b z),\\
 a\in\left\{\frac{3r-4}{2}:r=0,\ldots,4\right\},\qquad
 b\in\left\{\frac{3r-19}{2}:r=0,\ldots,19\right\}.
 \end{gathered}
\]
Here $a=\operatorname{Re}(\nu_3(X)+\nu_3(Y))$ and
$b=\operatorname{Re}\nu_3(Z)$ arise from rotations of order dividing
three on spaces of dimensions $s=4$ and $e=19$, respectively.
For such a rotation on an $H$-dimensional space, if its eigenvalue-one
multiplicity is $r$, the real part of its trace is
$r-(H-r)/2=(3r-H)/2$. These values are rational, so $m_3\in K$,
as required for the quadratic averages. The indicators themselves
need not be real.

For clarity the finite comparison is specified as follows. A column consists of the constant $1$, the two averages at each power $j$ in increasing order, and the indicators of $\theta=1,-1$. Use powers one and two for $f=20$ and powers one through three for $f=5$. The Cauchy theorem restricts the prime divisors of the order to $2,5$ and $2,3,5$, respectively. Orders dividing $20$ or $180$, respectively, and one additional column $(1,0,\ldots,0)$ suffice: for any exponent exceeding those in these bounds every averaged power vanishes by cyclotomic trace. For a resulting primitive component of order $5h$ the two averages are $(-1\pm\sqrt5)\mu(h)/(4\varphi(h))$ in their two orders, determined by its fifth-root component. For order $h$ prime to five they are both $\mu(h)/\varphi(h)$. There are $10$ distinct columns at $f=20$ and $25$ at $f=5$.

The exact certificates in \texttt{verify\_exchange.py} consist of two integral linear forms for $f=20$ and six for $f=5$. Each is nonnegative on every root column and has nonpositive last two coefficients. For each of the six or $600$ possible indicator tuples, respectively, one of the forms evaluates strictly negatively on the required moment vector, even upon substitution of the stated lower bounds for $P,N$. This contradicts its being a positive weighted sum of root columns. All comparisons use rational operations and exact signs in $\mathbb Q(\sqrt5)$.
\end{proof}



\section{Perfect categories: generators and the quintic field}

% ===== ade.tex =====
\subsection{A simple of Frobenius--Perron dimension at most two}
\begin{theorem}\label{thm:perfect-small-simple}
Suppose that the tensor unit is the only invertible simple in a rank-five complex fusion category of irrational Frobenius--Perron dimension. If some nonunit simple has Frobenius--Perron dimension at most two, then the Grothendieck ring is the adjoint $\mathrm{SU}(2)_9$ ring.
\end{theorem}
\begin{proof}
Every proper fusion subcategory is integral by Proposition~\ref{prop:proper-integral}. The integral classifications through rank four \cite{ABDP} show that every integral category of ranks two through four contains a nonidentity invertible. Thus any nonunit simple generates the whole category. Let $X$ be a nonunit simple of dimension $d\le2$. Then $X X^*$ contains a nonunit, since its unit multiplicity is one and $d>1$, so $X X^*$ generates as well.

On two copies of the set of simple objects form the undirected bipartite multigraph with adjacency matrix
\[
 H=\begin{pmatrix}0&N_X\\N_X^{\mathsf t}&0\end{pmatrix}.
\]
The matrix $N_X N_X^{\mathsf t}=N_{XX^*}$ has connected graph: every simple occurs in a power of the self-dual tensor generator $XX^*$. Thus the five vertices on the first side belong to a single connected component of the bipartite graph, and every vertex on the second side has a neighbor. The vector of simple dimensions, repeated on both sides, is a positive eigenvector of $H$ with eigenvalue $d$, so the graph has norm $d$. If $d=2$, the one-dimensional Perron eigenspace is defined over $\mathbb Q$. Normalizing at the tensor unit gives rational, hence integral, dimensions for all simples. This contradicts irrationality of the total dimension.

For $d<2$ the connected graph is of type $A$, $D$ or $E$. For a direct verification in this size, an edge of multiplicity at least two, a cycle, a vertex of valency four, or a shortest path joining two trivalent vertices together with two additional neighbors at both ends gives a subgraph of norm at least two. In the last case a positive eigenvector of eigenvalue two is constant two along the path and one at its four leaves. Hence the graph is a path or a tree with one trivalent vertex. In the latter case the three arm lengths $a\le b\le c$ sum to nine, and positivity of $2I-H$ gives $1/(a+1)+1/(b+1)+1/(c+1)>1$ by eliminating the three path blocks. The only possibility is $(1,1,7)$, the diagram $D_{10}$, whose bipartition has sizes four and six. The two sides here have five vertices each, leaving the path $A_{10}$.

With the path vertices labelled $j=1,\ldots,10$, the positive eigenvector is proportional to $\sin(j\pi/11)$. Each tensor-unit vertex is an endpoint by the unit rule. All five values on one side are distinct, so all five simples have distinct dimensions and are self-dual. The two copies of any simple are paired by $j\leftrightarrow11-j$, the only other vertex with the same eigenvector entry. Thus $N_X$ is the five-vertex path with a loop at its last vertex, with the unit at the opposite end. Multiplication by $X$ successively gives $b_0=1$, $b_1=X$, $b_{i+1}=Xb_i-b_{i-1}$ for $1\le i\le3$, and $Xb_4=b_3+b_4$. These relations determine the whole based ring. In the adjoint $\mathrm{SU}(2)_9$ ring, multiplication by the simple labelled $8$ has just this graph, with labels in order $0,8,2,6,4$, so this ring realizes the result.
\end{proof}



% ===== noncomm.tex =====
\subsection{Characters at a non-self-dual pair}\label{sec:onepair}

Let $R$ have basis $1,Y,Z,X,X^*$, where $Y,Z$ are self-dual and
$X\ne X^*$. The ring is commutative by Proposition~\ref{prop:rank-five-commutative}.
Its five characters consist of three real characters and one complex
conjugate pair. Indeed the linear extension of duality has trace three,
and it permutes the primitive idempotents by complex conjugation.

\begin{proposition}
Write $\chi$ for a nonreal character. Up to complex conjugation,
its row and formal codegree are
\[
 (1,u,v,(k+\sqrt{-h})/2,(k-\sqrt{-h})/2),\qquad
 f_\chi=h=2+2u^2+2v^2+k^2,
\]
where $u,v,k$ are integers and $h\equiv0,3\pmod 4$.
For any real character with coordinates $(1,y,z,x,x)$ one has
$1+uy+vz+kx=0$.
If the unit is the only invertible basis element, then $h\ge7$.
\end{proposition}
\begin{proof}
The anti-invariant lattice is $\mathbb Z(X-X^*)$; it is preserved by
multiplication by $Y,Z,X+X^*$. Evaluating on $\chi$ yields the
integer eigenvalues $u,v,k$. Put $w=\chi(X)$.
Orthogonality of the two nonreal rows gives
$1+u^2+v^2+w^2+\overline w^2=0$. Consequently their common codegree
is $|w-\overline w|^2=2+2u^2+2v^2+k^2$. The norm
$w\overline w=(k^2+h)/4$ is a rational algebraic integer.
Orthogonality to a real row gives the asserted dimension relation.
If $h=3$, then $u=v=0$ and $|k|=1$, so $\operatorname{FPdim}(X)=1$.
If $h=4$, then $k=0$ and exactly one of $u,v$ has absolute value one;
the corresponding self-dual simple is invertible.
\end{proof}




% ===== selfdual.tex =====
\subsection{Classification of the quintic case}
\begin{theorem}\label{thm:quintic-classification}
If the Frobenius--Perron dimension character of a rank-five complex fusion category has degree five, then its Grothendieck ring is the adjoint $\mathrm{SU}(2)_9$ ring.
\end{theorem}
\begin{proof}
Theorem~\ref{thm:quintic} identifies the rational fusion algebra with $K=\mathbb Q(\zeta_{11})^+$. All five simples are self-dual and the five characters form a single Galois orbit. Upon Galois conjugating the spherical category we may therefore take its spherical dimension character to be the Frobenius--Perron character. Write $A\subset\mathcal O_K$ for its image order, $D$ for the global dimension, $p=1/D$ and $h=[\mathcal O_K:A]$. By the trace pairing and \eqref{eq:qmoment},
\[
 \mathrm N(D)=11^4h^2,\qquad \mathrm{Tr}(p)=1,\qquad
 \mathrm{Tr}(p^2)\le 3/11.
\]

We first show that $\mathrm N(D)$ is odd. The prime two is inert in $K$, since its order in $(\mathbb Z/11)^\times/\{\pm1\}$ is five. If $2\mid\mathrm N(D)$ then $2\mid h$, and the image of $A$ in $\mathcal O_K/(2)=\mathbb F_{32}$ is a proper subfield: surjectivity would imply $A\otimes\mathbb Z_2=\mathcal O_K\otimes\mathbb Z_2$ by Nakayama's lemma. That subfield is $\mathbb F_2$. Write $D=2^a\delta$, where $a>0$ and $\delta$ is a unit at two. For every $\gamma\in\mathrm{Gal}(K/\mathbb Q)$, $D/\gamma(D)$ lies in $A$, being the dimension of the forgetful image of the induction constituent of codegree $\gamma(D)$ \cite[Theorem~2.13]{Ostrik3}. Its reduction is a unit and therefore one. Thus the residue of $\delta$ is fixed by the whole Galois group, which acts on $\mathbb F_{32}$ as its full Galois group. We get $\overline\delta=1$. Reduction of $\mathrm{Tr}(\delta^{-1})=2^a\mathrm{Tr}(p)=2^a$ now gives $1=0$ in $\mathbb F_2$.

The Cauchy theorem for spherical categories \cite[Theorem~3.9]{RankFinite} gives odd Frobenius--Schur exponent $N$. Since $D\in\mathbb Q(\zeta_N)$ generates the field of conductor eleven, $11\mid N$. Let $\sigma$ act on $\zeta_N$ by the inverse of two modulo $N$. Its restriction to $K$ is nontrivial. The higher Gauss-sum formula of \cite[Theorem~4.4, equation~(4.17), in the arXiv version]{NSW} gives
\[
 \sigma(p)=\frac{\nu_2(\mathcal C)}{D}
  =\frac{1}{D}\sum_X \operatorname{FPdim}(X)\nu_2(X)
  \le \frac{1+2\sqrt{D-1}}{D}.
\]
Here $|\nu_2(X)|\le1$ for every simple $X$, and we used Cauchy's inequality on the four nonunit dimensions. If $D\ge324$, two distinct reciprocal conjugates $u=p$ and $v=\sigma(p)$ satisfy $u\le1/324$ and $v<37/324$, since $1/D+2/\sqrt D$ decreases. The function $u^2+v^2+(1-u-v)^2/3$ decreases in each variable on this rectangle. Cauchy's inequality on the other three conjugates yields
\[
 \operatorname{Tr}(p^2)\ge
 u^2+v^2+(1-u-v)^2/3\ge
 \frac{1+37^2+286^2/3}{324^2}>\frac3{11}.
\]
Thus $D<324$. Order the reciprocal conjugates as $p_1\le\cdots\le p_5$,
where $p_1=1/D$. If $p_j\le a<1/5$ for $j<5$, the first $j$ entries
have mean at most $a$. Applying Cauchy's inequality to the first $j$ and
last $5-j$ entries gives
\[
 \operatorname{Tr}(p^2)\ge ja^2+(1-ja)^2/(5-j).
\]
Here the right side decreases with $a<1/5$. For $(j,a)=(2,1/20)$,
$(3,1/10)$ and $(4,1/8)$ the bound exceeds $3/11$. Hence
$p_2>1/20$, $p_3>1/10$, $p_4>1/8$ and $p_5\ge1/5$ by the mean.
When $D\le100$ we obtain $\mathrm N(D)<100\cdot20\cdot10\cdot8\cdot5$.
When $D>100$ we also have $p_2>1/10$: otherwise the first two entries
are at most $1/100,1/10$, giving
\[
 \operatorname{Tr}(p^2)\ge 1/100^2+1/10^2+(1-1/100-1/10)^2/3>3/11 .
\]
Thus $\mathrm N(D)<324\cdot10\cdot10\cdot8\cdot5$ in this case.
Both norm bounds are below $11^6$, so $h<11$. The formal $d$-number
property makes $(D)$ Galois invariant. At any rational prime $\ell\ne11$
all primes above $\ell$ are unramified and have the same valuation on
$D$, whence $5\mid v_\ell(\mathrm N(D))=2v_\ell(h)$.
Since $h<11$, this implies $h=1$.

Write $b_j=\zeta_{11}^j+\zeta_{11}^{-j}$ for $1\le j\le5$.
These elements are an integral basis, and
$\operatorname{Tr}(b_i b_j)=11\delta_{ij}-2$.
Write $p=-\sum v_jb_j$, where $\sum v_j=1$. Since $A=\mathcal O_K$, each $b_j$ belongs to the dimension order, and the coefficient-of-unit trace gives $t_j=11v_j=2-\operatorname{Tr}(p b_j)\in\mathbb Z$. We have $\sum t_j=11$ and $\sum t_j^2=11(2+\mathrm{Tr}(p^2))\le25$. The minimum integral square sum at sum eleven is 25, attained only by permutations of $(2,2,2,2,3)$. Hence $p=(2-b_j)/11$; the distinguished embedding minimizes $p$, so $j=1$.

The adjoint $\mathrm{SU}(2)_9$ ring has the same global dimension $11/(4\sin^2(\pi/11))$. The simple labelled $8$ has dimension $2\cos(\pi/11)=-b_5$, generating $\mathcal O_K$. Thus the two sets of simple dimensions are orthonormal integral bases for the same lattice $(\mathcal O_K,(x,y)\mapsto\mathrm{Tr}(pxy))$. Each squared-length-one integral vector is a signed basis vector; positivity in the distinguished embedding fixes every sign. The two based rings coincide.
\end{proof}



% ===== character_tools.tex =====
\section{Character lattices and twist support}
\label{sec:character-tools}
The field-by-field arguments use three common tools: the discriminant
of a projected character lattice, the local structure at a rational
character of codegree divisible by a prime only once, and the support
of central twists. We establish them together before applying them.
For the small-codegree restrictions we assume rank five and perfectness;
the other statements retain their indicated generality.

\begin{lemma}\label{lem:orbit-lattice} Let $L$ be the field of one Galois orbit of real characters, $J$ its codegree in one embedding and $v=\mathrm{Tr}_{L/\mathbb Q}(1/J)$. Orthogonally project the based ring onto the rational simple ideal of this orbit for the coefficient-of-unit form. The image $B$ is an order with form $\mathrm{Tr}_{L/\mathbb Q}(ab/J)$. By unimodularity its dual is the intersection of the based lattice with the ideal. If $n$ is the Gram determinant of this intersection, then
\[
  \operatorname N(J)=n\operatorname{disc}(B),\qquad
  \operatorname N(J)\ge\operatorname{disc}(L)/v .
\]
\end{lemma}
\begin{proof}
The determinant of the form on $B$ is $\operatorname{disc}(B)/\operatorname N(J)$ and the determinant on its dual is the reciprocal. The squared norm of the projection of the unit is $v$. Its denominator divides $n$ by the Gram formula for orthogonal projection, so $nv$ is a positive integer.
\end{proof}

\begin{lemma}\label{lem:r5-local-partner}
Let $A$ be a commutative based ring and $\rho$ a rational character of
codegree $f$ with $v_\ell(f)=1$. In the completed local factor at
$\rho\bmod\ell$, the total quotient algebra is
$\mathbb Q_\ell\times L$ for a single totally ramified field $L$, and
the order is $\mathbb Z_\ell\times_{\mathbb F_\ell}\mathcal O_L$.
\end{lemma}
\begin{proof}
The coefficient-of-unit trace is associative and perfect, including
on each local factor $A_0$ of $A\otimes\mathbb Z_\ell$. In $A_0$ the
intersection with the rational axis is $\ell\mathbb Z_\ell$: the
primitive idempotent has numerator $\sum_i\rho(b_i)b_i^*$ with unit
coefficient one and denominator $f$. The projection $B$ off the
rational axis is consequently a local order with quotient
$\mathbb F_\ell$ and $A_0=\mathbb Z_\ell\times_{\mathbb F_\ell} B$.
Let $\mathfrak m$ be the maximal ideal of $B$. If $B$ is not normal,
choose $y$ in its normalization, outside $B$, whose class is killed by
$\mathfrak m$. In fact $\mathfrak m y\subset\mathfrak m$ by
lying-over, on each branch of the normalization. The classes of
$(\ell,0),(0,\ell y)$ in $A_0/\ell A_0$ are in the socle and are
linearly independent: dividing a relation by $\ell$ first forces the
second coefficient to vanish modulo $\ell$, since $y\notin B$, then
forces the first to vanish by the fiber condition. Perfect duality
identifies the socle with the dual of the residue field, of dimension
one. Thus $B$ is normal, hence a single discrete valuation ring since
it is local. Its residue field is $\mathbb F_\ell$. This extends the
local argument of \cite[Lemma~8.5]{Rank4}.

The same socle argument applies to a local fiber product
$O\times_{\mathbb F_\ell}B$ with a perfect associative pairing,
where $O$ is any discrete valuation ring finite over $\mathbb Z_\ell$
with residue field $\mathbb F_\ell$. If $B$ were nonnormal, the
same classes $(\ell,0)$ and $(0,\ell y)$ would contradict the
one-dimensional socle modulo $\ell$. Thus $B$ must again be a
valuation ring with residue field $\mathbb F_\ell$. We use this
version when $O$ is a ramified quadratic integer ring.
\end{proof}

\begin{lemma}\label{lem:small-complement}
In a perfect rank-five fusion category, no character has codegree two.
A real quadratic pair with common rational codegree $h$ and field
discriminant $\Delta$ satisfies $h\ge\Delta\ge5$.
A nonreal pair cannot have common codegree three or four, and two
rational characters cannot both have codegree four.
\end{lemma}
\begin{proof}
For two real quadratic
conjugate rows $\chi,\chi'$ of common codegree $h$ and field
discriminant $\Delta$, write $\chi_i-\chi'_i=b_i\sqrt\Delta$
with $b_i\in\mathbb Z$. Orthogonality gives
$\Delta\sum b_i^2=2h$. The difference is perpendicular to the positive
dimension row, so has at least two nonzero entries. Thus $h\ge\Delta
\ge5$. For a nonreal character of rational codegree $h\le4$, its Galois
orbit is quadratic: a larger nonreal orbit has at least four members
and would exhaust the reciprocal mass. Write its field discriminant
as $-\Delta$. Orthogonality gives
$h=\Delta\sum_j b_j^2$, where the sum runs over dual pairs and
$\chi_j-\overline{\chi_j}=b_j\sqrt{-\Delta}$.
Thus $h=3$ or $4$ forces respectively $\Delta=3$ or $4$ and exactly
one nonzero $b_j$. If $h=3$, the unit and the two nonreal entries of
norm one exhaust the norm budget; their real parts are $\pm1/2$.
Orthogonality to the positive row forces that pair to be invertible.
If $h=4$, the nonreal entries are $i,-i$. The remaining norm budget
is one, so the two other simples cannot form a dual pair; a single
entry $\pm1$ forces an invertible. Finally
two rational rows of codegree four each have the unit and three
further entries $\pm1$. Orthogonality forces identical supports, since
different supports intersect in three places. Their sum has support
on the unit and one additional simple, forcing that simple to be
invertible. A rational codegree-two row likewise forces an invertible;
a quadratic pair of codegree two uses all reciprocal mass.
\end{proof}


% ===== odd_row.tex =====
\begin{lemma}[Integral character blocks and twist support]
\label{lem:prime-local-block}
Let $\mathcal D$ be a modular fusion category with positive spherical
dimensions, multiplicative central charge one, normalized matrix $S$,
and twist exponent
$N=\ell^a m$, where $\ell$ is prime and $(\ell,m)=1$.
Here central charge one means
$\sum_Y\theta_Y\dim(Y)^2=\sqrt{\dim(\mathcal D)}$, so the ordinary
twist matrix satisfies $STS=T^{-1}ST^{-1}$. This holds for every
Drinfeld center used below. For a label $J$, let $e_J$ be its
primitive character idempotent.
Suppose that a set $\mathcal J$ of labels has real $S$-rows and that
\[
 e_{\mathcal J}=\sum_{J\in\mathcal J}e_J
       \in K_0(\mathcal D)\otimes\mathbb Z_{(\ell)}.
\]
Then $T^m$ preserves the span of these rows. If restriction of this
span to the twist-one columns is injective, every row in the span
vanishes at twists whose orders are divisible by $\ell$.
\end{lemma}
\begin{proof}
We give the prime-local version of the indicator argument used in
\cite[Section~3]{Dyadic}. Choose a genuine object $V$ whose
multiplicities agree modulo $\ell$ with the rational coefficients of
$e_{\mathcal J}$. In split character coordinates its values reduce to
one on $\mathcal J$ and zero elsewhere. Thus
$V^{\ell^n}\to e_{\mathcal J}$ over the completed cyclotomic field.

Let $L_k^i$ be the generalized indicator with insertion $(i,i^*)$ in
$Z(\mathcal D)$ and second index one. The insertion has twist one.
The modular indicator formula \cite[Proposition~6.1]{NS} gives
\[
 L_k^i(e_{\mathcal J})
   =\sum_{J\in\mathcal J}|(ST^kS^{-1})_{iJ}|^2.
\]
Indeed, for each $J$, insert
$S_{0J}^{-2}e_J=\sum_z(S_{zJ}/S_{0J})[z]$ in that formula and use
multiplicativity of the real character $S_{-J}/S_{0J}$. A real row
has $J=J^*$, so $S$ and $S^{-1}$ have the same $J$-column.

Choose $b\ge\max(1,a)$ with $\ell^b\equiv1\pmod m$.
For $n\ge2b$, grouping the $m\ell^n$ strands of $V$ gives
\[
 L_m^i(V^{\ell^n})\equiv
 L_{m\ell^b}^i(V^{\ell^{n-b}})
                         \pmod{\ell^{n-2b+1}}.             \tag{*}
\]
To make the grouping explicit, put $M=m\ell^n$ and let
$R_n=E^{(M,1)}_{(i,i^*),V}$ be the one-step rotation. By
\cite[Lemma~2.5 and Proposition~2.8(ii)]{NS}, the two sides of (*)
are respectively $\operatorname{Tr}(R_n^{\ell^n})$ and
$\operatorname{Tr}(R_n^{\ell^{n-b}})$.
The insertion has twist one, so $R_n^M=1$
\cite[Lemma~2.7]{NS}. Its eigenvalues are therefore $M$-th roots
of unity. All power traces lie in $\mathbb Q(\zeta_N)$
\cite[Proposition~5.7]{NS}, so their multiplicities are constant
on Galois orbits over that field.
The two powers agree on the prime-to-$\ell$ components and on the
components of order at most $\ell^{n-b}$. On a component of order
$\ell^t$ with $t>n-b$, the units congruent to one modulo $\ell^b$
and fixing the $m$-th roots have orbits of size $\ell^{t-b}$.
Both powers are constant on each orbit, and its size is divisible
by $\ell^{n-2b+1}$. This proves (*).

Take the $\ell$-adic limit. Since $N\mid m\ell^b$, linearity gives
\[
 \sum_{J\in\mathcal J}|(ST^mS^{-1})_{iJ}|^2
                          =\mathbf1_{i\in\mathcal J}.
\]
Unitarity now shows that $ST^mS^{-1}$ preserves the coordinate
subspace on $\mathcal J$. The rows are real, so $T^m$ preserves
their span. On a twist-one column $T^m$ acts trivially. If restriction
to those columns is injective, $T^m$ is therefore the identity on
the span, which proves the support assertion.
\end{proof}

\begin{lemma}[Twist support of a rational row]\label{lem:odd-row}
In the same modular setting, suppose that every entry of the
normalized $S$-row $S_{J,-}$ is rational, and write
$q=S_{0J}^{-2}\in\mathbb Z_{>0}$ for its formal codegree.
This row is supported only at twists whose orders
have prime divisors among those of $q$. If $\theta_J=1$ and
$(k,q)=1$, then
\[
                    (ST^kS)_{Ji}=S_{Ji}
\]
for every twist-one label $i$. In particular, a rational row of
odd formal codegree has odd twist support and satisfies this identity
at $k=2$.
\end{lemma}
\begin{proof}
The character $S_{-J}/S_{0J}$ takes integer values. For every prime
$\ell\nmid q$, its primitive idempotent is therefore rational and
$\ell$-integral. The vacuum entry is nonzero, so restriction of this
one-dimensional row space to twist-one columns is injective.
Apply Lemma~\ref{lem:prime-local-block} for each such prime.

Choose a unit $u$ modulo $N$ equal to $k$ on the primary factors of
$N$ belonging to primes of $q$, and equal to one on the remaining
factors. The support assertion makes the $J$-rows of $ST^kS$ and
$ST^uS$ equal. In $\mathrm{SL}_2(\mathbb Z/N)$,
\[
 ST^uS=T^{-u^{-1}}\operatorname{diag}(u^{-1},u)ST^{-u^{-1}}.
\]
The diagonal element fixes the $J$-coordinate vector by rationality
of its modular row, and the outer twist phases are one at $(J,i)$.
This proves the identity.
\end{proof}
In the applications below, $J=A_\rho$ is the central induction
label of a rational character $\rho$ of the original fusion ring.
Its normalized row satisfies
\[
 S_{A_\rho,W}=\frac{\rho(F(W))}{f_\rho}\in\mathbb Q,
\]
since $\rho$ is integer-valued. Thus the required rationality holds
for the normalized entries themselves.

\section{Perfect categories in quadratic degree}
\label{sec:perfect-quadratic}

% ===== qsmall.tex =====
Throughout this section, $\mathcal C$ is a perfect rank-five fusion
category with quadratic field $K=\mathbb Q(D)$, and every nonunit
Frobenius--Perron dimension exceeds two. We work in its positive
spherical Galois conjugate. The argument first reduces the possible
character profiles to $\mathbb Q(\sqrt5)$, then uses twist moments to
exclude the two remaining profiles.

\begin{lemma}[An inert-prime obstruction]
\label{lem:quadratic-inert-three}
Suppose $K=\mathbb Q(D)$ is quadratic and three is inert in $K$.
The three complementary characters cannot all be rational with
codegrees $3,t,g$ and $3\nmid g$.
\end{lemma}
\begin{proof}
Let $\rho,\tau,\gamma$ be these characters. The $\gamma$-idempotent
is integral at three, while the inert quadratic branch cannot be the
local partner of $\rho$ in Lemma~\ref{lem:r5-local-partner}.
Hence the local factor containing $\rho$ consists of the two rational
branches $\rho,\tau$.

Let $A_\rho$ be its unit-induction constituent. The character formula
for the integral functional $x\mapsto3\langle x,F(A_\rho)\rangle$ is
\[
 d(x)+d'(x)+a\rho(x)+b\tau(x)+c\gamma(x),
       \qquad a,b,c\in\mathbb Z,\quad 3\mid c.
\]
Indeed the two dimension-orbit entries of the rational modular row
are $1/3$, and $3S_{A_\rho,B}$ is integral for every $B$.
Symmetry makes both $3S_{A_\rho,A_\gamma}$ and
$gS_{A_\rho,A_\gamma}$ integral, so $3\mid c$.
Evaluate the functional on the identity of the completed local factor
of $\rho,\tau$. It gives $a+b\equiv0\pmod3$. Evaluating on the
global unit instead gives $2+a+b+c\equiv0\pmod3$, a contradiction.
\end{proof}

\begin{proposition}[Quadratic profile reduction]
\label{prop:quadratic-profile-reduction}
Every such category is self-dual. If all three complementary formal
codegrees are rational, then $K=\mathbb Q(\sqrt5)$. Otherwise
$K=\mathbb Q(\sqrt5)$ and their multiset is one of
\[
 \{g,5+\sqrt5,5-\sqrt5\}\quad(4\le g\le8),\qquad
 \left\{4,\frac{15+3\sqrt5}{2},\frac{15-3\sqrt5}{2}\right\}.
\]
\end{proposition}
\begin{proof}
\par\smallskip\noindent\emph{Dimension characters.}\enspace
Write $c=\operatorname{disc}(K)$. Every nonunit generates the category,
and the positive eigenline of its connected fusion graph is defined
over its eigenvalue field. Hence every nonunit dimension generates
$K$. Write
\[
 d_i=(a_i+b_i\sqrt c)/2,\qquad D=(A+B\sqrt c)/2,\qquad N=DD'.
\]
All four $b_i$ are nonzero. From conjugacy, orthogonality and the
definition of $D$, with all sums in the following display taken over
the four nonunit simples,
\[
 A=c\sum b_i^2=4+\sum a_i^2,\qquad B=\sum a_i b_i>0,\qquad
 D\ge A/2\ge2c,\qquad N\mid A^2 .
\]
Here $B>0$ follows from the strict inequality $D>D'$ by spectral domination.
Every complementary formal codegree lies in $K$. These values either
are three rational integers $f\le g\le j$, or one rational integer
$g$ and a conjugate irrational pair $H,H'$. In the latter case the
characters giving $H,H'$ have field $K$, hence are real. Write
$H=(h+b\sqrt c)/2$, $n=HH'$, $b>0$. Subtracting the two character
rows gives $h=c\sum b_i'^2\ge2c$, since their difference is perpendicular
to the positive row. As before $n\mid h^2$. In both cases put
$v=\operatorname{Tr}(1/D)$. Put $t=0$ for the rational profile, and
$t=(h^2-4n)/n^2$ for the mixed profile. The simultaneous Fourier bound
on the quadratic pairs reads
\[
 \delta:=v^2+t-1/c\ \le\ 4/N .
\]
Indeed the normalized difference of the two dimension-orbit vectors in the ordinary center has diagonal character of conductor $c$. Its $S$-image on the twist-one induction constituents has squared norm at least $(1/D-1/D')^2+t$, so Lemma~\ref{lem:hecke-projection} applies. Using $(1/D-1/D')^2=v^2-4/N$ gives the displayed inequality. The sign of $\delta$ determines whether this inequality gives a
finite norm search. We first recall the row restrictions used to
enumerate the profiles.

\par\smallskip\noindent\emph{Small rational codegrees.}\enspace
Suppose first that all three complementary codegrees are rational.
They are integers at least three: a rational row of norm two forces
an invertible, and a conjugate pair of codegree two exhausts the
reciprocal mass. Put $s=v=\operatorname{Tr}(1/D)$.
Proposition~\ref{prop:ordinary-fourier} gives $s-2/D\le1/\sqrt c$.
The following arguments treat a rational codegree as a number, allowing it to be shared by Galois conjugate characters. A formal codegree three or four repeated on a quadratic pair is excluded
by the difference-of-rows argument in Lemma~\ref{lem:small-complement};
this uses only the size of the rows. Two rational rows of codegree four
likewise force an invertible. Two rational rows of codegree three each
have one positive and one negative nonunit entry, and their supports
meet in exactly one nonunit position, where the signs are opposite.
After relabeling, the rows are $(1,-1,1,0,0)$ and $(1,0,-1,1,0)$.
The three corresponding dimensions are $z+2,z+1,z$. A third real
rational row is
\[
 (1,u+1,u,u-1,v),\qquad h=3+3u^2+v^2 .
\]
If $v=0$, orthogonality to the Perron row gives $1+u(z+1)=0$,
which is impossible. If $u=0$, the last simple has rational
dimension $3/|v|$. For $h=4$ this would be an object $W$ of
dimension three. The three dimensions in the chain have equal
nonzero irrational parts, so $W^2$ can contain only the unit and
$W$, giving $9=1+3m$. Values $h=5,6$ do not occur with $v\ne0$.
These facts exclude all triples with
$0<s<4/21$, namely $(3,3,4),(3,3,5),(3,3,6),(3,4,4)$.
Hence $s\ge4/21$.


\par\smallskip\noindent\emph{Two codegrees equal to three.}\enspace
There is a sharper obstruction for the profile $(3,3,h)$ when $3\nmid h$.
Write $\rho,\tau,\psi$ for the rows of codegrees $3,3,h$. The two
codegree-three rows have distinct reductions at three, since their
squared difference is six. The $\psi$-idempotent is integral at three.
Hence Lemma~\ref{lem:r5-local-partner} requires three to split in $K$,
one local branch pairing with each of $\rho,\tau$. In particular
$\mathrm{Tr}_K(d_X)\equiv\rho(X)+\tau(X)\bmod3$ for every simple $X$.

In the normalized center matrix $S$, modular symmetry and coprimality
make $S_{A_\psi,A_\rho}$ and $S_{A_\psi,A_\tau}$ integers of absolute
value less than one, hence zero. The two entries joining either
codegree-three row to the dimension-orbit labels are $1/3$.
Since every $F(A_\chi)$ contains the unit once, the block on
$A_\rho,A_\tau$ is therefore
$\tfrac13\left(\begin{smallmatrix}a&1-a\\1-a&a\end{smallmatrix}\right)$
for some integer $|a|<3$. The coefficient of the first nonunit
simple $X$, where $\rho(X)=-1,\tau(X)=0$, in $F(A_\rho)$ is
$(\mathrm{Tr}_K(d_X)-a)/3$. Thus $a\equiv-1\bmod3$.

The rational row $u_Z=3S_{A_\rho,Z}$ has squared norm nine, and the
anomaly-free modular relation $STS=T^{-1}ST^{-1}$ gives
$\sum_Z u_Z^2\theta_Z=3a$. For $a=-1$ there is mass seven at twist
one from $I(1)$, leaving mass two to reach $-3$, which is impossible.
Thus $a=2$. The five entries on $I(1)$ are now $(1,1,2,-1,0)$, so
exactly two further simples have nonzero entry, each $\pm1$.
Their twists sum to $6-7=-1$, hence are the two primitive cube roots.
Applying $STS=T^{-1}ST^{-1}$ at $(A_\rho,\mathbf1)$ gives
\[
 \sum_{j=1}^2 u_j\frac{\dim Z_j}{D}\theta_{Z_j}
       =1-(s+1/3)=1/3+1/h .
\]
The imaginary parts force the two real coefficients to agree; their
common value is $-(1/3+1/h)$. The categorical dimensions are positive,
so $D/\dim Z_j=3h/(h+3)$.
Its square is a formal codegree of the modular center, hence it is
an algebraic integer and therefore an integer. This requires
$h+3\mid9$, impossible for $h\ge3$ coprime to three.

\par\smallskip\noindent\emph{Finite profile ranges.}\enspace
We now enumerate the possible complementary profiles, independently of
$D$. In the rational case, $v=s=1-1/f-1/g-1/j$. Since $D\ge j$ and
$D\ge2c$, the inequality $s-2/D\le 1/\sqrt c$ gives
\[
 4c/21\le\sqrt c+1,\quad
 1-5/f\le1/\sqrt c,\quad 1-1/f-4/g\le1/\sqrt c,\quad
 1-1/f-1/g-3/j\le1/\sqrt c.
\]
Hence $c\le37$, $f\le9$, and $g\le18$. Only $(c,f,g)=(5,3,3)$,
$(8,3,3)$ and $(5,3,4)$ can give an unbounded $j$ in the last
inequality. In the first two cases the two rational codegree-three
rows have different reductions modulo three. Both need local partners
and three is inert, leaving only one available axis. In the last case a second four is impossible. For $j>4$ the three
complementary characters are rational, and
Lemma~\ref{lem:quadratic-inert-three} applies with $g=4$.
Thus all three exceptional pairs can be discarded at once.

In the mixed case let $w=h/n$, $u=v-2/D$, $L=u+w=1-1/g-2/D$.
Lemma~\ref{lem:orbit-lattice} gives $n\ge c/w$. Thus
\[
 1\ge c(u^2+w^2)-4w\ge cL^2/2-2L-2/c .
\]
For $c\ge17$, $L\ge2/3-1/c$ makes the right side greater than one
(the lower bound is increasing in $c$). This leaves $c=5,8,12,13$.
Also $D\ge\sqrt n$ by spectral domination, and $u^2+w^2\le1/c+4/n$.
It follows that $(2/3-2/\sqrt n)_+^2\le2/5+8/n$, so
$\sqrt n\le30+\sqrt{990}$ and $n\le3777$.
For each $n$ take $2c\le h<n$, $c\mid h$, $n\mid h^2$,
$h^2-4n=cb^2>0$, $h-b\sqrt c>2$.
Discard $t>1/c$. For this finite list one checks
$1-w>\sqrt{1/c-t}$, so $D\ge g$ and the Fourier bound give the finite
range $1-w-3/g\le\sqrt{1/c-t}$. At $(c,g)=(5,3)$ neither quadratic
component can furnish the rational three-row with a totally ramified
partner.

\par\smallskip\noindent\emph{Profiles with a finite norm bound.}\enspace
There are 1299 rational and 614 mixed profiles with $\delta>0$.
In each take $N\le4/\delta$ and $A=Nv\in c\mathbb Z$, $A\ge4c$.
Require $N\mid A^2$, $A^2-4N=cB^2>0$, $D'>1$, $D>17$.
For each rational codegree $j$ impose $D/j\in\mathcal O_K$; for the
irrational pair use $n\mid N$ and the integral traces of $D/H,D/H'$.
The resulting sixteen cases are all rational. A repeated four is
impossible, including a cubic orbit as explained below. For
$(3,3,j)$ we require $3\mid j$ and $j=3+3u^2+v_0^2$ with
$u v_0\ne0$: $u=0$ would give a rational nonunit dimension.
At $(c;f,g,j)=(8;3,4,16)$ the inert prime three prevents a partner
for the three-row. The five cases left by these restrictions are
\[
\begin{array}{c|rrr|c}
 c&A&B&N&(f,g,j)\\\hline
 5&720&320&1600&(4,5,10)\\
 12&672&192&2304&(3,3,24),\ (3,4,8)\\
 12&48&12&144&(3,4,12),\ (3,6,6).
\end{array}
\]
None satisfies the three displayed dimension-vector identities with
$d_i>2$ and $b_i\ne0$. For example, the exhaustive integer test takes
$|a_i|\le\sqrt{A-4}$ and $|b_i|\le\sqrt{A/c}$, tests integrality and
positivity, and matches the triples
$(\sum a_i^2,\sum b_i^2,\sum a_i b_i)$ from two pairs of coordinates.
The exact profile and vector enumeration is
\texttt{quadratic\_profiles.py}.

\par\smallskip\noindent\emph{Self-duality of the remaining profiles.}\enspace
In the remaining profiles $\delta\le0$. There are 162 rational profiles
and six mixed profiles in the above lists. If a non-self-dual pair
occurs, its two nonreal rows have a repeated rational codegree
$j\ge7$, $j\equiv0,3\pmod4$ by Section~\ref{sec:onepair}.
The only entries with a repeated value at least seven are
$(3,7,7),(3,8,8),(3,9,9)$ at $c=5$. Nine has the wrong congruence.
In the other two cases the paired idempotent is integral at three,
and the remaining quadratic component is inert. The rational
codegree-three row therefore cannot have a local partner. This proves
self-duality in every profile still under consideration.

\par\smallskip\noindent\emph{Removing the larger quadratic fields.}\enspace
For a remaining rational profile, $\delta\le0$ says precisely that
$s\le1/\sqrt c$, and hence $c\le27$.
We exclude every $c\ge13$. In that range
$s\le1/\sqrt{13}<0.282$. With three integer codegrees at least
three the only possibilities are $(3,3,h)$ with $h\le19$,
$(3,4,h)$ with $4\le h\le7$, $(3,5,5)$, and $(4,4,4)$.
For $(4,4,4)$ the complementary rows cannot include a pair of rational
rows or a quadratic orbit. If they are a single cubic orbit, their
character field is cyclic cubic of discriminant at least $49$.
Lemma~\ref{lem:orbit-lattice} with $v=3/4$ would give
$64=\mathrm N(4)\ge49/(3/4)>64$. This excludes the triple.

We use a consequence of Lemma~\ref{lem:r5-local-partner}.
If a prime $\ell$ divides a rational codegree $f$ exactly once
and all other complementary codegrees are units at $\ell$, then
$\ell$ must ramify in $K$. The remaining complementary idempotent
is integral (use its Galois-invariant sum for a pair). An inert
quadratic branch cannot supply the totally ramified partner. In
the split case the other quadratic branch would be an isolated factor
with nonintegral primitive idempotent, since $D/f$ is integral.
This proves the claim.

In $(3,3,h)$ the preceding twist argument leaves $3\mid h$.
For $h\le19$ the row formula gives only $h=15$ with $u,v\ne0$,
or $h=12$ with $u=0$, giving an invertible of dimension one.
The surviving $h=15$ forces five to ramify and
$s=4/15$ forces $c<15$ by the Fourier bound. Five does not ramify
at the only remaining discriminant $13$.
For $(3,4,5)$, both three and five must ramify, whereas no positive
fundamental discriminant at most $28$ is divisible by fifteen.
For $(3,4,7)$ or $(3,5,5)$, the respective values of $s$ force
$c<15$ and three would have to ramify at $c=13$. The case
$(3,4,4)$ was excluded above.
Finally $(3,4,6)$ requires $c<17$ and hence $c=13$ in this range.
At two its codegree-six row needs a local partner. The quadratic
field is inert, the codegree-three axis splits off, and the rational
codegree-four axis cannot be the partner: in a fiber product of two
copies of $\mathbb Z_2$ over $\mathbb F_2$ both primitive idempotents
have denominator valuation one, whereas its denominator valuation is
two. This excludes $c=13$ and completes the proof that
$c\in\{5,8,12\}$ in the perfect rational-complement case.

\par\smallskip\noindent\emph{Discriminants twelve and eight.}\enspace
At $c=12$ one has $s\le1/\sqrt{12}<0.294$.
For $(3,3,h)$ this leaves $h\le25$, $3\mid h$. In
$h=3+3u^2+v^2$ the possibilities with $u,v\ne0$ are $15,24$,
and $u=0$ would give $h=12$ and an invertible. The value $15$
requires five to ramify. The other sorted profiles surviving the
preceding tests are $(3,4,6),(3,4,8),(3,5,5)$.
At $(3,4,6)$ the rational codegree-six row cannot pair at two
with the codegree-four row. If it paired with the single ramified
quadratic branch, the codegree-four row would be isolated, again
impossible. A rational row of codegree five must have all entries
$\pm1$, since the alternative support with a single nonunit entry
$\pm2$ contradicts Perron orthogonality. Two such rows cannot be
orthogonal in odd length. Thus in $(3,5,5)$ the last two rows form
a real quadratic pair. Their difference identity gives
$\Delta=5$ and $\sum b_i^2=2$, so the projected order is
$\mathcal O_{\mathbb Q(\sqrt5)}$. At five the intersection with
its axis is the trace dual for $\mathrm{Tr}(xy/5)$, namely the
prime $\mathfrak P$, because the different is $\mathfrak P$. The
codegree-three component splits off at five. This leaves the fiber
product $\mathcal O_{\mathbb Q_5(\sqrt5)}\times_{\mathbb F_5}B$
with $B$ a local order in $K_5$. If $B$ were nonnormal, the
classes $(5,0),(0,5y)$ in the proof of
Lemma~\ref{lem:r5-local-partner} would contradict perfectness.
Hence $B$ is a valuation ring with residue field $\mathbb F_5$,
contrary to the inertness of five in $\mathbb Q(\sqrt3)$.

The two further profiles $(3,3,24),(3,4,8)$ each have
$s=7/24>1/\sqrt{12}$, and therefore were covered by the finite
norm test. This eliminates $c=12$.

At $c=8$, $s<0.359$. If a codegree three occurs, it is rational.
Three is inert in $K$. With a complementary quadratic pair of common
codegree $h$, the inequality $2/3-2/h<0.359$ forces $h\le6$.
The real-quadratic row identity leaves only $h=5$, and its paired
idempotent splits off at three, leaving no local partner for the
rational row. Thus all rows are rational. Distinct codegree-three
rows cannot both find partners at this inert prime. The local partner
of a single such row has codegree $h\ge15$, $h\equiv6\bmod9$,
by the squared norm of the difference of the congruent rows.
Writing $k$ for the third codegree, the bound on $s$ forces $k=4$
and $h=15$. Five would then ramify, contrary to $c=8$.
If all codegrees are at least four, the same mass bound leaves a
repeated four, excluded above, or $(4,5,5)$. In the latter the
preceding local argument for the quadratic codegree-five pair again
excludes the inert prime five, the codegree-four axis splitting off.
Thus only $c=5$ remains.
\end{proof}



% ===== qmixed.tex =====
\begin{lemma}\label{lem:quadratic-mixed-profile}
If a complementary formal codegree is irrational, the only profile
left by Proposition~\ref{prop:quadratic-profile-reduction} is
$4,(15+3\sqrt5)/2,(15-3\sqrt5)/2$.
\end{lemma}
\begin{proof}
It remains to exclude the pair $5\pm\sqrt5$. Write its conjugate rows
$\chi_i=(a_i+b_i\sqrt5)/2$ with $H=5+\sqrt5$.
Their orthogonality and codegrees give
$\sum a_i^2=10$, $\sum b_i^2=2$ and $\sum a_i b_i=2$.
The two nonzero $b_i$ are $\pm1$. At both positions $a_i=b_i$;
one of the two remaining nonunit positions has $a_i=\pm2$ and
the other has $a_i=0$. In the scalar product with the positive
dimension row the two contributions proportional to
$(1+\sqrt5)/2$ cancel, since the difference of the two character
rows has scalar product zero. The remaining equation
$1\pm d_X=0$ gives an invertible, contrary to perfectness.
\end{proof}



% ===== mixed_polar.tex =====
\begin{proposition}\label{prop:mixed-polar}
A perfect self-dual rank-five fusion category with all four nonunit
dimensions greater than two cannot have a quadratic dimension field
and a complementary character of irrational formal codegree.
\end{proposition}
\begin{proof}
Use the unique remaining profile above. Put $p=1/D$, $p'=5/12-p$,
$a=1/H$ and $a'=1/H'$, so $a+a'=1/3$ and $aa'=1/45$.
Order the five multiplicity-one constituents of $I(1)$ as
$0,0',J,J',R$ according to their codegrees $D,D',H,H',4$.

The two real rows indexed by $J,J'$ have odd twist support by
Lemma~\ref{lem:prime-local-block}. Indeed their summed primitive
idempotent is rational and integral at two: the denominators
$H^2,(H')^2$ are units there. Their restrictions to the twist-one
columns $0,0'$ form the invertible matrix
$\left(\begin{smallmatrix}a&a'\\a'&a\end{smallmatrix}\right)$.
Write $N=2^s m$ with $m$ odd.
Let $h$ be their sum, with coordinates $h_B=S_{JB}+S_{J'B}$.
In the factorization of $ST^kS$ from Lemma~\ref{lem:odd-row}, choose
$k=2\pmod m$ and $k=1\pmod{2^s}$. The diagonal Galois action
permutes $J,J'$ without signs, by their character-row expressions.
Hence the sum of their two rows in $ST^2S$ equals $h$ on the
twist-one coordinates.

Write $w_B=S_{0B}$ and $r=S_{RJ}=S_{RJ'}\in\tfrac14\mathbb Z$;
the equality follows from rationality of the $R$-row and unsigned
Galois conjugacy. On our five induction constituents $I$ the two
real vectors restrict to
\[
 w|_I=(p,p',a,a',1/4),\qquad
 h|_I=(1/3,1/3,2/3-r,2/3-r,2r).
\]
Indeed each of the five rows sums to one on $I$, since the underlying
class of $I(1)$ is the Casimir $\sum_X XX^*$ and its character value is
the corresponding formal codegree. Symmetry and Galois conjugacy now
give the displayed restrictions.

For $t\in\mathbb R$ put mass $(w_B+t h_B)^2$ at each twist
$\theta_B$. Its total mass is $1+2t^2$ by orthonormality. Put
$M_2=\Re\sum_B \theta_B^2 w_B^2$; the indicator formula gives
$M_2\le (1+2\sqrt{D-1})/D$.
The first moment of the measure is
$p+2t/3+2(2/3-r)t^2$ by $STS$; the second moment has real part
$M_2+2t/3+2(2/3-r)t^2$ by the dyadic identity just proved.
Integrating $|1+z+z^2|^2$ and retaining the five twist-one
constituents on its lower side shows that for all $t$
\[
 0\ge 9U-3-4p-2M_2
 +(5/2+3r)t+(-4-12r+54r^2)t^2,
 \qquad U=p^2+(p')^2+a^2+(a')^2+1/16 .
\]
The quadratic coefficient must be nonpositive, forcing $r=0$ or
$1/4$. At either value the maximum over $t$ of the last two terms
is at least $25/64$, attained at $r=0$ by $t=5/16$.
Since $U=109/360-5p/6+2p^2$, it follows that
\[
 37/320\le 27/(2D)+4/\sqrt D.
\]
The right side is smaller than $37/320$ for $D\ge1600$.

Finally the trace and norm of $D$ are $5m,12m$ for some positive
integer $m$, since $p+p'=5/12$. The discriminant identity
$25m^2-48m=5b^2$ with $b\in\mathbb Z$ forces $m=5k$ and
$b^2=k(125k-48)$. Because $D<1600$ and
$D'=1/(5/12-1/D)<3$ for $D>17$, one has $1\le k\le64$.
The only square is $k=12$, $b=132$. Then $(A,B,N)=(300,132,720)$;
the three dimension-vector identities give a unique unordered list of
doubled coordinates $(4,2),(6,2),(10,4),(12,6)$. Every dimension,
including the unit, has even doubled coordinates. But $D/4$ has
coordinates $(75,33)$ and is the dimension of the restriction of
an actual induction constituent. This contradiction excludes the
last mixed profile. Both finite tests are in
\texttt{quadratic\_profiles.py}.
\end{proof}



% ===== q5rational.tex =====
\begin{proposition}[The rational complementary profile]
\label{prop:quadratic-rational-profile}
Under the hypotheses of Proposition~\ref{prop:quadratic-profile-reduction},
if the complementary formal codegrees are rational, their multiset
is $\{4,5,5\}$. The two codegree-five characters are conjugate over
$\mathbb Q(\sqrt5)$, and the codegree-four character is rational.
\end{proposition}
\begin{proof}
We already have self-duality and $K=\mathbb Q(\sqrt5)$.
\par\smallskip\noindent\emph{Excluding codegree three.}\enspace
A rational row $\rho$ of codegree three
must pair locally at three with a totally ramified component.
If the other two rows form a quadratic pair of common codegree $h$,
the mass bound $2/3-2/h\le1/\sqrt5$ gives $h\le9$. If
$3\nmid h$, their joint idempotent is integral at three, and $K$ is
inert. If $3\mid h$, a real quadratic pair of rational codegree $h$
in a field of discriminant $c'$ satisfies $2h=c'\sum b_i^2$,
where $\chi_i-\chi'_i=b_i\sqrt{c'}$.
At least two $b_i$ are nonzero by Perron orthogonality.
No positive fundamental discriminant satisfies these requirements for
$h=3,6,9$. A cubic orbit of codegree three is excluded by the
projection-discriminant estimate.

Consequently all three rows are rational whenever a three occurs.
Its local partner $\tau$ satisfies $f_\tau\equiv6\pmod9$ since
$3+f_\tau$ is the square norm of their difference. The value six is
impossible: the rows could then differ only at a single position by
three. The resulting row of norm six has its nonunit coefficients
either all positive or all negative, in the latter case of
absolute values one and two, contradicting orthogonality to
the positive dimension row in both cases. Thus $f_\tau\ge15$.
Write $g$ for the remaining codegree. If $3\nmid g$, Lemma~\ref{lem:quadratic-inert-three} gives a
contradiction. If $3\mid g$, the mass bound and $f_\tau\ge15$ force $g=3$
or $6$ (and $f_\tau=15$ for $g=6$). In either case $g$ needs
its own totally ramified local partner, although only the inert
quadratic component is left. This removes all triples containing three.

\par\smallskip\noindent\emph{The remaining finite list.}\enspace
With the entries sorted, the same mass bound gives
the following enclosing ranges for the three entries:
\[
\begin{array}{c|cccc}
(r,s)&(4,4)&(4,5)&(4,6)&(5,5)\\
t&4\text{--}20&5\text{--}10&6\text{--}7&5\text{--}6 .
\end{array}
\]
A repeated four was excluded above. The unique unramified
valuation-one primes three and seven remove $(4,5,6),(4,5,7),(4,6,7)$.
The ten in $(4,5,10)$ has no local partner at two, since a rational
partner would have the same denominator valuation and $K$ is inert.
Likewise the six in $(5,5,6)$ has no partner at three. A quadratic
pair of sixes is impossible by the row difference identity, whereas
rational sixes must pair at three and cannot be congruent, as their
squared difference has norm twelve.
For $(4,5,8)$, the trace discriminant has odd valuation at two since
$\mathrm N(D)$ has even valuation at the inert prime. This contradicts
the discriminant-index identity for $K\times\mathbb Q^3$.
A rational row of codegree nine has four nonunit entries whose
squared norm is eight, so all are even. Evaluating $XX^*$ modulo
two at such a row contradicts its unit coefficient one. This removes
$(4,5,9)$.

\par\smallskip\noindent\emph{Excluding three codegrees equal to five.}\enspace
For $(5,5,5)$ there cannot be two rational rows, since all five
entries on each would be $\pm1$. For a quadratic pair,
$c'=5$ and its entries are $(a_i+b_i\sqrt5)/2$ with
$\sum b_i^2=2$, $\sum a_i^2=10$, $\sum a_ib_i=0$.
The two nonzero $b_i$ have opposite signs by positivity, hence
their two $a_i$ are equal and are $\pm1$. Of the remaining
nonunit entries, one is $\pm1$, the other zero. A rational
codegree-five row has equal signs on the two positions supporting
$b_i$, and its scalar product with the quadratic row is a sum of
three signed units, a contradiction. If all three rows comprise a
cyclic cubic orbit, the discriminant of its projected order is at
most $5^3(3/5)=75$, so this order is maximal in the conductor-seven
field. At five its completion is the unramified cubic integer ring
$A$. The intersection of the based order with this axis is the
dual of $A$ under $\mathrm{Tr}(xy/5)$, namely $5A$. Its other
projected order, of rank two, must therefore surject onto
$A/5A=\mathbb F_{125}$, which is impossible.

The only residual profile is $(4,5,5)$. Two rational codegree-five
rows cannot be orthogonal, as observed above, so these two rows form
a quadratic pair. The difference identity forces its field to be
$\mathbb Q(\sqrt5)$.
\end{proof}



% ===== q455_line.tex =====
\begin{proposition}\label{prop:quadratic-455-line}
Let $\mathcal C$ be a perfect self-dual rank-five fusion category
with all four nonunit dimensions greater than two and complementary
profile $4,5,5$ as in Proposition~\ref{prop:quadratic-rational-profile}. Then the
Frobenius--Perron row is, up to permutation,
$(1,x,x,x+e,2x-e)$ with $e\in\{1,-1\}$. If $y=ex$ and
$J=\operatorname{Tr}(y)$, then
\[
 7N(y)=J-3,\qquad J\equiv66\pmod{140},\qquad
 7J^2-4J+12=35b^2\quad (b\in\mathbb Z).
\]
In the center induction block, for each of the two labels $A$ of
codegree five, the multiplicities of the two dimension-$x$ simples
in $F(A)$ are consecutive integers.
\end{proposition}
\begin{proof}
In a quadratic row write each of the four nonunit entries as
$(t_i+s_i\sqrt5)/2$. The row, its conjugate and their orthogonality
give $\sum s_i^2=2$, $\sum t_i^2=6$ and $\sum s_i t_i=0$.
Two entries are signed versions of $\varphi=(1+\sqrt5)/2$ and
$\varphi'=1-\varphi$, one is $\pm1$, and one vanishes. Orthogonality
of the positive row to both characters forces the first two dimensions
to coincide, and the signs are equal. Orthogonality to the rational
codegree-four row then determines its entries up to the sign of the last
one. Equivalently, after changing signs on individual basis elements,
the five columns are presented by the three character rows
\[
 (1,y,y,-1-y,2y-1),\quad (1,\varphi,\varphi',1,0),\quad
 (1,-1,-1,0,1).
\]
When returning to positive dimensions this gives the claimed row,
since $|y|>2$. Orthogonality of the first row to its conjugate gives
$3+7N(y)-J=0$. The dimension order is $\mathbb Z[y]$ and
$D=3+7y^2-2y=(6-J)+(7J-2)y$.
Both $D/4$ and $D/5$ belong to the dimension order by central
induction, hence so does $D/20$. Thus $J\equiv6\pmod{20}$;
combining with $J\equiv3\pmod7$ gives $J\equiv66\pmod{140}$.
The quadratic discriminant $J^2-4N(y)$ is five times a square.

For the final statement let $a=5S_{AA}$, an algebraic integer in $K$;
indeed it equals the complementary character evaluated on $F(A)$.
Its conjugate $a'$ is $5S_{A'A}$. The entry joining the rational
codegree-four label to $A$ vanishes by coprimality: multiplication
by $4$ and by $5$ both give algebraic integers and each conjugate of
the entry has modulus less than one, using modular row orthogonality.
The unit coefficient of $F(A)$ therefore gives $a+a'=3$.
Writing $a=(3+j\sqrt5)/2$ gives $j$ odd. The squared norm of the
column, restricted to the unit induction constituents, gives
$2+a^2+(a')^2\le25$, hence $|j|=1$.
Weighted column orthogonality recovers the signed multiplicities in
$F(A)$ at the first two nonunit coordinates: they are
$(J+(3+5j)/2)/5$ and $(J+(3-5j)/2)/5$.
Their difference is $j$. The original two signs agree, so the
unsigned multiplicities are consecutive as well.
\end{proof}




% ===== q455_close.tex =====
\begin{proposition}\label{prop:quadratic-455-end}
The quadratic profile with complementary codegrees $4,5,5$ is impossible
for a perfect self-dual rank-five fusion category with all four nonunit
dimensions greater than two.
\end{proposition}
\begin{proof}
Write $R,A,A'$ for the induction constituents of codegrees $4,5,5$.
Put $w_B=S_{0B}$, $r_B=S_{RB}$ and $h_B=S_{AB}+S_{A'B}$.
These vectors are real, and both $r,h$ have rational entries.
Use the order $0,0',R,A,A'$ for the labels $I$ of unit induction;
put $p=1/D$, so $p'=7/20-p$. Proposition~\ref{prop:quadratic-455-line} gives
\[
 w|_I=(p,p',1/4,1/5,1/5),\quad
 r|_I=(1/4,1/4,1/2,0,0),\quad
 h|_I=(2/5,2/5,0,3/5,3/5).
\]
Lemma~\ref{lem:odd-row} gives two-primary twist support for $r$.
Lemma~\ref{lem:prime-local-block} gives five-primary support for
$A,A'$. To check its hypotheses at each prime $\ell\ne5$, their
summed primitive idempotent is rational and $\ell$-integral, since
both modular formal codegrees are $25$. Their restrictions to the
twist-one columns $A,A'$ are independent: the diagonal entries are
$(3+j\sqrt5)/10$ and the off-diagonal entries $(3-j\sqrt5)/10$,
where $j=\pm1$. In particular $r_B h_B$ vanishes away from twist one.
The sum of these products at twist one is zero by row orthogonality.

For a real vector $z$ write $[z]_k=\sum_B z_B\Re(\theta_B^k)$;
products of vectors below are coordinatewise. The modular factorization
in Lemma~\ref{lem:odd-row} shows, at the twist-one rows $0,R,A,A'$,
that for $5\nmid k$,
\[
 [wh]_k=2/5,\qquad [h^2]_k=6/5 .
\]
Indeed replace $k$ by a unit modulo $N$ with its original value on
the five-part and value one on the prime-to-five part; its signed
Galois permutation fixes the sum of the two label vectors $A,A'$,
with positive signs by their vacuum coordinate $1/5$.
Likewise $[r^2]_k=1/2$ for odd $k$. Averaging over the odd powers
on the two-primary twists gives mass at one minus mass at minus one
equal to $1/2$ for the probability weights $r_B^2$.
Thus $[r^2]_2\ge0$, since the cosine of the double argument is one
at both of these twists and is everywhere at least minus one.

For every $k$ we claim
\[
             2[wr]_k=1-[r^2]_k .
\]
The coefficient-of-unit character formula on a label $B$ outside $I$
is $\operatorname{Tr}_K(w_B)=-r_B-h_B$. Choose a cyclotomic
automorphism nontrivial on $K=\mathbb Q(\sqrt5)$ and trivial on the
two-part of $N$. Its signed column permutation preserves $I$ and
each individual twist on the support of $r$. Its signs cancel
in products with the rational row $r$. Therefore on the complement
of $I$ one may replace twice $w_B r_B$ by
$\operatorname{Tr}_K(w_B)r_B$ under the cosine sum. Using
$[rh]_k=0$ gives the claim because
$\sum_I (2wr+r^2+rh)=1$.

Set $v=w+(r+h)/2$ and $M_k=[w^2]_k$. The contributions just found
give
\[
 [v^2]_0=7/4,\quad [v^2]_1=43/40+p,\quad
 [v^2]_2\le6/5+M_2,\quad [v^2]_3=43/40+M_3.
\]
At twist one the $v^2$ mass is at least
$(p+13/40)^2+(27/40-p)^2+3/4$ on $I$ alone.
The integral of $|7+12z+12z^2+7z^3|^2$ is
$386[v^2]_0+624[v^2]_1+336[v^2]_2+98[v^2]_3$.
Dropping a positive $p^2$ term in the lower bound yields
\[
 \frac{119}{200}\le \frac{8174}{5D}+336M_2+
                  98\left(M_3-\frac{19}{49}\right).
 \tag{*}
\]
The arithmetic of the dimension row forces the twist exponent to be
prime to three. Indeed write $T=(7J-2)/20\in\mathbb Z$ by
Proposition~\ref{prop:quadratic-455-line}. Completing the square gives
$49b^2=16(5T^2+1)$. Thus $4\mid b$, and $U=7b/4$ satisfies
\[
 U^2-5T^2=1,\qquad 7\mid U.
\]
Every positive Pell solution satisfies
$|U|+|T|\sqrt5=(9+4\sqrt5)^n$ for $n\ge1$:
the least unit greater than one in this norm-one group is
$9+4\sqrt5$, as is seen by trying $|T|=1,2,3,4$.
Division by its powers gives the assertion.
Put $\epsilon=9+4\sqrt5$. Modulo seven, $\epsilon^2=2\sqrt5$
and $\epsilon^4=-1$. The rational coefficients of
$1,\epsilon,\epsilon^2,\epsilon^3$ reduce to $1,2,0,5$.
Consequently $n\equiv2\pmod4$. Modulo three, $\epsilon^2=-1$, so
$3\nmid U$ and $3\nmid b$. By the formula for $D$ in the preceding
proposition,
\[
 \operatorname{Tr}(D)=7J^2-4J+12=35b^2,\qquad
 \operatorname{N}(D)=(20/7)\operatorname{Tr}(D)=100b^2 .
\]
Categorical Cauchy now makes three prime to the twist exponent.
The modular factorization at the vacuum yields
$M_3\le\max(p,p')<7/20$, since the signed Galois permutation sends
the vacuum row to its dimension orbit.

Finally $n\ge2$ implies $|T|\ge72$, whence $|J|\ge205$.
The bound $|J|\le2\sqrt D$ gives $D>60$, so
$D'=1/(7/20-1/D)<3$. Since $|y'|<\sqrt3$,
$D\ge y^2\ge(|J|-\sqrt3)^2>40000$. In (*) use
$M_2\le(1+2\sqrt{D-1})/D$. Its right side is then at most
\[
 1971/40000+672/200+98(7/20-19/49)
 = -11629/40000 .
\]
This contradiction excludes the entire profile.
\end{proof}




% ===== ind.tex =====
\section{Perfect categories in cubic degree}
\label{sec:perfect-cubic}
Assume $K=\mathbb Q(D)$ is cubic. It is cyclic; write $F$ for its
conductor. The two characters outside the dimension orbit have
rational integral codegrees $r\le t$: each belongs to $K$ by central
induction and to a character field of degree at most two.
The small-codegree restrictions are those of
Lemma~\ref{lem:small-complement}. Our aim is the dimension bound
$D<36$, which will suffice for the final finite classification.

\subsection{Local restrictions and a conductor bound}
Suppose that $3$ is inert in $K$ and $r=3$ belongs to a rational
character $\rho$. By Lemma~\ref{lem:r5-local-partner}, its partner must
be another rational character $\tau$, whose row is congruent to
$\rho$ modulo three. Thus $v=(\tau-\rho)/3$ is integral,
$v_0=0$, and $v$ is perpendicular to the strictly positive dimension
row. It is nonzero, so has squared norm at least two. Orthogonality
of the two rows gives
\[
 t=9\|v\|^2-3\ge15,\qquad t\equiv6\pmod9.                  \tag{*}
\]
\begin{lemma}\label{lem:r5-odd-partner}
Suppose the complementary rows $\rho,\tau$ are rational with respective
codegrees $r,t$. An odd prime $\ell$ satisfying
$v_\ell(t)=1$ and $\ell\nmid r$ ramifies in $K$.
\end{lemma}
\begin{proof}
Let $A_t$ be the induction constituent for $\tau$. Put
$x=\rho(F(A_t))/r$ and $y=\tau(F(A_t))/t$, where here $F$ denotes
the forgetful functor. Its Frobenius--Perron dimension is $D/t$.
Expanding the unit coefficient gives $1=3/t+x+y$. For each simple
$X$, the same character expansion gives
\[
 [X,F(A_t)]=\mathrm{Tr}_K(d_X)/t+
                   x\rho(X^*)+(1-3/t-x)\tau(X^*).
\]
The two rational rows are fixed by duality and $x$ is integral at
$\ell$ because $\ell\nmid r$. Hence
$\mathrm{Tr}_K(d_X)\equiv3\tau(X)\pmod\ell$.

If $\ell$ were unramified, it would be inert or completely split in
$K$. Lemma~\ref{lem:r5-local-partner} gives $\tau$ one totally
ramified partner. The character $\rho$ has integral primitive idempotent
at $\ell$, so the partner must be a cubic-field branch. This excludes
the inert case. In the split case denote the partner by $K_1$; its
residue character is $\tau\bmod\ell$. The two remaining cubic
branches $K_2,K_3$ must lie in a single local factor. Indeed, if
either were isolated, that factor would contain its primitive
idempotent, whose unit coefficient $1/\sigma D$ is not
$\ell$-integral since $D/t$ and its conjugates are algebraic integers.
The branches $K_2,K_3$ thus give the same residue character $u$.
Reducing the field trace now gives
$\tau(X)+2u(X)=3\tau(X)$ in $\mathbb F_\ell$ for all $X$.
As $\ell$ is odd, $u=\tau$, contradicting the fact that the two
local factors correspond to distinct residual characters.
\end{proof}

\begin{proposition}\label{prop:cubic-no-three}
In a perfect cubic category, if one complementary codegree is three, the other equals six.
\end{proposition}
\begin{proof}
Suppose $r=3$. The small repeated-codegree observations above imply that both complementary rows are rational. Use the positive spherical realization and put $J=A_\rho$, $H=A_\tau$ for the two complementary induction constituents of codegrees $3,t$. Use the unitary modular representation of the center with normalized $S$ and ordinary twist $T$. Its anomaly is one. For a central simple $W$ set $a_W=S_{0W}=\dim(W)/D$ and $b_W=S_{JW}=\rho(FW)/3$. Thus $b$ is rational and $3b_W$ is integral. On each of the three vacuum-orbit constituents of $I(1)$ one has $b=1/3$ by Galois symmetry. Moreover $\sum_{W\subset I(1)} b_W=1$: by additivity the sum is $\rho(FI(1))/3=1$ since the Casimir acts in $\rho$ by its codegree. All multiplicities in $I(1)$ are one. Consequently, with $m=3b_H\in\mathbb Z$, we have $b_H=m/3$, $b_J=-m/3$.

The weights $b_W^2$ sum to one. Their first twist moment is $b_J=-m/3$, by $STS=T^{-1}ST^{-1}$ at $(J,J)$, and their second twist moment equals the same value by Lemma~\ref{lem:odd-row}. Indeed the central modular character indexed by $J$ is rational and its formal codegree is $S_{0J}^{-2}=9$. The five labels of $I(1)$ all have twist one. Integrating $|1+z+z^2|^2$ with these weights gives
\[
 3+2m^2=9\sum_{W\subset I(1)}b_W^2\le 3-2m.
\]
Hence $m=0$ or $-1$, with equality in both cases. Each label outside $I(1)$ for which $b_W\ne0$ therefore has primitive cube-root twist. Orthogonality of the vacuum row to $b$ gives $\sum_W a_W b_W=0$. Taking the real part of $STS$ at $(0,J)$ consequently yields
\[
 \frac13=\sum_W a_Wb_W\operatorname{Re}(\theta_W)
       =\frac32\sum_{W\subset I(1)}a_Wb_W .
\]
On the three vacuum-orbit labels the $a_W$ add to $s=1-1/3-1/t$, and $a_J=1/3$, $a_H=1/t$. Substitution gives
\[
 \frac29=\frac{s}{3}-\frac{m}{9}+\frac{m}{3t},
 \qquad mt=3(m-1).
\]
The case $m=0$ is impossible. The case $m=-1$ forces $t=6$.
\end{proof}

\begin{proposition}\label{prop:r5-cubic-frontier}
Order the complementary codegrees as $r\le t$ and set $s=1-1/r-1/t$.
Then $s\ge1/2$ and $s\le 3/D+2/\sqrt F$. If $F=13$, then $s\ge7/12$.
\end{proposition}
\begin{proof}
Set $p_\sigma=1/\sigma(D)$, $s=1-1/r-1/t$, and let $\eta$ be a nontrivial
cubic character cutting out $K$. Proposition~\ref{prop:ordinary-fourier}
gives $|\sum_\sigma\eta(\sigma)p_\sigma|\le F^{-1/2}$.
Cubic Fourier inversion therefore gives
\[
                     3/D\ge s-2/\sqrt F .
\]
Proposition~\ref{prop:cubic-no-three} gives $s\ge1/2$ when $r=3$; otherwise both codegrees are at least four. At conductor thirteen, three is inert since $3^4\equiv3\pmod{13}$, so $(*)$ excludes $(3,6)$. The observations above exclude $(4,4)$, and Lemma~\ref{lem:r5-odd-partner} excludes $(4,5)$ since five is unramified. The least remaining mass is $1-1/4-1/6=7/12$. The cyclic cubic conductors below nineteen are $7,9,13$, directly from the even Dirichlet characters of order three.
\end{proof}

\begin{corollary}\label{cor:r5-cubic-even}
Suppose $F\in\{7,9\}$ and $r,t$ are odd. Then $s\ge3/5$ and
\[
              (s-1/\sqrt F)D\le3+3\sqrt{D-1}.
\]
\end{corollary}
\begin{proof}
Assume $r,t$ are odd. The idempotent onto the complementary
rational algebra is integral at two. This follows from the
primitive idempotent formula for rational characters; for a
quadratic pair, use the sum of the two conjugate idempotents,
whose integral numerators have integral field trace. Thus the
projected cubic order $B$ is an orthogonal direct summand over
$\mathbb Z_2$, self-dual for $\mathrm{Tr}(xy/D)$. The prime two
is inert in the fields of conductors seven and nine.
If two divided $\mathrm N(D)$, then $B$ would differ from its
normalization, since the unramified trace dual of that normalization
is itself. By Nakayama its image in the residue field
$\mathbb F_8$ is a proper unital subring, hence $\mathbb F_2$.
The elements $U=D/D'$, $V=D/D''$ lie in $B$ by central induction
and are algebraic units by the formal codegree theorem. Thus both
reduce to one. But $1+U+V=sD$ reduces to zero on the right, since
$s=1-1/r-1/t$ is integral at two. This contradiction proves oddness
of the norm.

Categorical Cauchy and the higher Gauss sum identity for the ordinary
center therefore give
\[
 \frac1{\sigma D}
  =\frac{1}{D}\left|\sum_i d_i\nu_2(X_i)\right|
  \le \frac{1+2\sqrt{D-1}}{D},
\]
where $\sigma$ acts by inverse squaring on the cyclotomic roots
of the odd exponent. Its restriction to either cubic field is nontrivial.
By cubic Fourier inversion the remaining reciprocal is at most $(s+2/\sqrt F)/3$. Summing this bound, $1/D$ and the displayed second-indicator bound gives the claimed inequality. Proposition~\ref{prop:cubic-no-three} rules out a codegree three paired with an odd codegree. Thus $r,t\ge5$.
\end{proof}



% ===== cubic_two.tex =====
\begin{theorem}\label{thm:cubic-at-two}
Suppose $F\in\{7,9\}$ and at least one of $r,t$ is even. Then both
complementary characters are rational. After interchanging their
codegrees one has
\[
 a=v_2(r)<b=v_2(t),\qquad a\equiv0\pmod2,\qquad
 b\ge 2,\qquad v_2(D)=b.
\]
Here the valuation on the cubic field $K$ is normalized by $v_2(2)=1$.
\end{theorem}
\begin{proof}
Central induction implies
$v_2(D)\ge\max(v_2(r),v_2(t))>0$. The prime two is inert in $K$.
Write $B$ for the image of $R\otimes\mathbb Z_2$ in $K\otimes\mathbb Q_2$,
and $O$ for the integers in this unramified cubic field.

Suppose $B=O$. The local factor of $R\otimes\mathbb Z_2$ mapping onto $O$
has residue field $\mathbb F_8$. None of the complementary characteristic-zero
branches can belong to that factor: its residue field injects into the
residue field of the normalization of each branch, and every
complementary branch has degree at most two. The factor is therefore
$O$ itself. The coefficient-of-unit pairing restricts to a perfect
pairing on each local factor. On $O$ it is $\operatorname{Tr}(xy/D)$,
whose dual lattice is $DO$ by unramifiedness. This contradicts $v_2(D)>0$.
Hence $B$ is proper. Its image in $O/2O=\mathbb F_8$ is the subfield
$\mathbb F_2$, since a full image would give $B=O$ by Nakayama.
The central dimensions $D/D'$ and $D/D''$ belong to $B$ and are units
in $O$ by the $d$-number property. Their residues are consequently one.
Therefore
\[
  (1-1/r-1/t)D=1+D/D'+D/D''\equiv1\pmod{2O}.
\]
If $v_2(r)=v_2(t)=b>0$ the first factor on the left has valuation at
least $1-b$, a contradiction. Otherwise its valuation equals minus
the larger valuation $b$, proving $v_2(D)=b$.

A conjugate pair of complementary characters has equal rational
codegrees, so both characters must be rational. In particular every
basic element is self-dual. The based-order discriminant is the product
$r t\,N(D)$ of the codegrees. It is a square since the rational
algebra is $\mathbb Q\times\mathbb Q\times K$ with $K$ cyclic cubic.
Its valuation at two equals $a+b+3b=a+4b$, so $a$ is even. Finally,
if $b=1$ the other rational character has odd codegree and its
primitive idempotent is integral at two. It cannot supply the partner
for the character of valuation one in Lemma~\ref{lem:r5-local-partner},
and the inert cubic branch cannot supply a totally ramified partner
either. This proves $b\ge2$.
\end{proof}



% ===== cubic_three.tex =====
\begin{proposition}\label{prop:cubic-four-odd}
Under the hypotheses of Theorem~\ref{thm:cubic-at-two}, if one complementary codegree is four then the other is even.
\end{proposition}
\begin{proof}
Write $r=4$, $t$ odd, and denote the rational rows by $\rho,\tau$. Each $\rho_i$ is $\pm1$ except for a single zero, say at $Z$. At two the primitive idempotent of $\tau$ is integral. By Theorem~\ref{thm:cubic-at-two} the cubic order has residue field $\mathbb F_2$ and $v_2(D)=2$. The rational branch of $\rho$ cannot be isolated, since its idempotent has nonintegral unit coefficient $1/4$. Thus it shares a residue character with the cubic branch: $d_i=\rho_i+2x_i$, $x_i\in O=\mathcal O_K\otimes\mathbb Z_2$. The identity $(1-1/4-1/t)D=1+D/D'+D/D''\equiv1\pmod{2O}$ in the proof of Theorem~\ref{thm:cubic-at-two} gives $D/4\equiv1$. Orthogonality $\sum_i\rho_i d_i=0$ gives $\sum_{i\ne Z} x_i\equiv0$ modulo two. Substituting the $d_i$ into $D=\sum_i d_i^2$ yields $D/4\equiv1+x_Z^2$, so $d_Z\in4O$. Weighted column orthogonality now gives
\[
 1=\mathrm{Tr}(d_Z^2/D)+\tau_Z^2/t,\qquad
 0=\mathrm{Tr}(d_Zd_i/D)+\tau_Z\tau_i/t\quad(i\ne Z).
\]
Reducing modulo two shows $\tau_Z$ is odd and all other $\tau_i$ have the same residue $\mathrm{Tr}(d_Z/4)$. The unit entry is one, so all five $\tau_i$ are odd and $t\equiv5\pmod8$.

For $J=A_\rho$ and $H=A_\tau$, symmetry gives $S_{JH}\in\tfrac14\mathbb Z\cap\tfrac1t\mathbb Z=\mathbb Z$. The vacuum entry of the $J$ row is nonzero, so $|S_{JH}|<1$ and $S_{JH}=0$. The $H$ row therefore has $S_{HH}=1-3/t$ by summing its entries on $I(1)$ as in Proposition~\ref{prop:cubic-no-three}. The dimension bound on modular character values gives $|S_{HH}|\le D/t^2$, whence $D\ge t(t-3)$.

Write $U=\sum_{B\subset I(1)}(\dim B/D)^2$. The first twist moment of the vacuum-row squares is $1/D$ and the second is bounded by $(1+2\sqrt{D-1})/D$, by the central second-indicator formula. Thus integration of $|1+z+z^2|^2$ gives $U\le1/3+4/(9D)+2(1+2\sqrt{D})/(9D)$. Now $t\ge125$ would imply $D\ge15250$. Cauchy's inequality on the two nonvacuum cubic entries gives the contradiction
\[
 U\ge\frac1{16}+\frac12\left(\frac34-\frac1{125}-\frac1{15250}\right)^2
  >\frac13+\frac{4+2(1+2\cdot124)}{9\cdot15250}\ge U.
\]
Consequently $t=5,13,21,29,37,45,53,61,69,77,85,93,101,109$ or $117$. In this list each integer has a prime of valuation one outside $\{3,7\}$, except $21$, which has both three and seven to valuation one. All such primes must ramify in $K$ by Lemma~\ref{lem:r5-odd-partner}, whereas the conductor $7$ or $9$ has only one prime divisor. This excludes every entry.
\end{proof}



% ===== cubic_profiles.tex =====
\begin{lemma}\label{lem:cubic-eight-depth}
In the rational-complement case of Theorem~\ref{thm:cubic-at-two}
an odd codegree and a codegree of two-adic valuation three cannot occur together.
\end{lemma}
\begin{proof}
Let $r$ be odd, $v_2(t)=3$, $O$ the cubic integer ring at two, $B$ the cubic image, and $\tau$ the row of codegree $t$. The $r$-idempotent splits off its rational factor. In the
remaining order $A$, the intersection with the rational axis is
$8\mathbb Z_2$, by the primitive idempotent formula. Thus $A$ is
the fiber product of $\mathbb Z_2$ and $B$ over $\mathbb Z/8$ and
$\psi(d_i)=\tau_i$ modulo eight. The order discriminant has
valuation $v_2(t)+3v_2(D)=12$. Since the cubic extension is unramified,
$[\mathbb Z_2\times O:A]=2^6$. Dividing by
$[\mathbb Z_2\times B:A]=8$ gives $[O:B]=8$.

On $O/\mathbb Z_2$ the elementary divisors of $B/\mathbb Z_2$ are
therefore $2,4$. Choose $2u\in B$ such that $\bar u$ is outside
$\mathbb F_2$. Since $1,\bar u,\bar u^2$ are independent in
$\mathbb F_8$, the span of $1,2u,4u^2$ has index eight in $O$,
and is all of $B$. Write $2s$ for the image of $2u$ under
$\psi$, using an integer representative
(it is even by reduction modulo the maximal ideal). For each basic
element put $e_i=(d_i-\tau_i)/2$. The gluing condition, expanded in
the displayed basis of $B$, gives
\[
 \bar e_i\in\{0,\bar u-\bar s\}\subset\mathbb F_8 .
\]
These residues cannot all vanish: otherwise all $d_i$ would belong to
$\mathbb Z_2+4O$, of index $16$, whereas they span $B$.

Orthogonality and the norm of the rational row give
\[
 \sum_i\tau_i e_i=-t/2,\qquad
 \sum_i e_i^2=(D+t)/4.
\]
Here $D/8$ reduces to one, since it is an odd rational unit times
the central dimension $D/t\in B$. As $t/8$ is odd, the second sum
belongs to $4O$. Elements equal modulo $2O$ have equal squares modulo
$4O$. The number of nonzero residues $\bar e_i$ must thus be divisible
by four. The unit contributes $e_0=0$, so all four other residues
equal the same nonzero element. The identity $\sum_i\tau_i^2=t$ with
$4\mid t$ forces exactly four odd $\tau_i$, including the unit.
Reducing the first displayed sum modulo $2O$ now gives the contradiction
$3(\bar u-\bar s)=0$.
\end{proof}

\begin{proposition}\label{prop:cubic-finite-profiles}
Assume $F\in\{7,9\}$ and at least one complementary codegree is even. Then either $r=4$, or $r\ge7$ and $s\ge65/84$.
\end{proposition}
\begin{proof}
Let $B_r,B_t$ be the central induction constituents of the two rational
characters and $w$ the value of the first character on the restriction
of $B_t$. In the normalized center matrix
$S_{B_r B_t}=w/r$; the three entries in column $B_t$ from the
dimension orbit are $1/t$. The unit-coefficient expansion of its
restriction gives $S_{B_t B_t}=1-3/t-w/r$. Since $w$ is an integer
and the row of $B_r$ has a nonzero vacuum entry, its unit norm gives
$w\le r-1$. Spectral domination of a modular character gives
$|S_{B_t B_t}|\le D/t^2$. Hence
\[
 D\ge D_0:=\max(t,t^2/r-3t).
\]
The sum of the squared reciprocal codegrees is at least
$r^{-2}+(1-1/r-1/t-1/D)^2/2$, by Cauchy's inequality on the two real
characters outside the positive character and the rational pair. The
degree-two Fej\'er inequality for central twists bounds this sum above
by $1/3+2/(3D)+4/(9\sqrt D)$: the first normalized moment is $1/D$,
and the second has modulus at most $(1+2\sqrt{D-1})/D$.
For $r=5$ and $t\ge76$ the preceding inequality fails: insert $D_0=76^2/5-3\cdot76$ in its two sides at $t=76$. The lower bound then increases with $t$ and the upper bound decreases. Thus $r=5$ implies $t\le75$.

A codegree three forces the other to be six by Proposition~\ref{prop:cubic-no-three}, contradicting the valuation conditions of Theorem~\ref{thm:cubic-at-two}.
The local theorem excludes $r=6$: a valuation one cannot occur, as the
smaller valuation must be even and the larger at least two.
If $r=5$, five is inert in both possible cubic fields. The
local-partner lemma forces the two rational characters to be congruent
modulo five. Their difference divided by five is a nonzero integral
vector perpendicular to the positive dimension row, of squared norm at
least two. Thus $5+t\equiv0\pmod{25}$ and $t\ge45$. On the other hand
$4\mid t$, and the only positive multiple of four at most $75$
congruent to $20$ modulo $25$ is $20$. This excludes $r=5$.
For $r=7$ the other codegree is a multiple of four and cannot be eight by Lemma~\ref{lem:cubic-eight-depth}. Thus $t\ge12$ and $s\ge65/84$. For $r=8$ the valuation conditions and Lemma~\ref{lem:cubic-eight-depth} similarly give $t\ge12$. For $r\ge9$ one already has $s\ge7/9>65/84$.
\end{proof}



% ===== cubic_four_depth.tex =====
\begin{proposition}\label{prop:cubic-four-depth}
Suppose the perfect cubic category has conductor $F\in\{7,9\}$ and
complementary codegrees $4,t$ with $t$ even. Then $v_2(t)=4$. If $Z$ is the simple where the codegree-four row vanishes, then $d_Z\equiv \tau_Z\pmod{8O}$ and $\tau_Z\equiv2\pmod4$, where $O$ is the ring of integers of the cubic completion and $\tau$ is the row of codegree $t$.
\end{proposition}
\begin{proof}
Write $b=v_2(t)>2$ by Theorem~\ref{thm:cubic-at-two}, and denote the codegree-four row by $\rho$. The two rational branches and the inert cubic branch lie in one local factor. Indeed the separate rational idempotents have nonintegral unit coefficients, and the cubic idempotent has coefficient $1-1/4-1/t$ of valuation $-b$. The cubic image has residue field $\mathbb F_2$, so $d_i,\rho_i,\tau_i$ have common residues modulo two. There are four odd entries in $\tau$, in the support of $\rho$. Since $8\mid t$, their squares show $\tau_Z\equiv2\pmod4$.

The image $C$ on the two rational axes is $\mathbb Z_2\times_{\mathbb F_2}\mathbb Z_2$, since the differences $\tau_i-\rho_i$ are even and one of them has valuation one (at $Z$). Its intersection with the order supported on those axes is its trace-dual lattice for $xx'/4+yy'/t$, by perfectness of the full pairing. The index of this lattice in $C$ is $2^b$: in $(1,1),(0,2)$ coordinates the Gram determinant is $1/t$. Thus the index of the full order in $B\times C$ is $2^b$, where $B$ is its cubic projection. The full index in $O\times\mathbb Z_2^2$ has valuation $(2+b+3b)/2=1+2b$ by the discriminant formula and unramifiedness. As $C$ itself has index two, $[O:B]=2^b$.

Let $j\ge1$, and suppose $d_i\equiv\tau_i\pmod{2^jO}$. Put $e_i=(d_i-\tau_i)/2^j$; in particular $e_0=0$. The norm of the rational row and its orthogonality to $d$ give
\[
 \sum\tau_i e_i=-t/2^j,\qquad\sum e_i^2=(D+t)/2^{2j}.
\]
If $b\ge2j+1$, these sums vanish modulo two and four, respectively, since $D/t\equiv1\pmod{2O}$. Reducing both modulo two first gives $\bar e_Z=0$ and residues $a,c,a+c$ at the other three nonunit entries. Squares modulo four depend only on residues modulo two. Therefore $a^2+ac+c^2=0$ in $\mathbb F_8$, which lacks a nontrivial cube root of unity. All $\bar e_i$ vanish and the congruence lifts to $j+1$. Induction puts $B\subset\mathbb Z_2+2^{\lceil b/2\rceil}O$. Its computed index forces $b=2k$ and $B=\mathbb Z_2+2^kO$.

Now put $q=2^k$, $t=q^2m$ with $m$ odd, and $e_i=(d_i-\tau_i)/q$. The same two sums vanish modulo two, proving again that $e_Z=2z$ for some $z\in O$. Write $\tau_Z=2v$, $v$ odd, and $\lambda=t/D$, a local unit with residue one. The unit and $Z$ columns of the character table give
\[
 \mathrm{Tr}(\lambda)=3t/4-1,\quad
 \mathrm{Tr}(\lambda d_Z)=-\tau_Z,\quad
 \mathrm{Tr}(\lambda d_Z^2)=t-\tau_Z^2.
\]
Substitute $d_Z=2v+2qz$ to get
\[
 \mathrm{Tr}(\lambda z)=-3qmv/4,\qquad
 \mathrm{Tr}(\lambda z^2)=m(1+3v^2)/4 .
\]
The latter is odd, since $v^2\equiv1\pmod8$. Their residues coincide because $\lambda\equiv1$ and finite-field trace is invariant under squaring. Hence $4\mid q$ but $8\nmid q$, giving $k=2$ and $b=4$.
\end{proof}

\begin{corollary}\label{cor:cubic-profile}
Suppose $F\in\{7,9\}$, $r=4$, and at least one complementary codegree is even. At $F=9$ one has $s\ge11/16$. At $F=7$ the category has field $\mathbb Q(\zeta_7)^+$ and complementary rational codegrees $4,16$. In this latter case two nonunit simples have dimensions $d_Z=2d_X$.
\end{corollary}
\begin{proof}
Proposition~\ref{prop:cubic-four-odd} and Proposition~\ref{prop:cubic-four-depth} give $v_2(t)=4$, hence $s\ge11/16$. For $F=7$ suppose $t\ge112$. Then $D\ge t^2/4-3t\ge2800$ as in Proposition~\ref{prop:cubic-finite-profiles}. The reciprocal sum is $s\ge83/112$. Fourier inversion bounds each reciprocal by $(s+2/\sqrt7)/3<503/1000$. The two other than $1/D$ are therefore each at least $83/112-1/2800-503/1000>235/1000$. Since $D/t$ is an algebraic integer, its norm is at least one, giving $D>112^3(235/1000)^2>77000$. Now the twist-one mass is at least $1/16+(83/112-1/77000)^2/2$, which exceeds its upper bound $1/3+2/(3\cdot77000)+4/(9\sqrt{77000})$ from Proposition~\ref{prop:cubic-finite-profiles}. Thus $t<112$, leaving $16,48,80$. In conductor seven, Lemma~\ref{lem:r5-odd-partner} excludes $48,80$ at three and five.

For $t=16$, let $e_i=(d_i-\tau_i)/4$. We have $\bar e_0=\bar e_Z=0$, and residues $a,c,a+c$ at the three remaining nonunit entries. They are independent apart from their sum (as elements of the additive group): otherwise the dimensions would fail to span $B=\mathbb Z_2+4O$, since their reductions after subtracting the integers and dividing by four would span at most one dimension modulo $\mathbb F_2$. The identities $\sum\rho_i e_i=0$ and $\sum\tau_i e_i=-4$ imply $\sum_i((\tau_i-\rho_i)/2)e_i=-2$. Thus $(\tau_i-\rho_i)/2$ have equal parity on the three nonzero residues. If they were odd, $\sum\rho_i\tau_i\equiv1-3\not\equiv0\pmod4$. Hence $\tau_i\equiv\rho_i\pmod4$ on the support of $\rho$. Now $\tau_Z=\pm2$, and the four odd squares in $16-\tau_Z^2=12$ are $1,1,1,9$. Therefore $\tau_X=-3\rho_X$ at exactly one nonunit $X$, and $\tau_i=\rho_i$ at the other three support entries. Subtracting the two rows against $d$ gives $\tau_Z d_Z=4\rho_Xd_X$, proving $d_Z=2d_X$.
\end{proof}



% ===== cubic_polar.tex =====
\begin{proposition}\label{prop:cubic-polar}
The conductor-seven profile with complementary codegrees four and sixteen in Corollary~\ref{cor:cubic-profile} is impossible.
\end{proposition}
\begin{proof}
Write $R,T$ for the induction constituents of codegrees $4,16$, and
$\rho,\tau$ for their characters, with the signs and notation of
Corollary~\ref{cor:cubic-profile}. Thus $d_Z=2d_X$,
$\rho_Z=0$, $\tau_Z=2\rho_X$, and $\tau_X=-3\rho_X$,
where $\rho_X=\pm1$. Write $I$ for the five constituents of $I(1)$,
in the order of the three dimension characters followed by $R,T$.
The three conjugates of $p=1/D$ have sum $s=11/16$.
Put $b=4S_{RT}\in\mathbb Z$. Each row in $I$ has sum one on $I$,
by applying its character to the Casimir $F(I(1))$.
Therefore $S_{RR}=(1-b)/4$ and $S_{TT}=(13-4b)/16$.
The real row $4S_{R,-}$ has squared norm sixteen and squared mass
$c=3+(1-b)^2+b^2$ on $I$, so $c\le16$ and
$2c-16\le4(1-b)$ by its first twist moment. The integer solutions
are $b=-1,0,1$; $b=-1$ makes $S_{TT}>1$.
Taking the based inner product of $Z-2X$ with $F(R)$ gives the integer
$\rho_X(5b-1)/2$, by the formal-codegree character formula: the
dimension orbit contributes zero, and the other two terms use
$\rho(FR)=1-b$, $\tau(FR)=4b$. Hence $b=1$.

By Lemma~\ref{lem:odd-row}, the rational modular rows
$r_B=S_{RB}$ and $t_B=S_{TB}$, of formal codegrees $4^2,16^2$,
are supported on twists of two-power order.
For every odd $k$, choose a unit modulo $N$ equal to $k$ modulo its
two-part and one modulo its odd part. The modular factorization in
Lemma~\ref{lem:odd-row} and rationality of the two rows show, for
$h=r+t$, that
\[
 \sum_B \theta_B^k h_B^2=h_R+h_T=17/16.
\]
Averaging over odd $k$ on the two-part shows that its twist-one mass
minus its twist-minus-one mass is $17/16$.
On $I$ the entries of $h$ are $(5/16,5/16,5/16,1/4,13/16)$, of
squared norm $65/64$. Thus its twist-one mass outside $I$ is at
least $3/64$.

Write $w_B=S_{0B}$. On the set $E$ of twist-one labels outside $I$,
the coefficient-of-unit character formula on $F(B)$ gives
$\operatorname{Tr}_K(w_B)=-h_B$. This set is stable under the
modular Galois action, including the possible column signs.
The rational row $h$ is fixed under the same signed permutations.
Consequently $3\sum_E w_B h_B=-\sum_E h_B^2$. By Cauchy's
inequality the extra mass $\sum_E w_B^2$ is at least $1/192$.
For the probability measure of weights $w_B^2$ write $P$ for its
atom at one and $M_j=\Re\sum_B w_B^2\theta_B^j$. We obtain
\[
 P\ge 1/16+1/256+(s-p)^2/2+1/192
 \ge 473/1536-(11/16)p.
\]

There is no nonunit of dimension at most two, by Theorem~\ref{thm:perfect-small-simple}, so $D>17$. Cubic Fourier inversion gives $p_j\le(s+2/\sqrt7)/3<482/1000$. The two nonidentity reciprocals are therefore each greater than $s-1/17-482/1000>146/1000$, and at least one is greater than $(s-1/17)/2>314/1000$. The integrality of $\mathrm N(D/16)$ gives $D\ge16^3p'p''>187$. Inserting this back gives $p_j>s-1/187-482/1000>1/5$ for both. Their codegrees consequently lie between two and five.
Put $x=D^{-1/2}$ and $c_0=137/242$. We shall use
\[
 M_3\le 7x+(c_0^{1/4}+5x/2)^4 .
\]
Indeed $|\nu_3(Y)|\le [Y^3:1]=N_{YY}^Y$ by rotation and
self-duality. In the character formula for
$D^{-1}\sum_i d_i N_{ii}^i$, the positive character contributes
$\sum d_i^4/D^2$. Each other dimension character contributes at most
$\sqrt5/\sqrt D$, by Cauchy's inequality and
$\sum_i\sigma(d_i)^6\le\sigma(D)^3$.
The two rational character contributions are at most
$(1/2+\sqrt{796}/16)/\sqrt D$, by the same inequality and their
explicit absolute entries. Their combined coefficient on $x$ is at most $2\sqrt5+1/2+\sqrt{796}/16<7$. Here the sum of sixth powers of $\tau$ is
$729+64+1+1+1=796$.

The nonunit dimensions are $u,2u,v,w$ with $w=\pm u\pm v\pm1$
by the $\rho$-row. Change the unit coordinate to zero and remove
the $\pm1$ perturbation from $w$, obtaining a real vector $d_0$.
It satisfies $\|d_0\|_4^4\le c_0\|d_0\|_2^4$:
after changing signs and dividing by $u$, the entries other than the unit are $1,2,v,1+v$ up to signs, and
\[
 c_0(5+v^2+(1+v)^2)^2-
       (17+v^4+(1+v)^4)
 =\frac8{121}(2v+1)^2(v^2+v+36)\ge0 .
\]
Thus $\|d\|_4/\sqrt D\le c_0^{1/4}+5/(2\sqrt D)$ by the
triangle inequality and $c_0^{1/4}\sqrt2+2^{1/4}<5/2$.
This proves the stated bound on the third moment.

The other moments satisfy $M_1=p$ and
$M_2\le(1+2\sqrt{D-1})/D$. Integrating
$|7+11z+11z^2+7z^3|^2$ against the twist measure gives
\[
 1296P\le340+550 M_1+308 M_2+98 M_3.
\]
Since $c_0^{1/4}<7/8$, expansion of the fourth power gives the following necessary inequality:
\[
 \frac{13995}{3872}\le
 \frac{250691}{128}x+\frac{292011}{64}x^2
 +\frac{42875}{8}x^3+\frac{30625}{8}x^4 .
\]
The right side is strictly increasing with $x$. The first integer $D$ at which it fails is $296143$, by exact squaring. Thus $D<296143$. The checker derives this upper bound from the coefficients.
 
We finish with the integral lattice enumeration in \texttt{cubic\_polar\_sieve.py}. Write $D=A+B\alpha+C\alpha^2$ with integer
coefficients, using $\alpha=2\cos(2\pi/7)$ and its integral
power basis. Relabel by a field automorphism to put the largest
embedding at the largest root of $Q(X)=X^3+X^2-2X-1$.
The three roots belong respectively to
$(1.246979,1.246980)$, $(-0.445042,-0.445041)$ and
$(-1.801939,-1.801937)$ by substitution.
Lagrange interpolation gives $C=\sum_i D_i/Q'(\alpha_i)$,
so $0<C<D/5+5$ by the displayed intervals and $D>187$. The difference of the last two embeddings gives
$|B-C(1+\alpha_0)|<3/(\alpha_1-\alpha_2)<2.22$.
For each such $C$, the checker bounds the integers $B$ by this inequality and then $A$ by $2<A+B\alpha_1+C\alpha_1^2<5$, using the rational endpoints of the root intervals. For each triple use the norm polynomial and its $A$-derivative to
compute $n=\operatorname N(D)$ and the second elementary symmetric
sum $e_2$. Enforce $16e_2=11n$ and that $n$ is a positive square,
the latter by the based discriminant $4\cdot16\cdot n$ and the
field discriminant $49$. The enumeration gives no triple.
\end{proof}

\begin{theorem}\label{thm:cubic-intrinsic}
If the global Frobenius--Perron dimension $D$ of a perfect rank-five fusion category has degree three, then its dimension field has conductor seven, its complementary codegrees are $7,7$, and $D$ is the largest root of
\[
                X^3-42X^2+245X-343 .
\]
In particular $35<D<36$. This includes the case of a nonselfdual simple.
\end{theorem}
\begin{proof}
Theorem~\ref{thm:perfect-small-simple} implies that every nonunit dimension exceeds two, so $D>17$. Write $s=1-1/r-1/t$. Proposition~\ref{prop:r5-cubic-frontier}, Corollary~\ref{cor:r5-cubic-even}, Proposition~\ref{prop:cubic-finite-profiles} and Corollary~\ref{cor:cubic-profile} place every profile in one of the lines below, with the sole further profile excluded by Proposition~\ref{prop:cubic-polar}. The column $s_0$ gives a lower bound for $s$. In an odd line set $H(s,F)$ equal to the first integer $h\ge2$ with $(s-1/\sqrt F)h>3+3\sqrt{h-1}$. In every other line let $H(s,F)$ be the first such integer satisfying $(s-2/\sqrt F)h>3$. These bounds are finite because the coefficients are positive in their respective lines. Both right sides after division by $h$ decrease with $h\ge2$. All admissible dimensions therefore satisfy $D<H(s,F)$.
\[
\begin{array}{c|c|c|c}
 F & \text{codegree restriction}&s_0&H(s_0,F)\\\hline
 F\ge19 & \text{all}&1/2&\le73\\
 13 & \text{all}&7/12&105\\
 7,9&r,t\text{ odd}&3/5&208,148\\
 7,9&r\ge7,\ \text{at least one even}&65/84&168,29\\
 9&r=4,\ v_2(t)=4&11/16&145
\end{array}
\]
The individual bounds $H(s,F)$ suffice to enumerate the minimal polynomial. Here are the full loops. Put $N=\mathrm N(D)$ and $P=X^3-AX^2+sNX-N$. All three conjugates of $D$ exceed one. Of the two others at least one is less than five, since their reciprocal sum exceeds $1/2-1/17>2/5$. Consequently $N<5H^2$ and $A<2H+5$, where $H=H(s,F)$; for $F\ge19$ one can always use $H(s,19)$. Central induction gives $\operatorname{lcm}(r,t)^3\mid N$, bounding $r,t$ already with $H(s_0,F)$. Loop over these integers, using also that $r=3$ forces $t=6$ and $(4,4)$ is impossible. For each multiple $N$ in range retain integral $sN$ and $A$ with $N\mid A^3$ by the $d$-number criterion. The squared Fourier modulus is $s^2-3A/N\le1/F$. The discriminant of $P$ is a positive square by cyclicity. The field discriminant $F^2$ divides both $\operatorname{disc}(P)$ and the absolute based-order discriminant $rtN$. For $F\ge19$ it suffices to require their greatest common divisor to be divisible by the square of some integer at least nineteen. When $r\ne t$ the two complementary rows are rational and $rtN$ itself must be a square. Finally require $P$ irreducible with all three roots in $(1,H)$ and at least one in $(17,H)$.

The exact program \texttt{small\_cubic.py} leaves only $(r,t,F,A,sN,N)=(7,7,7,42,245,343)$. Its largest root lies in $(35,36)$ by rational root isolation.
\end{proof}




% ===== field_frontier.tex =====
\section{Perfect categories in quartic degree}
\label{sec:residual-fields}
Assume throughout this section that $\mathcal C$ is perfect, every
nonunit simple has dimension greater than two, and $K=\mathbb Q(D)$
is quartic. Thus $D>17$. We first reduce the field and rational
codegree to a short list, then exclude that list by twist moments
and exact arithmetic.

\begin{proposition}\label{prop:quartic-field-list}
Suppose $K=\mathbb Q(D)$ is quartic and $D>17$. Let $f$ be the complementary rational codegree. The following table contains all pairs $(f,K)$ allowed by the norm, ramification and Fourier tests below:
\[
\begin{array}{c|l|l}
 f&\text{cyclic conductors}&\text{biquadratic subfield discriminants}\\\hline
 4 &15,16,20,35,40 &(5,8,40),(5,12,60),(5,13,65),\\
   &&(5,17,85),(5,21,105),(8,12,24)\\
 5 &15,20& \\
 16&15&(5,8,40)\\
 20&15&
\end{array}
\]
In quartic degree all characters are real and hence every simple is self-dual. The kernels defining the five cyclic fields, in conductor order, are
$\{\pm1\},\{\pm1\},\{\pm1\},\{\pm1,\pm11,\pm16\},\{\pm1,\pm11\}$.
\end{proposition}
\begin{proof}
Write $p=1/D$, $p_g=g(p)$, $s=\sum_g p_g=1-1/f$ and $N=\mathrm N(D)$. If $c_1\le c_2\le c_3$ are the nontrivial character conductors of $K$, then $\Delta=\operatorname{disc}(K)=c_1c_2c_3$. Let $B$ denote the dimension order and $h=[\mathcal O_K:B]$. The based order $R\subset\mathbb Z\times B$ has regular discriminant $fN$, the product of its formal codegrees. The primitive rational idempotent has unit coefficient $1/f$ and $f$ annihilates its denominators. Consequently $R\cap(\mathbb Q\times0)=f\mathbb Z\times0$ and $[\mathbb Z\times B:R]=f$. Also $D/f$ is the dimension of a constituent of central induction and is integral over $\mathbb Z$. Thus
\begin{equation}\label{eq:quartic-norm-index}
 N=f\Delta h^2,\qquad f^4\mid N,\qquad (4/e_\ell)\mid v_\ell(N),
\end{equation}
where $e_\ell$ is the ramification index in $K$. The last assertion holds because the $d$-number property gives equal valuations for $D$ at the primes above $\ell$. Precisely $4/e_\ell$ of the four field characters, including the trivial one, are unramified at $\ell$.

Codegree one is impossible by the mass equation. At $f=2$ the rational character has a single nonzero nonunit entry $\pm1$, forcing an invertible by orthogonality. At three the inertia is tame and its order divides two. Thus $v_3(N)$ and $v_3(\Delta)$ are both even, the latter by the tame discriminant formula. Equation~\eqref{eq:quartic-norm-index} excludes $f=3$. Hence $s\ge3/4$. Also $f\equiv1\bmod8$ is impossible. The four nonunit integer entries of the rational row have square sum $f-1$ and cannot all be even by orthogonality with the dimension row. Four squares with sum divisible by eight have all four entries even.

Put $S=\sum_i c_i^{-1/2}$. Proposition~\ref{prop:ordinary-fourier} and Fourier inversion at $p$ give $s-4p\le S$. Pairwise application of the same bound gives
\[
 p_g\ge a_g-p,\qquad a_g=(s-b_g)/2\quad(g\ne1).
\]
For an involution $g$, $b_g$ is the inverse square root of the conductor of the unique nontrivial character trivial on $g$. Indeed its Fourier coefficient is $2(p+p_g)-s$. For an element $g$ of order four, projecting the faithful complex Fourier coefficient onto a diagonal gives $b_g=\sqrt{2/c_2}$. Here the faithful conductor $c_2=c_3$ is at least fifteen, as one sees from $(\mathbb Z/n)^\times/\{\pm1\}$ for $n<15$. Thus every $b_g\le1/\sqrt5$. In particular all $a_g>1/17$. For every lower bound $N_0\le N$ we obtain the necessary conditions
\begin{equation}\label{eq:quartic-interval}
 p\in I=\left[\max\{0,(s-S)/4\},1/17\right],\qquad
 \Phi(p)=N_0\,p\prod_{g\ne1}(a_g-p)\le1.
\end{equation}
In fact $1/N=p\prod_{g\ne1}p_g$. On the interior of $I$ every displayed linear factor is positive, and $\log\Phi$ is concave. The minimum of $\Phi$ on a nonempty $I$ is therefore attained at an endpoint, including when the left endpoint is zero. Larger upper bounds on $S,b_g$ and smaller lower bounds on $s,N$ preserve necessity.

Two applications give finite conductor and codegree ranges. If $c_2\ge47$, use $s_0=3/4$, $N_0=20\cdot47^2$ from $N\ge f\Delta$, $S_0=1/\sqrt5+2/\sqrt{47}$ and the three bounds $1/\sqrt5,\sqrt{2/47},\sqrt{2/47}$ for $b_g$ in \eqref{eq:quartic-interval}. The respective endpoint values exceed $1.31$ and $10.9$. Thus $c_2\le46$. Enumerating biquadratic fields from their two smallest fundamental discriminants and cyclic fields from the index-four kernels containing $-1$ of conductor at most $46$ gives
$S\le S_* =1/\sqrt5+2/\sqrt{15}$. If $f\ge28$, take $s_0=27/28$, $N_0=28^4$, $S_0=S_*$ and $b_g\le1/\sqrt5$ in \eqref{eq:quartic-interval}. The endpoint values exceed $1.78$ and $288$. Hence $f\le27$.

For each remaining pair use \eqref{eq:quartic-norm-index} to compute $N_0$. At each prime $\ell\mid f\Delta$ the exponent is the smallest integer $v\ge\max\{v_\ell(f\Delta),4v_\ell(f)\}$ congruent to $v_\ell(f\Delta)$ modulo two and divisible by $4/e_\ell$. Incompatible congruences exclude the pair. Otherwise the product of the resulting prime powers divides $N$. Testing the two endpoints with this $N_0$ gives exactly the displayed table. The complete calculation in \texttt{quartic\_fields.py} uses $10^{-8}(1+\lfloor10^8\sqrt{m/n}\rfloor)$ for each upper bound on $\sqrt{m/n}$, so every comparison is rational.
\end{proof}


\subsection{Local restrictions and the rational row}

% ===== quartic_two_support.tex =====
\begin{proposition}\label{prop:quartic-two-support}
In the perfect quartic case with rational codegree four, write $R$ for
its unit-induction constituent in the center. Then $S_{RB}=0$ whenever
$\theta_B^2\ne1$. In the integral row $(4S_{RB})_B$, the sum of the
squares of the entries is eight on each of the two twists $+1,-1$.
\end{proposition}
\begin{proof}
Use the positive spherical structure and put $r_B=S_{RB}=\rho(FB)/4$,
$w_B=S_{0B}=\dim(FB)/D$, and $l_B=[F(B):1]$.
Here $\rho$ is the rational character of the based ring. All $r_B$ are
rational, and the coefficient-of-unit character formula reads
$\operatorname{Tr}_{K/\mathbb Q}(w_B)=l_B-r_B$.
Write $I$ for the five multiplicity-one labels of $I(1)$.
For each of its four dimension-character labels, $r_B=1/4$ by
Galois symmetry and the rational positive vacuum entry $r_0=1/4$.
Also $\sum_{B\in I}r_B=1$, by applying $\rho$ to the Casimir
$F(I(1))$. Thus $r_R=0$.

Let $N$ be the twist exponent and, for any integer $j$, set
$m_j=\sum_B w_B r_B\theta_B^j$. The central indicator formula gives
\[
 4m_2=\sum_i\rho(X_i)\nu_2(X_i) .
\]
Indeed expand $\rho(FB)$ over the simple constituents $X_i$ of $FB$
and apply the indicator formula to each $I(X_i)$.
The rational character row has exactly four entries of magnitude one
and one zero: the alternative of a lone nonunit entry of square three
is impossible. All simples are self-dual, so their second indicators
are $\pm1$. In particular $4m_2$ is an even integer.

For $u\in(\mathbb Z/N)^\times$ let $\sigma_u$ act by the $u$th
power on roots of unity, and write $B\mapsto\widehat u B$ for the
modular label permutation with column sign $\epsilon_u(B)$. Thus
\[
 \sigma_u(w_B)=\epsilon_u(B)w_{\widehat uB},\quad
 r_B=\epsilon_u(B)r_{\widehat uB},\quad
 \theta_{\widehat uB}=\theta_B^{u^2}.
\]
The last identity uses the anomaly-free normalization of the center.
Changing variables gives $\sigma_u(m_j)=m_{j/u}$, with exponents
taken modulo $N$. Rationality of $m_2$ therefore gives
$m_{2/u^2}=m_2$, and a second change of variables gives
\[
 \sum_B\sigma_u(w_B)r_B\theta_B^2=m_{2/u^2}=m_2.
\]
Choose a cyclotomic lift for each of the four automorphisms of $K$,
sum, and apply the trace formula. Every label of $I$ has twist one,
so this yields
\[
 4m_2+\sum_B r_B^2\theta_B^2
 =\sum_B l_B r_B\theta_B^2=1 .
\]
The second term has absolute value at most $\sum_B r_B^2=1$ by
unitarity. Since $4m_2$ is even, that term is $1$ or $-1$.
Its vacuum summand has positive weight and squared twist one, excluding
$-1$ by equality in the triangle inequality. Equality at $1$ forces
$\theta_B^2=1$ on the whole support. Finally the first twist moment
$\sum_B r_B^2\theta_B=S_{RR}=0$ by $STS=T^{-1}ST^{-1}$.
The squared masses at both signs are therefore $1/2$; multiplying by
sixteen proves the integer count.
\end{proof}



% ===== quartic_twenty.tex =====
\begin{proposition}\label{prop:quartic-twenty}
The perfect quartic profile cannot have rational codegree $20$.
For rational codegree $16$, the corresponding central row is supported
either on twists $1,-1$ or on twists $1,i,-i$.
\end{proposition}
\begin{proof}
Let $f\in\{16,20\}$, let $J$ denote the rational induction
constituent, and use $r_B=S_{JB}=\rho(FB)/f$ and $w_B=S_{0B}$.
Each of the four dimension-orbit induction constituents has $r_B=1/f$
by modular Galois symmetry. Since $\sum_{B\subset I(1)}r_B=1$,
the first diagonal twist moment is
$\sum_B r_B^2\theta_B=S_{JJ}=1-4/f$.
For $m_j=\sum_B w_B r_B\theta_B^j$, the indicator formula gives
$f m_j=\sum_i\rho_i\nu_j(X_i)$, where $\rho_i=\rho(X_i)$.
Set $k=fm_2$. All $X_i$ are self-dual, so $\nu_2(X_i)=\pm1$.
Thus $k\equiv\sum\rho_i^2=f\equiv0\pmod2$.
The Galois trace calculation of Proposition~\ref{prop:quartic-two-support}
uses only the rationality of $m_2$, the quartic dimension field and
$\sum_{B\subset I(1)}r_B=1$.
It yields $\sum_B r_B^2\theta_B^2=1-4k/f$.
For the probability weights $r_B^2$, the first moment has real part
$1-4/f$. The inequality $\mathbb E(\cos^2 t)\ge(\mathbb E\cos t)^2$
therefore gives
\[
 0\le k\le \frac f2\bigl(1-(1-4/f)^2\bigr)=4-8/f.
\]
Thus $k=0$ or $2$. At $k=0$, equality in the triangle inequality for
the second moment implies $\theta_B^2=1$ throughout the support.

At $k=2$ take the cyclotomic trace of the two rational moment
identities, keeping the rational weights $r_B^2$ fixed.
A twist of order $n$ contributes its normalized Ramanujan sum
$a_j(n)=\mu(n/(n,j))/\varphi(n/(n,j))$ to the $j$th moment.
Lemma~\ref{lem:odd-row}, applied to the modular formal codegree
$f^2$, shows that every $n$ on the support has prime factors in
$\{2,5\}$.
For $n=1,2,4,5,10,20$ the pairs $(a_1,a_2)$ are respectively
\[
 (1,1),\ (-1,1),\ (0,-1),\ (-1/4,-1/4),\
 (1/4,-1/4),\ (0,1/4);
\]
all other orders with these prime factors give $(0,0)$.
Hence $1-a_2\le2(1-a_1)$ throughout the support, with equality
only at orders $1,4$. The traced moments give equality after averaging
since their deficits are $4/f,8/f$, so $\theta_B^4=1$ throughout,
and the twist of order two does not occur when $k=2$.

For clarity, there is a general prime-divisor obstruction underlying the last step. For a rational character of a commutative fusion ring in a positive spherical category, let $L$ be the least common multiple of the orders of the twists on its central induction row. Then every prime divisor of its formal codegree $f$ divides $L$. Indeed, if $p\mid f$ and $p\nmid L$, choose a positive power $n=p^t\equiv1\pmod L$. The central indicator formula gives
\[
  \sum_i\rho_i\nu_n(X_i)=fm_n=fm_1=1.
\]
At a prime over $p$ in $\mathbb Q(\zeta_n)$, every eigenvalue of the $n$-fold rotation on $\operatorname{Hom}(1,X_i^{\otimes n})$ reduces to one. Consequently $\nu_n(X_i)\equiv[X_i^n:1]$. Put $a=\sum_i\rho_i X_i$; a rational character takes the same value on a simple and its dual, so $a$ is the numerator of its primitive idempotent, $a^2=fa$ and $[a:1]=1$. Frobenius in the commutative fusion ring modulo $p$ gives
\[
 1\equiv\sum_i\rho_i[X_i^n:1]\equiv[a^n:1]=f^{n-1}\equiv0\pmod p.
\]
For $f=20$ the support just proved has $L\mid4$, and we can take $p=n=5$. This excludes $f=20$.
\end{proof}



% ===== quartic_16_depth.tex =====
\begin{proposition}\label{prop:quartic-dyadic-depth}
In the cyclic quartic field of conductor fifteen the rational codegree cannot be sixteen.
\end{proposition}
\begin{proof}
Write $f=16$, $a=v_2(f)=4$, $e=v_2(D)$, and $O$ for the integers in the completion of $K$ at two, an unramified extension with Frobenius $\sigma$ of order four. Let $B$ be the projected dimension order over $\mathbb Z_2$ and $\rho$ the rational row. The projected kernel on the rational axis is $f\mathbb Z_2$: its primitive numerator has unit coefficient one. Thus there is a unital gluing map $\phi:B\to\mathbb Z_2/f\mathbb Z_2$ carrying $d_i$ to $\rho_i$. Its residue is the unique residue character of $B$, since $O$ is local and integral over $B$, so all $d_i$ reduce to the corresponding $\rho_i$ in $\mathbb F_{16}$. The four central dimension ratios $D/\gamma D$ in $B$ are units and reduce to one. The trace identity $(1-1/f)D=\sum_\gamma D/\gamma D$ then gives $e>a$ and $D/2^e\equiv1\pmod{2O}$ (the residue of this unit is fixed by Frobenius). The rational row has four odd entries, including the unit, and a single even entry at $Z$, by summing its five squares modulo four.

Suppose that $u_i=(d_i-\sigma d_i)/2^j$ are integral. Orthogonality of rational and quartic characters, and of distinct quartic characters, gives
\[
 \sum_i\rho_i u_i=0,\qquad
 \sum_i u_i^2=(D+\sigma D)/2^{2j},\qquad
 \sum_i u_i\sigma u_i=-\sigma D/2^{2j},\qquad
 \sum_i u_i\sigma^2u_i=0.
\]
For $e\ge2j+1$ the second sum vanishes modulo four and the third modulo two. Reduction of the first two sums modulo two gives $\bar u_Z=0$ and three residues $v,w,v+w$ at the remaining nonunit entries. Squares modulo four depend only on the residues modulo two, so $v^2+vw+w^2=0$. Unless all are zero, the three residues form $c\mathbb F_4^\times$ for some $c\ne0$. The sum of their cubes is then $c^3$, contradicting the third identity modulo two. Thus the congruence improves to $j+1$. Starting with $j=0$ reaches $j=\lceil e/2\rceil$, where integrality of the third sum forces $e=2k$ even. It follows that $d_i$ is a scalar modulo $q=2^k$: the Frobenius-fixed elements of $O/2^kO$ are $\mathbb Z/2^k$ (equivalently use the successive Teichm\"uller digits). Hence $B\subset\mathbb Z_2+qO$.

At depth $j=k$, the first two sums again vanish modulo two, so $u_Z\equiv0$ and the three remaining residues have the same form $v,w,v+w$. The third sum now has residue one, forcing $v,w,v+w$ nonzero. The last sum gives $vw^4+wv^4=0$, hence $v/w\in\mathbb F_4^\times$ and $v^2+vw+w^2=0$. Therefore $D+\sigma D\in4q^2O$, while $d_Z$ is a scalar $z$ modulo $2q$. Write $y=(d_Z-z)/(2q)$ and $\lambda=q^2/D$, a local unit of residue one. Orthogonality at the unit and $Z$ columns gives
\[
 \mathrm{Tr}(\lambda y)=-\frac q2\left(z+\frac{\rho_Z-z}{f}\right),\qquad
 \mathrm{Tr}(\lambda y^2)=\frac14\left(1+z^2-\frac{(\rho_Z-z)^2}{f}\right).
\]
The scalar $z$ is even. Integrality of the second trace gives $v_2(\rho_Z-z)=a/2$. The four nonunit squares sum to fifteen, forcing absolute values $3,2,1,1$ and $\rho_Z=\pm2$. Thus $z\equiv2\pmod4$, $(\rho_Z-z)/4$ is odd, and the second trace is odd. The traces have the same parity, since their reductions are the finite-field traces of $\bar y,\bar y^2$. The first has valuation $k-3$, so $k=3$.

The determinant identity $\mathrm N(D)=f\operatorname{disc}(B)$ at two gives $[O:B]=2^{2e-a/2}=2^{10}$. Since $B\subset\mathbb Z_2+8O$ with index two in that lattice, it contains $16O$. Let $U=D/\sigma D$. For the corresponding central simple the modular entry on the rational row is $1/f$ by Galois conjugation of the vacuum entry, so its rational character value is one. Its dimension is $U$, whence $\phi(U)=1$. Choose a scalar $c\equiv U\pmod{8O}$. We saw $D+\sigma D\in4q^2O$, so $c\equiv-1\pmod4$. Then $x=U-c$ belongs to $8O$ and $\phi(x)$ has valuation one modulo sixteen. But $x^3\in512O\subset32B$, forcing $\phi(x)^3=0$ modulo sixteen, a contradiction.
\end{proof}



% ===== quartic_f_table.tex =====
\begin{corollary}\label{cor:quartic-coupled}
For a perfect category with $D>17$ and quartic dimension field, the possible pairs $(f,K)$ are those in Proposition~\ref{prop:quartic-field-list} after deleting $(20,15)$ and $(16,15)$, where $15$ denotes the cyclic field of that conductor.
\end{corollary}
\begin{proof}
Proposition~\ref{prop:quartic-twenty} excludes codegree twenty. Proposition~\ref{prop:quartic-dyadic-depth} excludes codegree sixteen at conductor fifteen.
\end{proof}



% ===== quartic_exhaustion.tex =====
\subsection{Dimension bounds and the final exhaustion}\label{sec:quartic-exhaustion}
We now exclude every pair in Corollary~\ref{cor:quartic-coupled}. We first bound the global dimension or the least nonunit dimension, then enumerate the algebraic integers allowed by those bounds.
We work on the self-dual perfect locus on which every nonunit simple has dimension greater than two. In particular $D>17$.
\begin{lemma}\label{lem:quartic-conjugates}
Every nonidentity embedding satisfies $1<g(D)<11$.
\end{lemma}
\begin{proof}
Write $p_g=1/g(D)$, $p=1/D$ and $s=1-1/f\ge3/4$.
For an involution $g$, the quadratic character trivial on $g$ gives
$|s-2(p+p_g)|\le1/\sqrt5$. For an element of order four, projecting a faithful Fourier coefficient on a diagonal gives $|s-2(p+p_g)|\le\sqrt{2/c}$, where its conductor $c$ is at least fifteen.
Consequently $p_g\ge(3/4-1/\sqrt5)/2-1/17>1/11$. The unit term in the squared norm of the character row gives $g(D)>1$.
\end{proof}



% ===== quartic_sixteen.tex =====
\begin{proposition}\label{prop:quartic-sixteen}
The biquadratic field of discriminants $(5,8,40)$ does not occur with rational formal codegree sixteen.
\end{proposition}
\begin{proof}
Here $K=\mathbb Q(\sqrt2,\sqrt5)$, with group $G$ of exponent two.
Put $p=1/D$, $p_g=g(p)$, $s=\sum_g p_g=15/16$ and $w_B=S_{0B}$.
For $M_j=\sum_B w_B^2\theta_B^j$, the indicator formula gives
\[
 M_1=p,\qquad M_2=p\sum_i d_i\nu_2(X_i),\qquad
 |M_2|\le p+2\sqrt p.
\]
Moreover $\operatorname{Tr}_K(M_2)=1-k/16$, where
$k=\sum_i\rho_i\nu_2(X_i)\in\{0,2\}$ by Proposition~\ref{prop:quartic-twenty}.
Indeed $\operatorname{Tr}_K(d_i/D)=\delta_{i,1}-\rho_i/16$ by column orthogonality.

We bound the vacuum mass $W=\sum_{\theta_B=1}w_B^2$ from above.
For a root of unity $z$ and $j=1,2$, let $A_j(z)$ be the average of
$\operatorname{Re}(z^{ju})$ over cyclotomic units $u$ fixing $K$.
Let $C_j(z)$ be four times the average over all units, and
let $C_{5,1}(z)$ be four times the same average at $j=1$ with the
quadratic character of conductor five inserted. All averages can be
taken at any common multiple of forty and the order of $z$. We use
the elementary pointwise inequality
\[
 1000\,\mathbf1_{z=1}\le
 337+472(A_1(z)+A_2(z))+7 C_1(z)-76 C_2(z)-49 C_{5,1}(z).
 \tag{*}
\]
Here are the full ranges for its finite verification.
Let $b_{c,j}$ be the all-unit average with the real primitive character
of conductor $c=1,5,8,40$ inserted. Then $A_j=\sum_c b_{c,j}$,
$C_j=4b_{1,j}$ and $C_{5,1}=4b_{5,1}$.
Writing $z^j=\exp(2\pi i t/n)$ in lowest terms, one has
\[
 b_{c,j}=
 \begin{cases}
 \mu(d)\chi_c(td)\sqrt c/\varphi(n),&
 n=cd,\ (c,d)=1,\\
 0,&\text{otherwise},
 \end{cases}
\]
by the primitive quadratic Gauss sum and the Chinese remainder
theorem (including $c=1$). If the order of $z$ is divisible by
$32$ or by $25$, all the averages in (*) vanish. Otherwise write
its order as $n_0 m$ with $n_0\mid80$ and $(m,10)=1$.
The odd factor $m$ multiplies every average by
$\mu(m)/\varphi(m)$, after changing the numerator modulo $n_0$
by $m^{-1}$. This multiplier lies in $[-1/2,1]$. By linearity it
therefore suffices to check $m=1,3$ and all reduced numerators.
For $n_0=m=1$ one also checks the indicated unit mass, which is
stronger than the inequality at the endpoint with zero unit mass.
The accompanying \texttt{verify16.py} checks these values with rational
bounds for the three square roots.

For cyclotomic units $u$, the signed permutation of the modular
columns gives $\sigma_u(M_j)=M_{j/u}$. Both $M_1,M_2$ belong to $K$.
Averaging (*) against the probability weights $w_B^2$ consequently yields
\[
 1000 W\le337+472(p+M_2)+7s
        -76(1-k/16)-49\widehat p(\chi_5).
\]
Label the nonidentity reciprocals $u,v,w$ by the conductors $5,8,40$ of the characters whose kernels contain the corresponding involutions. The Fourier coefficient retained above is $\widehat p(\chi_5)=2(p+u)-s$. Using $k\le2$, $s=15/16$ and $|M_2|\le p+2\sqrt p$ gives
\[
                  1000W+98u\le323+846p+944\sqrt p .
\]
The four dimension constituents and the rational constituent of $I(1)$ all have twist one. Consequently $W\ge p^2+u^2+v^2+w^2+1/256$. Using $x^2\ge2tx-t^2$ at $t=250/1000,299/1000,390/1000$ for $u,v,w$ respectively gives
\[
  1000W+98u\ge 1000/256-304001/1000+598(s-p)+182w .
\]
Indeed the first two linear coefficients coincide after adding $98u$. The Fourier bound for conductor forty gives $w\ge(s-1/\sqrt{40})/2-p>37/96-p$.

It remains to compare the two inequalities at the reciprocal size imposed by the norm. Equation~\eqref{eq:quartic-norm-index} gives $N=16(5\cdot8\cdot40)h^2$ and $16^4\mid N$, so $2^{16}\cdot25\mid N$. The three pair bounds with $p<1/17$ give $u>1/6$, $v>1/5$ and $w>3/10$. Hence $D=Nuvw>2^{16}\cdot25/100=128^2$. Substituting $w>37/96-p$ in the lower bound, the difference between that bound and the upper bound decreases with $p$. Already at $p=1/128^2$ it equals $620053/3072000>0$. This excludes the field.
\end{proof}



% ===== quartic_moment_bounds.tex =====
\begin{proposition}\label{prop:quartic-four-effective}
At rational formal codegree four, the following are strict upper bounds for the global Frobenius--Perron dimension. A single integer denotes the cyclic quartic field of that conductor.
\[
\begin{array}{c|rrrr}
K&16&20&35&40\\\hline
D<&13883&3121&1019&495
\end{array}
\qquad
\begin{array}{c|r}
K&D<\\\hline
(5,8,40)&10500\\
(5,12,60)&1111\\
(5,13,65),\ (5,17,85),\ (5,21,105)&182,\ 260,\ 317\\
(8,12,24)&44123
\end{array}
\]
\end{proposition}
\begin{proof}
Let $p=1/D$, $p_g=g(p)$ and $s=\sum_{g\in G}p_g=3/4$ for
$G=\operatorname{Gal}(K/\mathbb Q)$. Let $w_B=S_{0B}$ and
$r_B=S_{RB}$ for the rational induction constituent $R$.
We integrate functions on the circle with the probability weights
$w_B^2$ at the twists $\theta_B$. Write $W_+,W_-$ for the masses
at $1,-1$. Proposition~\ref{prop:quartic-two-support} gives
$\sum_{\theta_B=\pm1}r_B^2=1/2$ at each sign.
Of this mass, $1/4$ lies on the induction constituents $I$ at
twist one. On either of the Galois-stable sets
$E=\{B:\theta_B=\pm1\}\setminus I$ the trace identity
$\operatorname{Tr}_{K/\mathbb Q}(w_B)=-r_B$ and Cauchy's inequality
give $\sum_E w_B^2\ge\sum_E r_B^2/16$. Indeed each signed column
permutation preserves the sum of squares on $E$.
Consequently
\[
 W_+\ge \sum_g p_g^2+1/16+1/64,\qquad W_-\ge1/32.        \tag{1}
\]

For a root of unity $z$ let $A_{g,j}(z)$ be the average of
$\operatorname{Re}(z^{ju})$ over cyclotomic units $u$ whose inverse
restricts to $g$ on $K$. Let $C_j=\sum_g A_{g,j}$, and write
$M_j=\sum_B w_B^2\theta_B^j$. By the signed permutation of the
vacuum row, $\sigma_u(M_j)=M_{j/u}$. In particular
$\int A_{g,1}=p_g$, $\int A_{1,2}=M_2$ and
$|M_2|\le p+2\sqrt p$ by the second indicator formula.
Suppose $K\cap\mathbb Q(\zeta_j)=\mathbb Q$. From
$M_j=p\sum_i d_i\nu_j(X_i)$, $\nu_j(X_i)\in\mathbb Q(\zeta_j)$
and $\operatorname{Tr}_K(p d_i)=\delta_{i,1}-\rho_i/4$,
the independent averages over the two fields give
\[
 \int C_j=1-\tfrac14\operatorname{avg}\operatorname{Re}
                     \sum_i\rho_i\nu_j(X_i)
    =\begin{cases}3/4,&j\text{ odd},\\1,&j\text{ even}.\end{cases}   \tag{2}
\]
In fact the scalar sum on the left is already rational: it is
$4\sum_B w_Br_B\theta_B^j$, which is one for odd $j$ and zero
for even $j$ by the support of $r$, the first indicator formula and
row orthogonality. We use the $j\le6$ for which none of the
quadratic subfield conductors divides $j$. This condition ensures
the required disjointness.

We describe how to verify the ten numerical bounds by finite rational
inequalities.
For each field, a rational vector $t=(t_g)$ and a pointwise
majorant for
\[
 \mathbf1_{z=1}-2\sum_g t_g A_{g,1}(z)
\]
are given by \texttt{verify\_quartic\_moments.py}. The majorant is a
linear combination of $1$, the $C_j$ from (2), $A_{1,1}$,
$\pm A_{1,2}$, real projections of the Fourier coefficients of
$(A_{g,1})_g$, and $-\mathbf1_{z=-1}$.
The coefficients of the inequalities are nonnegative. Their integrals
are bounded respectively by $1/D$, $1/D+2/\sqrt D$, the Fourier
estimate $1/\sqrt{c_\chi}$ for each unit projection at a nontrivial
character of conductor $c_\chi$, and $-1/32$. The remaining integral
values are exact from (2). After moving the constant terms to the right, the condition for a contradiction is
\[
 A/D+B/\sqrt D<\Delta,
\]
with $A,B\ge0$ and $\Delta>0$ given by the majorant. Indeed (1) and $\sum p_g^2\ge2\sum t_gp_g-\sum t_g^2$ give the lower bound $1/16+1/64-\sum t_g^2$. The table records, for each fixed certificate and its stated rational enclosures, the least integer $H$ for which $\Delta H>A$ and $(\Delta H-A)^2>B^2H$. Monotonicity then gives $D<H$.

For clarity we specify why the certificates check every possible
order of $z$. Let $c$ be the field conductor, $J$ the least common
multiple of the utilized exponents and
$L=J\operatorname{lcm}(c,\operatorname{rad}J)$.
A prime power in the order of $z$ involving a prime of $L$ but exceeding
its power in $L$ makes all the relevant averages zero. Indeed after
raising to any utilized $j$, its exponent still exceeds both one and
the conductor exponent, so the corresponding primitive additive sum
on the unit coset vanishes. An order factor $m$ prime to $L$
multiplies every cosine average by $\mu(m)/\varphi(m)$ after a common
invertible change of numerator, by the Chinese remainder theorem.
This number ranges from $-1/(\ell-1)$ to one, where $\ell$ is the
least prime not dividing $L$. It is enough to check both endpoints
and orders dividing $L$. At the upper endpoint using the original
two real atoms strengthens the inequality if an extra factor is present.
For orders with all cosine averages zero the majorant is its positive
constant. For an order dividing $L$, elementary quadratic characters
or the four explicit cyclic unit cosets give the averages. The verifier
encloses each required cosine by its degree-$26$ Taylor polynomial on
$[-\pi,\pi]$ with a rational error bound, and checks the integrated
bounds with rational square-root enclosures. It calls no optimization
algorithm.
\end{proof}



% ===== five_support_masses.tex =====
\begin{lemma}\label{lem:quartic-five-row-masses}
Let $J$ be the rational induction constituent in the quartic case of formal codegree five. Put $r_B=S_{JB}$, and let $A_j=25\sum_{\operatorname{ord}(\theta_B)=j}r_B^2$. The only nonzero $A_j$ are at $j=1,5,25$, and the triples $(A_1,A_5,A_{25})$ are $(9,16,0)$ and $(5,0,20)$. For $W(j)=\sum_{\operatorname{ord}(\theta_B)=j}S_{0B}^2$, one has
\[
 W(1)\ge\sum_g p_g^2+\frac1{25}+\frac{A_1-5}{400},
 \qquad W(5)\ge\frac{A_5}{400},\qquad W(25)\ge\frac{A_{25}}{400}.
\]
\end{lemma}
\begin{proof}
The entries $5r_B=\rho(FB)$ are integers of squared norm $25$. Lemma~\ref{lem:odd-row} excludes all primes but five from the orders of supported twists. If one has order $5^a$, its square-power Galois orbit contains at least $\varphi(5^a)/2$ distinct twists, each with the same $|5r_B|\ge1$. Thus $a\le2$, and $A_{25}$ is divisible by ten. The first twist moment of $r_B^2$ is $S_{JJ}=1/5$. Its cyclotomic average gives $A_1-A_5/4=5$. In conjunction with $A_1+A_5+A_{25}=25$, this leaves $A_{25}=0,10,20$.

To distinguish the three possibilities put $w_B=S_{0B}$, $m_k=\sum_B w_Br_B\theta_B^k$ and $a_k(n)=\mu(n/(k,n))/\varphi(n/(k,n))$. The fifth-indicator formula gives $\nu=5m_5=\sum_i\rho(X_i)\nu_5(X_i)\in\mathbb Z[\zeta_5]$. Modular Galois symmetry gives $\sigma_u(m_5)=m_{5/u}$ for each unit $u$ of the exponent. Its normalized cyclotomic trace is therefore
\[
 \operatorname{Av}(\nu)=5\sum_B w_Br_B a_5(\operatorname{ord}(\theta_B)).
\]
For any function $H$ of the twist order, the signed permutation rule for $w,r$ gives
\[
 4\sum_B w_Br_B H_B
 =\sum_B\operatorname{Tr}_{K/\mathbb Q}(w_B)r_B H_B
 =\sum_B (l_B-r_B)r_B H_B,
\]
where $l_B=[F(B):1]$. All five induction labels have twist one, their $r_B$-entries are $1/5$, and $a_5(1)=a_5(5)=1$, $a_5(25)=-1/4$. Hence $\operatorname{Av}(\nu)=A_{25}/16$, which implies
$\operatorname{Tr}_{\mathbb Q(\zeta_5)/\mathbb Q}(\nu)=A_{25}/4\in\mathbb Z$. This excludes $A_{25}=10$ without any reality assertion for $\nu$.

For any fixed order $j$, remove the five unit-induction labels and call the remaining labels of order $j$ the set $E$. This is Galois-stable, and the same trace identity gives $4\sum_E w_Br_B=-\sum_Er_B^2$. Cauchy's inequality yields $\sum_Ew_B^2\ge\frac1{16}\sum_Er_B^2$. Restoring the induction labels, present only at order one, proves the three lower bounds.
\end{proof}



% ===== quartic_third.tex =====
\begin{lemma}\label{lem:quartic-least-third}
Suppose $(c,f)=(15,4),(15,5)$ or $(20,5)$ in the cyclic quartic case. Let $m$ be the least nonunit dimension and let $t=\Re M_3$ be the real part of the third vacuum twist moment. For any $u\le m$ set
\[
 \begin{gathered}
 P=(1+4u^2)^{-1},\qquad s=1-1/f,\qquad
 V_i=\left((s-b_i)/2-P\right)^{-1},\\
 (b_1,b_2,b_3)=(1/\sqrt5,\sqrt{2/c},\sqrt{2/c}).
\end{gathered}
\]
Put $V=\sum_i V_i$, $B=\sqrt{\max_i V_i-1}$, and $L=3/4$ at $f=4$, $L=11/20$ at $f=5$. Whenever each $V_i>0$ and $u>B+\tfrac12\sqrt{V/L}$, one has
\[
 t\le \tfrac14+(BL+\sqrt{VL})/u.
\]
\end{lemma}
\begin{proof}
Write $p=1/D$, $d_i=\operatorname{FPdim}(X_i)$,
$\rho_i=\rho(X_i)$, and $t_i=\Re\nu_3(X_i)$. These last numbers lie
in $\frac12\mathbb Z$, since the rotation has order dividing three.
The rational character has values in $\{0,1,-1\}$. At $f=5$ the sole
alternative would have nonunit entries $(\pm2,0,0,0)$, which violates
orthogonality to the dimension row since every nonunit dimension exceeds two.
We claim $\sum_i\rho_i\nu_3(X_i)=1$. By the indicator formula this sum
is $f\sum_B S_{0B}S_{JB}\theta_B^3$, where $J$ is the rational
induction constituent. At $f=4$ its row is supported on twists $1,-1$,
so the sum equals its value at exponent one. At $f=5$ it is supported on
twists of five-power order. Replace three by an exponent equal to three
on the five-part and one on its prime-to-five complement. This exponent
is a unit modulo the twist exponent, and the modular factorization in
Lemma~\ref{lem:odd-row}, evaluated at the twist-one vacuum and $J$, gives
$S_{0J}=1/5$ for the mixed moment.

Now $t=\Re M_3=p\sum_i d_i t_i$ belongs to $K$ and $|g(t)|\le1$ for every
embedding $g$. Indeed modular Galois covariance takes $M_3$ to a vacuum
twist moment, and cyclotomic conjugation commutes with real part.
Inverting the character table and using the unit coefficient in $X_i^3$
gives respectively
\[
 t_i=\sum_g g(d_i)g(t)+\rho_i/f,\qquad
 N_{ii}^i=\sum_g g(p)g(d_i)^3+\rho_i^3/f .
\]
Rotation on $\operatorname{Hom}(1,X_i^{\otimes3})$ shows $t_i\le N_{ii}^i$.
Take $d_i=m$. Put $a=m^2/D<1/4$ and $T=1$ at $f=4$, $T=4/5$ at $f=5$. At $f=5$ all four nonunit rational entries have absolute value one. Column orthogonality gives $\sum_{g\ne1}g(p)g(d_i)^2=1-\rho_i^2/f-a\le T-a$. The pair Fourier estimates in Lemma~\ref{lem:quartic-conjugates} and $p\le P$ give $g(D)\le V_i$ in the corresponding embeddings, so $|g(d_i)|\le B$. The rational terms cancel in $t_i\le N_{ii}^i$ because $\rho_i^3=\rho_i$. Thus weighted Cauchy gives
\[
 t\le a+\frac{B(T-a)+\sqrt{V(T-a)}}{m}.
\]
Indeed the cubic contribution at the three small embeddings is at most $B\sum_{g\ne1}g(p)g(d_i)^2$. Their linear contribution is at most $\sum_{g\ne1}|g(d_i)|\le\sqrt{V(T-a)}$. Replace $m$ in the denominator by $u$. The resulting function is increasing for $0\le a\le1/4$ by the condition on $u$ and $L=T-1/4$. Evaluating at $1/4$ proves the result.
\end{proof}



% ===== quartic15_bound.tex =====
\begin{proposition}\label{prop:four-fifteen-bound}
Suppose the rational codegree is four and $K=\mathbb Q(\zeta_{15})^+$.
Then some nonunit simple has dimension less than $125$.
\end{proposition}
\begin{proof}
Set $p=1/D$, $w_B=S_{0B}$ and $M_j=\sum_Bw_B^2\theta_B^j$. Write
$W=\sum_{\theta_B=1}w_B^2$. Proposition~\ref{prop:quartic-two-support}
gives a lower bound on $W$. On the set $E$ of extra labels of twist one
outside $I(1)$, the rational row $r=S_{J,-}$ has squared mass $1/4$,
since four induction labels each have entry $1/4$. Here
$\operatorname{Tr}_K w_B=-r_B$. Signed modular Galois permutations
preserve $E$ and $r$, hence $4\sum_Ew_Br_B=-\sum_Er_B^2$. Cauchy's
inequality gives $\sum_Ew_B^2\ge1/64$. The constituents of $I(1)$
therefore give
\[
 W\ge 5/64+(3/4-p)^2/3 .
\]

For $g\in\{1,2,4,8\}$ and a root of unity $z$, write $A_{g,j}(z)$
for the average of $\Re z^{ju}$ over units congruent to $\pm g$
modulo fifteen, taken at a common multiple of fifteen and the order
of $z$. Write $C_j(z)$ for four times the average over all units
and $H=-A_{1,1}-A_{2,1}+A_{4,1}+A_{8,1}$. The following pointwise
inequality holds:
\[
\begin{split}
10000\,\mathbf1_{z=1}\le {}&2003+3334 A_{1,1}+2236 A_{1,2}
 +1098 A_{1,3}+44H\\
 &+22 C_1+22 C_2+120 C_3+230 C_4-61 C_6.
\end{split} \tag{*}
\]
For reproducibility, the entire finite check uses only orders $dm$
with $d\mid360$ and $m=1,7$. Indeed, reducing the fraction $z^j=
\exp(2\pi i t/n)$, write $n=ab$, $a=(15,n)$. The coset average
vanishes when $(b,15)>1$. Otherwise it equals
$\mu(b)\cos(2\pi g t b^{-1}/a)/\varphi(b)$, with the cosine interpreted
as one when $a=1$. The all-unit average is $\mu(n)/\varphi(n)$.
Thus if the $2,3,5$-part of the original order does not divide $360$,
all displayed averages vanish. A remaining factor $m$ coprime to
thirty scales all averages, after a common numerator change, by
$\mu(m)/\varphi(m)\in[-1/6,1]$. It suffices to check both endpoints,
and the right endpoint at $z=1$ also implies the version with zero
left side. The exact check in \texttt{verify15.py} isolates
$\cos(2\pi a/15)$ by the real cyclotomic polynomials and verifies (*) at
every reduced numerator.

With respect to the measure $w_B^2$, the expectations of $A_{1,j}$
for $j=1,2,3$ are $\Re M_j$. Indeed $\Re\nu_j(X_i)$ is rational for
these $j$, and modular Galois covariance gives
$\sigma_u(M_j)=M_{j/u}$. The expectation of $H$ is at most
$\sqrt{2/15}$ by the faithful Fourier bound and Cauchy's inequality
on its real and imaginary parts. The expectations of $C_j$ are $3/4$
for $j=1,3$ and one for $j=2,4,6$. To check this, the real parts
of the $j$th indicators are rational for all five values. Column
orthogonality gives
\[
 \operatorname{Tr}_K(\Re M_j)
 =1-\tfrac14\sum_i \rho(X_i)\Re\nu_j(X_i).
\]
The final sum is one for odd $j$ and zero for even $j$: the mixed
moment of $w$ and $r$ has these values after multiplication by four,
since $r$ is supported on the two signs.

Suppose the least nonunit dimension is $m\ge u$. Then $p\le P=(1+4u^2)^{-1}$. The second-indicator bound gives $|M_2|\le (1+2\sqrt{D-1})/D\le(1+4u)P$, since the first expression decreases with $D\ge17$. Substitute this, the third-moment bound of Lemma~\ref{lem:quartic-least-third}, and $p\le P$ into (*). The upper bound falls below $5/64+(3/4-P)^2/3$ at $u=125$. The rational comparison in \texttt{quartic\_generator\_bound.py} verifies this is the first integer $u$ for which these estimates separate.
\end{proof}



% ===== quartic_five_height.tex =====
\begin{proposition}\label{prop:quartic-five-small}
Suppose the rational codegree is five, and put $m=\min_{X\ne1}\operatorname{FPdim}(X)$. Then $m<155$ at conductor fifteen and $m<153$ at conductor twenty.
\end{proposition}
\begin{proof}
Let $c=15$ or $20$ be the conductor of $K$, $G=\operatorname{Gal}(K/\mathbb Q)$, $p=1/D$, $p_g=g(p)$, $s=\sum_g p_g=4/5$, and $d_{i,g}=g(d_i)$. Put $w_B=S_{0B}$ and $M_j=\sum_Bw_B^2\theta_B^j$. By the central indicator formula,
\[
 M_j=p\sum_i d_i\nu_j(X_i),\qquad M_1=p,\qquad |M_2|\le p+2\sqrt p.
\]
Put $t=\Re M_3$.

For a root of unity $z$ let $A_j(z)$ be the mean of $\Re z^{ju}$ over cyclotomic units $u$ that fix $K$, and $C_j(z)$ four times the mean over all units, using an ambient exponent divisible by $c$ and the order of $z$. We use the following pointwise inequality for either conductor:
\begin{align*}
 10^5\bigl({\bf1}_{z=1}+\tfrac15{\bf1}_{\operatorname{ord}(z)=25}\bigr)
 \le{}&20100+32214A_1+21426 A_2+10787 A_3\\
 &+188 C_1+1558 C_2+1039 C_3+1511 C_4-403 C_6 .
\end{align*}
We give the complete finite verification ranges. For $z=\exp(2\pi i a/n)$ in lowest terms, $C_j=4\mu(n/(j,n))/\varphi(n/(j,n))$. To compute $A_j$ average over $u=1\bmod c$: the other kernel coset $u=-1\bmod c$ has the same cosine mean. Inclusion-exclusion at the primes dividing $n$ but not $c$ gives
\[
 A_j=\left(\prod_{q\mid n,\ q\nmid c}(1-1/q)\right)^{-1}
 \sum_{\substack{d\mid\operatorname{rad}(n),\ (d,c)=1\\n\mid jdc}}
 \frac{\mu(d)}d
 \cos\left(\frac{2\pi ja d\bar d}{n}\right),\qquad d\bar d\equiv1\pmod c.
\]
Indeed on the multiples of $d$ the residual progression is constant after applying the exponential precisely when $n\mid jdc$. For $j\in\{1,2,3,4,6\}$ all averages vanish if $32$, $27$ or $125$ divides $n$. Otherwise write $n=n_0 h$ with $n_0\mid 2^43^25^2$ and $(h,30)=1$. Every average acquires the factor $\mu(h)/\varphi(h)$ after changing the numerator modulo $n_0$ by $h^{-1}$. This factor lies in $[-1/6,1]$, so by linearity the two values $h=7,1$ suffice. The indicator on the left can only decrease when a prime-to-thirty factor is present. The exact rational interval check in \texttt{quartic\_five\_twists.py} verifies all divisors and reduced numerators.

Averaging the pointwise inequality against $w_B^2$ replaces $A_1,A_2,A_3$ by $p,M_2,t$. For the exponents of the $C_j$ occurring here, $K\cap\mathbb Q(\zeta_j)=\mathbb Q$, so the cyclotomic mean of $M_j=p\sum_i d_i\nu_j(X_i)$ factors term by term. Since $\operatorname{Tr}_K(pd_i)=\delta_{i,0}-\rho_i/5$ and $\sum_i\rho_i\nu_j(X_i)=1$ for $5\nmid j$ by the modular factorization used in Lemma~\ref{lem:quartic-least-third}, the $w_B^2$-average of every $C_j$ is $4/5$. For $m\ge u$ put $P=(1+4u^2)^{-1}$. Substitute $p\le P$, $|M_2|\le(1+4u)P$ and Lemma~\ref{lem:quartic-least-third} in the averaged inequality. Lemma~\ref{lem:quartic-five-row-masses} gives the lower bound
\[
 W(1)+W(25)/5\ge\sum_g p_g^2+1/20\ge (s-P)^2/3+1/20.
\]
The upper bound is strictly smaller at $u=155$ for $c=15$ and at $u=153$ for $c=20$. The verifier \texttt{quartic\_generator\_bound.py} computes these first integer thresholds from the inequalities.
\end{proof}



% ===== biquad_sieve.tex =====
\begin{proposition}\label{prop:quartic-biquadratic-sieve}
None of the six biquadratic fields with complementary codegree four
occurs for a perfect rank-five fusion category in the dimension range
of Proposition~\ref{prop:quartic-four-effective}.
\end{proposition}
\begin{proof}
Write $G=\{1,h_1,h_2,h_3\}$ and let $c_i$ be the fundamental
discriminant of the quadratic field fixed by $h_i$. Put
$s=3/4$, $D_i=h_iD$ and $x_i=D/D_i$. The quadratic
dimension-block bounds give
\[
 1<D_i<11,
 \qquad
 |2s(1+x_i)/P-s|\le1/\sqrt{c_i},\qquad P=1+x_1+x_2+x_3=sD .
\]
The first bound is Lemma~\ref{lem:quartic-conjugates}.
The determinant identity makes $\mathrm N(D)$ a rational square since
both $f=4$ and the biquadratic discriminant $c_1c_2c_3$ are squares.
Consequently $z_i=DD_i/\sqrt{\mathrm N(D)}$ belongs to the quadratic
fixed field, where it is totally positive of norm one. It is an
algebraic unit because $z_i^2=(D/D_j)(D_i/D_k)$ is a product of
$d$-number ratios. We have
\[
 x_i=z_jz_k,\qquad D/121<z_i^2<11D .
\]
Write $2z_i=t_i+b_i\sqrt{c_i}$, so $t_i^2-c_i b_i^2=4$.
For a field with dimension bound $H$, enumerate the positive powers of the smallest norm-one unit through trace $\lfloor\sqrt{11H}\rfloor+2$. Include also the unit one and the inverses of powers of trace at most four: if $z_i<1$, then $1/z_i<11/\sqrt{17}$ and $t_i=z_i+1/z_i<4$.

Here is the exact divisibility test. Set
$C=t_1t_2t_3-b_1b_2b_3\sqrt{c_1c_2c_3}$ and
$T=\operatorname{Tr}(P)=4+\sum_{i<j}t_it_j$. With $x_{\mathrm{id}}=1$, define
$Q_i=\sum_{g\in G} x_{g h_i}^2/x_g$, $x_{h_i}=z_jz_k$.
The unit coefficient of $F(A_{h_i})^2$ gives the rational integer
$(f-1)Q_i/P$. Therefore $2T$ divides $(f-1)$ times
\[
 2\operatorname{Tr}(Q_i)
 =2\big((t_j^2-2)(t_k^2-2)+t_jt_k\big)
     +(t_j+t_k)(C-2t_i).
\]
This follows by expanding
$Q_i=z_j^2z_k^2+(z_jz_k)^{-1}
       +z_i z_j^2/z_k+z_i z_k^2/z_j$ and taking traces.
On the triples passing the three divisions, enclose $z_i$ by solving $z_i+1/z_i=t_i$ with its indicated sign and enclose $D=P/s$ by the same interval arithmetic. The condition $D>17$ and the three displayed Fourier inequalities exclude every triple. The complete calculation appears in \texttt{biquad\_sieve.py}.
\end{proof}



% ===== cyclic_dimension.tex =====
\begin{proposition}\label{prop:cyclic-bounded-dim}
At rational codegree four, none of the cyclic fields of conductors
$16,20,35,40$ occurs for a perfect self-dual rank-five fusion category
in the dimension range of Proposition~\ref{prop:quartic-four-effective}.
\end{proposition}
\begin{proof}
Write $c$ for the conductor, $k=\mathbb Q(\sqrt d)$ for the quadratic subfield,
$\delta=\operatorname{disc}(k)$, and $s=3/4$. Here $d=2$, $\delta=8$
at $c=16$, and $d=\delta=5$ otherwise. The conductor--discriminant formula
gives $\operatorname{disc}(K)=\delta c^2$. Let $h=[\mathcal O_K:B]$ for
the dimension order $B$ and set $N=\mathrm N(D)=4\delta c^2 h^2$.
Fix a generator $\sigma$ of the Galois group and put
\[
 R=D\sigma^2(D),\quad U=D+\sigma^2(D),\quad z=R/\sqrt N .
\]
Use the positive square root at the distinguished embedding. Both $R$ and $z$
belong to $k$. The square $z^2=R/\sigma(R)$ is a product of two
$d$-number ratios. Consequently $z$ is an algebraic unit in $k$,
of norm $-1$ because $\sqrt N$ is a rational multiple of $\sqrt d$.
Write $H$ for the indicated dimension bound. The other three conjugates
are between one and eleven by Lemma~\ref{lem:quartic-conjugates}, so
\[
 h^2<11^3 H/(4\delta c^2),\quad \sqrt D/11<z<\sqrt{11H},\quad
 1<\sigma(R)<121,\quad 2<\sigma(U)<22 .
\]
Every such $z$ is an odd power of $(1+\sqrt5)/2$ or $1+\sqrt2$, as follows from the quadratic norm equation. Include negative powers as well: if $z<1$ then $1/z<11/\sqrt{17}$. These bounds enumerate $R$ by integer recurrences.

For each $(z,h)$ write $2R=h(C_0+E_0\sqrt d)$ and $2U=A+B_1\sqrt d$;
the coefficients are integers. The trace identity $U/R+\sigma(U/R)=s$
is the linear equation
\[
        A C_0-d B_1 E_0=2sN/h.
\]
It suffices to enumerate $21\le A\le H+33$ in this single congruence
class after division by $\gcd(C_0,dE_0)$. The real inequalities above
give $4<A-B_1\sqrt d<44$. All four conjugates of $D$ are roots of
\[
 P(X)=X^4-A X^3+\bigl(\mathrm N_{k/\mathbb Q}(U)+h C_0\bigr)X^2-sN X+N.
\]
Require integral coefficients, $N\mid A^4$ and
$N\mid(\mathrm N_{k/\mathbb Q}(U)+h C_0)^2$ by the $d$-number property.
Also $d\,\mathrm N_{k/\mathbb Q}(U^2-4R)$ is a rational square:
$(D-\sigma^2 D)(\sigma D-\sigma^3D)$ is anti-invariant under $\sigma$
and lies in $\sqrt d\,\mathbb Q$. Finally put $p=1/D$. The squared quadratic Fourier coefficient is
$s^2-4\mathrm N_{k/\mathbb Q}(U)/N$, since one pair of reciprocal
conjugates has sum $U/R$. By Parseval, twice the squared faithful
coefficient is $4\operatorname{Tr}(p^2)-2s^2+4\mathrm N_{k/\mathbb Q}(U)/N$.
It must be at most $2/c$. The integer enumeration in
\texttt{cyclic\_dimension.py} gives zero tuples for each of the four ranges.
\end{proof}



% ===== small_simple.tex =====
\begin{proposition}\label{prop:quartic-small-simple}
In the cyclic quartic cases $(F,f)=(15,4),(15,5),(20,5)$, where $F$ is the conductor, the generator ranges $2<d\le125,155,153$, respectively, are impossible for the least nonunit dimension.
\end{proposition}
\begin{proof}
Let $d$ be that least dimension, let $X$ be the corresponding simple and put $n=N_{XX}^X$, $r=\rho(X)$. Each nonunit simple generates the category. Thus the fusion matrix of $X$ is irreducible and its Perron eigenvalue $d$ is simple. Solving for its eigenvector normalized to have unit coordinate one puts every simple dimension in $\mathbb Q(d)$, so $d$ generates $K$. The nonidentity embeddings $g$ satisfy $1<g(D)<11$, hence $|g(d)|<\sqrt{10}$ by the unit term in $g(D)$.

Write $P(T)=T^4+a_3T^3+a_2T^2+a_1T+a_0$ for the minimal polynomial of $d$, and $p=1/D$. Orthogonality and the integer $n$ give
\[
 f\operatorname{Tr}(p d^j)=b_j,\quad
 (b_0,b_1,b_2,b_3)=(f-1,-r,f-r^2,fn-r^3).
\]
Multiplying $\sum_g f g(p)/(T-g(d))$ by $P(T)$ and expanding at infinity yields $fp P'(d)=H(d)$, for
\[
 H(T)=b_0T^3+(a_3b_0+b_1)T^2
 +(a_2b_0+a_3b_1+b_2)T+a_1b_0+a_2b_1+a_3b_2+b_3.
\]
Let $B$ be the dimension order. Taking norms and using $\operatorname N(D)=f\operatorname{disc}(B)$ gives
\[
 \operatorname N(H(d))=f^3[B:\mathbb Z[d]]^2.
\]
This positive integer divides $f^3\operatorname{disc}(P)/\operatorname{disc}(K)$. Moreover
$n\le d/4+\sqrt{10}+1/f<d/4+4$, since $D>4d^2$ and $\sum_{g\ne1}g(p)g(d)^2\le1$.

The enumeration uses $t=2\cos(2\pi/F)$, with $\mathcal O_K=\mathbb Z[t]$. Its minimal polynomial $P_t$ is $T^4-T^3-4T^2+4T+1$ for $F=15$ and $T^4-5T^2+5$ for $F=20$. We may take the largest root $t_0$ at the distinguished embedding by Galois conjugation. Write $d=A+Bt+Ct^2+Et^3$ here with integer coefficients (the letter $B$ in this expression denotes a coefficient). For completeness, let $z_i$ be the four conjugates of $d$, with $z_0=d$. Lagrange interpolation gives $E=\sum_i z_i/P_t'(t_i)$. The second divided difference of the three small $z_i$ equals $C+(t_1+t_2+t_3)E$. Their first divided difference on $t_1,t_3$ equals $B+(t_1+t_3)C+(t_1^2+t_1t_3+t_3^2)E$. Finally the bound on $z_2$ bounds $A$. These four inequalities give nested finite intervals for the coefficients, which are rounded outward in \texttt{small\_simple.py} using rational isolating intervals of length $10^{-12}$. The count includes the few points retained by outward rounding. For every candidate the program tests each allowed $r$ and each integer $0\le n<d/4+4$, the norm conditions above, integrality and the $d$-number condition for $D=fP'(d)/H(d)$, and $D\ge1+4d^2$ and $1<g(D)<11$ at the other embeddings.
\[
\begin{array}{c|r|r|r|r}
(F,f)&d\text{ tested}&(d,r,n)\text{ after norm tests}
 &\text{after dimension tests}&\text{after Fourier test}\\\hline
(15,4)&931&20&2&0\\
(15,5)&1158&3&0&0\\
(20,5)&856&23&3&0
\end{array}
\]
For the last column let $\sigma$ generate $G$. At $(15,4)$ the two remaining tuples have $(p-\sigma^2p)^2+(\sigma p-\sigma^3p)^2=29/240$. At $(20,5)$ the three values are $9/80,9/80,1/8$. In each case the value exceeds the faithful Fourier bound $1/F$. All tests use integer or rational arithmetic.
\end{proof}



% ===== remaining.tex =====
\section{The sharp bounds and the finite classification}\label{sec:remaining}
\begin{theorem}\label{thm:sharp-final}
At rank five, $\mu\le4$ and $D\le48+24\sqrt3$. Both inequalities are sharp.
\end{theorem}
\begin{proof}
Use the positive spherical conjugate supplied by Theorem~A. Proposition~\ref{prop:weak-five} proves both upper bounds for weakly integral categories. Assume $D$ is irrational. Four invertibles give the near-group rings of Corollary~\ref{cor:four-final}, for which the inequalities follow from $x^2=4+mx$ and $m\le4$. Three invertibles and a fixed involution acting trivially on all other simples give integral categories by Theorems~\ref{thm:three-inv} and~\ref{thm:fixed-inv-two-classification}. For an involution exchanging two simples apply Proposition~\ref{prop:sharp-exchange}.

It remains to treat a perfect category. Two dual pairs force integrality by Proposition~\ref{prop:rank-five-commutative}. A simple of dimension at most two gives the adjoint $\mathrm{SU}(2)_9$ ring by Theorem~\ref{thm:perfect-small-simple}; the same ring is the only possibility in quintic degree by Theorem~\ref{thm:quintic-classification}. Its fusion rules are multiplicity-free and $D=11/(4\sin^2(\pi/11))<35$, for example by $\sin t\ge t-t^3/6$ at $t=\pi/11$. Otherwise every nonunit dimension exceeds two, hence $D>17$.

In degree two, the complementary codegrees and their Fourier estimates give the two exclusions in Propositions~\ref{prop:mixed-polar} and~\ref{prop:quadratic-455-end}. In degree three Theorem~\ref{thm:cubic-intrinsic} gives $D<36$.
Since the other three nonunit dimensions exceed two, every nonunit simple
satisfies $d_i^2<D-13<23$. Thus $d_i<5$, and
Lemma~\ref{lem:mu-d} gives the integer bound $\mu\le4$.
The final enumeration identifies the surviving cubic ring. In degree four the four dimension characters are real and the remaining character is rational, so all simples are self-dual. Corollary~\ref{cor:quartic-coupled} gives the possible field and rational codegree. Proposition~\ref{prop:quartic-sixteen} excludes codegree sixteen. At codegree four the moment bounds and arithmetic sieves in Propositions~\ref{prop:quartic-four-effective}, \ref{prop:quartic-biquadratic-sieve} and~\ref{prop:cyclic-bounded-dim} exclude all biquadratic fields and the cyclic conductors $16,20,35,40$. For the remaining cyclic pairs the least nonunit dimension bound and the generator calculation in Section~\ref{sec:quartic-exhaustion} give a contradiction.

For the multiplicity endpoint use the unitary near-group
realizations $C_4+4$ and $C_2^2+4$ \cite[Examples~9.5--9.6]{Izumi}.
For the dimension endpoint, the following construction also verifies
the tensor used in the model comparison. Take either of Izumi's unitary near-group
categories of type $C_3+6$ \cite[Theorem~10.19]{Izumi}. In the notation of
\cite[Lemma~13.1 and Example~13.10]{Izumi}, the symmetry group contains
\[
 (\theta,u)=\bigl(-1,R(\pi/4)L\bigr),\qquad L=\operatorname{diag}(1,-1).
\]
This is an involution: $\theta$ inverts $C_3$ and $u$ is a reflection.
Its equivariantization is unitary, with simple objects $1,g,Z,X,Y$
of dimensions $1,1,2,d,d$, where $d=3+2\sqrt3$.
Here the orbit--stabilizer description of simples and the corresponding
tensor-product formula are given in
\cite[Corollary~2.13 and Theorem~3.9]{BurciuNatale}.
The intertwiner $S_0\in\operatorname{Hom}(1,\rho^2)$ is fixed, so
$X,Y$ are self-dual. On
$\operatorname{Hom}(\rho,\rho^2)=\mathbb C[C_3]\otimes K_0$
the action has trace zero, since the reflection on $K_0$ does.
Its two eigenspaces therefore both have dimension three. This gives
\[
 X^2=Y^2=1+Z+3X+3Y,\qquad XY=g+Z+3X+3Y,
\]
with $gX=Y$, $gZ=Z$, $Z^2=1+g+Z$ and $ZX=ZY=X+Y$.
Thus $\mu=3$ and $D=6+2d^2=48+24\sqrt3$.
\end{proof}
\begin{theorem}\label{thm:frontier-final}
Exactly the seventeen based rings in Table~\ref{tab:models-five} admit a complex fusion categorification.
\end{theorem}
\begin{proof}
Apply Theorem~\ref{thm:sharp-final} and enumerate the based rings in its range. The unit, duality, Frobenius reciprocity and associativity equations give 269 commutative based rings of multiplicity at most four up to permutation of the nonunit basis elements. The three duality types exhaust rank five. The program \texttt{final\_census/verify.py} regenerates these rings, then retains exactly those with $D\le M=48+24\sqrt3$. For the dimension test let $C_R=\sum_i N_iN_i^{\mathsf t}$. Its largest eigenvalue is $D$ by the common positive dimension vector. All eigenvalues are real, so $D\le M$ if and only if every coefficient of $\det((t+M)I-C_R)$ is nonnegative. The test uses exact signs in $\mathbb Q(\sqrt3)$.

\begin{center}
\begin{tabular}{rrrr}
\toprule
$\mu$ & Within both bounds & Excluded & Realized\\\midrule
1&16&6&10\\2&30&27&3\\3&46&44&2\\4&22&20&2\\\midrule
Total&114&97&17\\\bottomrule
\end{tabular}
\end{center}
For the 114 retained tensors, Theorem~A reduces categorifiability
to the existence of a pseudo-unitary realization. The tests below
exclude 97 tensors:
\[
\begin{array}{lr}
\text{formal-codegree $d$-number test}&76\\
\text{subcategory dimension divisibility}&2\\
\text{central dimension norms or fields}&5\\
\text{unit-induction and second-moment constraints}&10\\
\text{higher-indicator trace constraints}&4.
\end{array}
\]
Each of the seventeen other tensors is compared with its model tensor
under all permutations of the four nonunit basis elements. The
comparison gives an explicit based-ring isomorphism. The models are
those of Table~\ref{tab:models-five}, with unitary realizations supplied
by the constructions cited after that table; the endpoint construction
above specifies the equivariantization. The command
\texttt{python3 final\_census/verify.py} regenerates the enumeration,
checks every exclusion and verifies all model identifications.
This proves Theorem~B.
\end{proof}

\subsection{Mathematical content of the finite certificates}
We record the necessary conditions checked in the last step, so that
its dependence on positive spherical dimensions is explicit.
The $d$-number test is \cite[Theorem~1.2 and Lemma~2.7]{OstrikCodegrees}; subcategory divisibility
is \cite[Proposition~8.15]{ENO}. Central induction requires every
codegree to lie in the simple-dimension field and every $D/f$ to be
an algebraic integer. These give the first three tests.

For the remaining tests let $C_R=\sum_i N_iN_i^{\mathsf t}$ and
$d=(d_i)_i$. A central simple $Y$ has restriction column
$v_Y=([F(Y):X_i])_i$, and induction--restriction
\cite[Proposition~5.4]{ENO} gives
\[
 C_R=\sum_Y v_Yv_Y^{\mathsf t},\qquad \dim Y=v_Y\cdot d.
\]
The five constituents of $I(1)$ have $v_{Y,0}=1$, twist one and
dimensions $D/f$. Their possible columns are enumerated from these
conditions and the bounds $v_{Y,i}^2\le(C_R)_{ii}$ and
$v_{Y,i}\le(C_R)_{0i}$. If $L$ is the matrix of these five columns,
the residual Gram matrix $C_R-LL^{\mathsf t}$ must be entrywise
nonnegative and positive semidefinite.

For a nonunit simple $X_i$, put $a_i=(LL^{\mathsf t}d)_i$, the
known twist-one dimension in $I(X_i)$. On removing that mass, the total, first and second moments
are $Dd_i-a_i$, $-a_i$ and $D\nu_2(X_i)-a_i$.
The inequality $\mathbb E(\cos^2 t)\ge(\mathbb E\cos t)^2$ gives
\[
 a_i\le
 \begin{cases}
 Dd_i/3,&X_i\not\simeq X_i^*,\\
 Dd_i(d_i+1)/(3d_i+1),&X_i\simeq X_i^*.
 \end{cases}
\]
Every possible $L$ in the ten corresponding certificates violates a
Gram or moment condition.

In the four remaining certificates, the residual Gram matrix bounds
all possible further restriction columns. For each proposed central
dimension $z=v\cdot d$, the verifier also requires $D/z$ to be an
algebraic integer and $z$ to be a $d$-number. Indeed $(D/z)^2$ is a
modular formal codegree, and $D$ is itself a $d$-number.
For a non-self-dual simple $X_i$, one has $\nu_2(X_i)=0$.
For a self-dual simple $X_i$, the second indicator is $+1$ whenever
$X_i$ occurs with odd multiplicity in some $X_j\otimes X_j^*$:
the latter has its canonical symmetric self-duality, so a self-dual
summand of indicator $-1$ must occur evenly. In the other self-dual
cases both signs are retained by the verifier. The central trace identities
\cite[Theorem~4.1]{NSExp}
\[
 \sum_Y v_{Y,i}\dim(Y)\theta_Y^j=D\nu_j(X_i)
\]
are averaged in the rational or quadratic dimension field. Categorical
Cauchy \cite[Theorem~3.9]{RankFinite} restricts the primes in the twist
orders. For a fixed finite list of powers $j$, sufficiently high
primary orders have zero averages; the finite orders and this zero
case give all possible moment columns. For $j>1$ with $K\cap\mathbb Q(\zeta_j)=\mathbb Q$, let $H_{ij}=[X_i^{\otimes j}:1]$ and let $p(j)$ be
the least prime divisor of $j$. The pivotal rotation on
$\operatorname{Hom}(1,X_i^{\otimes j})$ has order dividing $j$
\cite[Theorem~5.1]{NSPiv}, and gives the exact averaged bounds
$-H_{ij}/(p(j)-1)\le\operatorname{Av}\nu_j(X_i)\le H_{ij}$,
by the Ramanujan averages of $j$-th roots of unity.
Each stored integer separating form is nonnegative on every admissible
column and strictly negative on the required moment vector.
All column lists, indicator assignments and pairings are recomputed
exactly. No floating-point feasibility decision is used.



% ===== refs.tex =====
\begin{thebibliography}{99}
\raggedright

\bibitem{ABDP} M. A. Alekseyev, W. Bruns, J. Dong and S. Palcoux,
\emph{Classifying integral Grothendieck rings up to rank 5 and beyond},
\emph{Lett. Math. Phys.} \textbf{116} (2026), article~141.
\href{https://doi.org/10.1007/s11005-026-02164-7}{doi:10.1007/s11005-026-02164-7};
\href{https://arxiv.org/abs/2507.07023v2}{arXiv:2507.07023v2}.

\bibitem{AndersenParadowski}
H.~H. Andersen and J. Paradowski,
\emph{Fusion categories arising from semisimple Lie algebras},
Comm. Math. Phys. \textbf{169} (1995), no.~3, 563--588.
\href{https://doi.org/10.1007/BF02099312}{doi:10.1007/BF02099312}.

\bibitem{BW} J. W. Barrett and B. W. Westbury,
\emph{Spherical categories}, Adv. Math. \textbf{143}
(1999), no.~2, 357--375.
\href{https://doi.org/10.1006/aima.1998.1800}{doi:10.1006/aima.1998.1800};
\href{https://arxiv.org/abs/hep-th/9310164}{arXiv:hep-th/9310164}.

\bibitem{BO5} P. Bruillard and C. M. Ortiz-Marrero,
\emph{Classification of rank 5 premodular categories},
J. Math. Phys. \textbf{59} (2018), no.~1, article~011702.
\href{https://doi.org/10.1063/1.5020256}{doi:10.1063/1.5020256};
\href{https://arxiv.org/abs/1612.08769}{arXiv:1612.08769}.

\bibitem{Bruillard4} P. Bruillard,
\emph{Rank 4 premodular categories}, with an appendix written jointly with
C. Galindo, S.-H. Ng, J. Plavnik, E. Rowell and Z. Wang,
New York J. Math. \textbf{22} (2016), 775--800.
\href{https://nyjm.albany.edu/j/2016/22-36.html}{Journal article}.

\bibitem{BNRW} P. Bruillard, S.-H. Ng, E. C. Rowell and Z. Wang,
\emph{On classification of modular categories by rank}, Int. Math. Res. Not. (2016), 7546--7588.
\href{https://arxiv.org/abs/1507.05139}{arXiv:1507.05139}.

\bibitem{RankFinite} P. Bruillard, S.-H. Ng, E. C. Rowell and Z. Wang, \emph{Rank-finiteness for modular categories}, J. Amer. Math. Soc. \textbf{29} (2016), 857--881. \href{https://arxiv.org/abs/1310.7050}{arXiv:1310.7050}.

\bibitem{BV} A. Bruguières and A. Virelizier,
\emph{On the center of fusion categories}, Pacific J. Math.
\textbf{264} (2013), 1--30. \href{https://www.math.virelizier.com/sources/articles/Center-fusion.pdf}{Full text}.

\bibitem{Dyadic} S. Burciu, J. Dong and S. Palcoux (contributors), \emph{Prime factors of fusion-category dimensions: weak group-theoreticity, rank and decomposition}, Section~3, \href{https://hexagonmath.org/2611.00233v3}{hexagon:2611.00233v3}.

\bibitem{BurciuNatale}
S. Burciu and S. Natale,
\emph{Fusion rules of equivariantizations of fusion categories},
J. Math. Phys. \textbf{54} (2013), no.~1, article~013511.
\href{https://doi.org/10.1063/1.4774293}{doi:10.1063/1.4774293};
\href{https://arxiv.org/abs/1206.6625v2}{arXiv:1206.6625v2}.
The cited statements refer to the corrected 2019 arXiv version.

\bibitem{DLN}
C. Dong, X. Lin, and S.-H. Ng, \emph{Congruence property in conformal
field theory}, Algebra Number Theory \textbf{9} (2015), no.~9, 2121--2166.
\href{https://doi.org/10.2140/ant.2015.9.2121}{doi:10.2140/ant.2015.9.2121}.

\bibitem{DCW} J. Dong, G. Chen and Z. Wang,
\emph{Fusion categories containing a fusion subcategory with maximal rank},
J. Algebra \textbf{604} (2022), 107--127.
\href{https://doi.org/10.1016/j.jalgebra.2022.02.020}{doi:10.1016/j.jalgebra.2022.02.020};
\href{https://arxiv.org/abs/2109.08787}{arXiv:2109.08787}.

\bibitem{DPP} J. Dong, S. Palcoux and A. Plessis,
\emph{Classifying Grothendieck rings of pivotal fusion categories of rank four},
preprint (2026).
\href{https://arxiv.org/abs/2609.33272v2}{arXiv:2609.33272v2}.

\bibitem{EIP} C. Edie-Michell, M. Izumi and D. Penneys,
\emph{Classification of $\mathbb Z/2\mathbb Z$-quadratic unitary fusion categories},
with an appendix by R. Johnson, S.-H. Ng, D. Penneys, J. Roat,
M. Titsworth and H. Tucker,
Quantum Topol. \textbf{16} (2025), no.~3, 493--568.
\href{https://doi.org/10.4171/QT/201}{doi:10.4171/QT/201}.

\bibitem{EGO} P. Etingof, S. Gelaki and V. Ostrik, \emph{Classification of fusion categories of dimension $pq$}, Int. Math. Res. Not. (2004), 3041--3056. \href{https://arxiv.org/abs/math/0304194}{arXiv:math/0304194}.

\bibitem{EGNO} P. Etingof, S. Gelaki, D. Nikshych and V. Ostrik,
\emph{Tensor Categories}, Mathematical Surveys and Monographs
\textbf{205}, American Mathematical Society, Providence, RI, 2015.
\href{https://doi.org/10.1090/surv/205}{doi:10.1090/surv/205}.

\bibitem{ENO}
P. Etingof, D. Nikshych, and V. Ostrik, \emph{On fusion categories},
Ann. of Math. (2) \textbf{162} (2005), no.~2, 581--642.
\href{https://doi.org/10.4007/annals.2005.162.581}{doi:10.4007/annals.2005.162.581}.

\bibitem{EvansGannon}
D.~E. Evans and T. Gannon,
\emph{Near-group fusion categories and their doubles},
Adv. Math. \textbf{255} (2014), 586--640.
\href{https://doi.org/10.1016/j.aim.2013.12.014}{doi:10.1016/j.aim.2013.12.014};
\href{https://arxiv.org/abs/1208.1500}{arXiv:1208.1500}.

\bibitem{GannonSchopieray}
T.~Gannon and A.~Schopieray,
\emph{Algebraic number fields generated by dimensions in fusion rings},
Commun. Number Theory Phys. \textbf{18} (2024), 705--743.
\href{https://doi.org/10.4310/CNTP.241202233757}{doi:10.4310/CNTP.241202233757};
\href{https://arxiv.org/abs/1912.12260}{arXiv:1912.12260}.

\bibitem{GN} S. Gelaki and D. Nikshych,
\emph{Nilpotent fusion categories}, Adv. Math. \textbf{217}
(2008), no.~3, 1053--1071.
\href{https://doi.org/10.1016/j.aim.2007.08.001}{doi:10.1016/j.aim.2007.08.001};
\href{https://arxiv.org/abs/math/0610726}{arXiv:math/0610726}.

\bibitem{HH} T. J. Hagge and S.-M. Hong,
\emph{Some non-braided fusion categories of rank three},
Commun. Contemp. Math. \textbf{11} (2009), no.~4, 615--637.
\href{https://doi.org/10.1142/S0219199709003521}{doi:10.1142/S0219199709003521};
\href{https://arxiv.org/abs/0704.0208v2}{arXiv:0704.0208v2}.

\bibitem{HR} S.-M. Hong and E. C. Rowell,
\emph{On the classification of the Grothendieck rings of non-self-dual modular categories},
J. Algebra \textbf{324} (2010), 1000--1015.
\href{https://doi.org/10.1016/j.jalgebra.2009.11.044}{doi:10.1016/j.jalgebra.2009.11.044}.

\bibitem{Izumi} M. Izumi,
\emph{A Cuntz algebra approach to the classification of near-group categories},
Proc. Centre Math. Appl. Austral. Nat. Univ. \textbf{46} (2017), 222--343.
\href{https://arxiv.org/abs/1512.04288v2}{arXiv:1512.04288v2}.

\bibitem{Larson} H. K. Larson,
\emph{Pseudo-unitary non-self-dual fusion categories of rank 4},
J. Algebra \textbf{415} (2014), 184--213.
\href{https://doi.org/10.1016/j.jalgebra.2014.05.032}{doi:10.1016/j.jalgebra.2014.05.032};
\href{https://arxiv.org/abs/1401.1879v2}{arXiv:1401.1879v2}.

\bibitem{LiuYB}
Z. Liu, \emph{Yang--Baxter relation planar algebras},
\href{https://arxiv.org/abs/1507.06030v3}{arXiv:1507.06030v3}
(2015; revised 2016).

\bibitem{LPR}
Z.~Liu, S.~Palcoux and Y.~Ren,
\emph{Classification of Grothendieck rings of complex fusion categories
of multiplicity one up to rank six},
Lett. Math. Phys. \textbf{112} (2022), article~54.
\href{https://arxiv.org/abs/2010.10264v5}{arXiv:2010.10264v5}.

\bibitem{MengYu} Y. Meng and Z. Yu,
\emph{The multiplicity bound for near-group fusion categories of type
$\mathbb Z/4\mathbb Z+k$},
\href{https://arxiv.org/abs/2611.12191}{arXiv:2611.12191} (2026).

\bibitem{MugerCenter} M. M\"uger,
\emph{From subfactors to categories and topology II:
The quantum double of tensor categories and subfactors},
J. Pure Appl. Algebra \textbf{180} (2003), no.~1--2, 159--219.
\href{https://doi.org/10.1016/S0022-4049(02)00248-7}{doi:10.1016/S0022-4049(02)00248-7};
\href{https://arxiv.org/abs/math/0111205}{arXiv:math/0111205}.

\bibitem{NSPiv} S.-H. Ng and P. Schauenburg,
\emph{Higher Frobenius--Schur indicators for pivotal categories},
in \emph{Hopf algebras and generalizations},
Contemp. Math. \textbf{441}, Amer. Math. Soc., Providence, RI,
2007, 63--90.
\href{https://doi.org/10.1090/conm/441/08500}{doi:10.1090/conm/441/08500};
\href{https://arxiv.org/abs/math/0503167}{arXiv:math/0503167}.

\bibitem{NSExp} S.-H. Ng and P. Schauenburg,
\emph{Frobenius--Schur indicators and exponents of spherical categories},
Adv. Math. \textbf{211} (2007), no.~1, 34--71.
\href{https://doi.org/10.1016/j.aim.2006.07.017}{doi:10.1016/j.aim.2006.07.017};
\href{https://arxiv.org/abs/math/0601012}{arXiv:math/0601012}.

\bibitem{NS} S.-H. Ng and P. Schauenburg, \emph{Congruence subgroups and generalized Frobenius--Schur indicators}, Comm. Math. Phys. \textbf{300} (2010), 1--46. \href{https://arxiv.org/abs/0806.2493}{arXiv:0806.2493}.

\bibitem{NSW} S.-H. Ng, A. Schopieray and Y. Wang,
\emph{Higher Gauss sums of modular categories}, Selecta Math. (N.S.)
\textbf{25} (2019), no.~4, article~53.
\href{https://doi.org/10.1007/s00029-019-0499-2}{doi:10.1007/s00029-019-0499-2};
\href{https://arxiv.org/abs/1812.11234v2}{arXiv:1812.11234v2}.
The theorem and equation numbers used here refer to this arXiv version.

\bibitem{Ostrik2}
V. Ostrik, \emph{Fusion categories of rank 2}, with an appendix by P.~Etingof,
Math. Res. Lett. \textbf{10} (2003), no.~2--3, 177--183.
\href{https://arxiv.org/abs/math/0203255}{arXiv:math/0203255}.

\bibitem{OstrikPremod3} V. Ostrik,
\emph{Pre-modular categories of rank 3},
Mosc. Math. J. \textbf{8} (2008), no.~1, 111--118.
\href{https://doi.org/10.17323/1609-4514-2008-8-1-111-118}{doi:10.17323/1609-4514-2008-8-1-111-118};
\href{https://arxiv.org/abs/math/0503564}{arXiv:math/0503564}.

\bibitem{OstrikCodegrees}
V. Ostrik, \emph{On formal codegrees of fusion categories},
Math. Res. Lett. \textbf{16} (2009), no.~5, 895--901.
\href{https://doi.org/10.4310/MRL.2009.v16.n5.a11}{doi:10.4310/MRL.2009.v16.n5.a11}.

\bibitem{Ostrik3}
V. Ostrik, \emph{Pivotal fusion categories of rank 3}, with an appendix
written jointly with D. Nikshych, Mosc. Math. J. \textbf{15} (2015),
no.~2, 373--396.
\href{https://doi.org/10.17323/1609-4514-2015-15-2-373-396}{doi:10.17323/1609-4514-2015-15-2-373-396}.

\bibitem{Rank2} \emph{Grothendieck rings of complex fusion categories: rank 2}, contributed by S. Palcoux, Hexagon (2026). \href{https://hexagonmath.org/2611.00641v1}{hexagon:2611.00641v1}.

\bibitem{Rank3} \emph{Grothendieck rings of complex fusion categories: rank 3}, contributed by S. Palcoux, Hexagon (2026). \href{https://hexagonmath.org/2611.00642v1}{hexagon:2611.00642v1}.

\bibitem{Rank4} \emph{Grothendieck rings of complex fusion categories: rank 4}, contributed by S. Palcoux, Hexagon (2026). \href{https://hexagonmath.org/2611.00645v1}{hexagon:2611.00645v1}.

\bibitem{Rank5Bounded} \emph{Rank-five Grothendieck rings through multiplicity twenty}, research manuscript with computational certificates, 17 October 2026, main theorem and model table. \href{https://github.com/sebastienpalcoux/grothendieck-ring-classification/tree/d35cb684d2f27cca18f89ab24c769af756c16160/manuscripts/rank5}{Manuscript and data}.

\bibitem{RSW} E. Rowell, R. Stong and Z. Wang,
\emph{On classification of modular tensor categories},
Commun. Math. Phys. \textbf{292} (2009), no.~2, 343--389.
\href{https://doi.org/10.1007/s00220-009-0908-z}{doi:10.1007/s00220-009-0908-z};
\href{https://arxiv.org/abs/0712.1377}{arXiv:0712.1377}.

\bibitem{SchCenter} P. Schauenburg,
\emph{The monoidal center construction and bimodules},
J. Pure Appl. Algebra \textbf{158} (2001), no.~2--3, 325--346.
\href{https://doi.org/10.1016/S0022-4049(00)00040-2}{doi:10.1016/S0022-4049(00)00040-2}.

\bibitem{SchNG} A. Schopieray, \emph{Categorification of integral group rings extended by one dimension}, J. Lond. Math. Soc. (2) \textbf{108} (2023), no.~4, 1617--1641. \href{https://doi.org/10.1112/jlms.12789}{doi:10.1112/jlms.12789}.

\bibitem{SchFixed} A. Schopieray, \emph{Fixed-point-free fusion automorphisms}, \href{https://arxiv.org/abs/2306.01666}{arXiv:2306.01666}.

\bibitem{SchopierayBlock} A. Schopieray, \emph{The dimension block of a spherical fusion category}, \href{https://arxiv.org/abs/2609.35287}{arXiv:2609.35287}.

\bibitem{SWY} A. Schopieray, Y. Wang and Z. Yu, \emph{Stretched near-groups}, \href{https://arxiv.org/abs/2611.11240}{arXiv:2611.11240}.

\bibitem{Siehler}
J. Siehler, \emph{Near-group categories},
Algebr. Geom. Topol. \textbf{3} (2003), 719--775.
\href{https://doi.org/10.2140/agt.2003.3.719}{doi:10.2140/agt.2003.3.719};
\href{https://arxiv.org/abs/math/0209073}{arXiv:math/0209073}.

\bibitem{TambaraYamagami}
D. Tambara and S. Yamagami,
\emph{Tensor categories with fusion rules of self-duality for finite
abelian groups}, J. Algebra \textbf{209} (1998), no.~2, 692--707.
\href{https://doi.org/10.1006/jabr.1998.7558}{doi:10.1006/jabr.1998.7558}.

\bibitem{WD} K. Wang and J. Dong,
\emph{On the Grothendieck ring of fusion categories of rank 5}, Comm. Algebra
\textbf{51} (2023), 4214--4222. \href{https://doi.org/10.1080/00927872.2023.2202777}{doi:10.1080/00927872.2023.2202777}.

\bibitem{Was} L. C. Washington,
\emph{Introduction to Cyclotomic Fields}, 2nd ed., Graduate Texts in Mathematics \textbf{83}, Springer, 1997.

\bibitem{Wenzl} H. Wenzl,
\emph{$C^*$ tensor categories from quantum groups},
J. Amer. Math. Soc. \textbf{11} (1998), no.~2, 261--282.
\href{https://doi.org/10.1090/S0894-0347-98-00253-7}{doi:10.1090/S0894-0347-98-00253-7}.
\end{thebibliography}


\end{document}
